\documentclass[11pt,letterpaper]{article}
\usepackage[T1]{fontenc}
\usepackage[utf8]{inputenc}
\usepackage{newtxtext}
\usepackage{amsmath,amsthm,mathtools}
\usepackage{newtxmath}
\usepackage[margin=1in,headheight=15pt]{geometry}
\usepackage{microtype}
\usepackage{enumitem,booktabs,array,longtable}
\usepackage[dvipsnames]{xcolor}
\usepackage{fancyhdr}
\usepackage[unicode,colorlinks=true,linkcolor=MidnightBlue,citecolor=MidnightBlue,urlcolor=MidnightBlue]{hyperref}
\usepackage[nameinlink,noabbrev]{cleveref}
\usepackage{bookmark}
\usepackage{xurl}
\hypersetup{pdftitle={The First-kappa Coefficients: Omitted Types and Completion in Surreal Hahn Fields, and Support-Bounded Surreal Fields},pdfauthor={OpenAI ChatGPT; research manuscripts prepared for Vladimir Reshetnikov},pdfsubject={Cardinally bounded Hahn fields, support-bounded surreal fields, singular cardinals, model theory, surreal numbers}}
\pagestyle{fancy}
\fancyhf{}
\fancyhead[L]{\small The first-$\kappa$ coefficients}
\fancyhead[R]{\small Surreal Hahn fields}
\fancyfoot[C]{\thepage}
\renewcommand{\headrulewidth}{0.3pt}
\setlength{\parindent}{1.2em}
\setlength{\parskip}{0.2em}
\setlength{\emergencystretch}{2em}
\setcounter{tocdepth}{2}
\setlist{itemsep=0.2em,topsep=0.4em}
\numberwithin{equation}{section}
\newtheorem{theorem}{Theorem}[section]
\newtheorem{proposition}[theorem]{Proposition}
\newtheorem{lemma}[theorem]{Lemma}
\newtheorem{corollary}[theorem]{Corollary}
\theoremstyle{definition}
\newtheorem{definition}[theorem]{Definition}
\newtheorem{example}[theorem]{Example}
\newtheorem{question}[theorem]{Question}
\theoremstyle{remark}
\newtheorem{remark}[theorem]{Remark}
% Added during assembly: all environments share the theorem counter, and cleveref
% would otherwise print every one of them as "theorem".
\makeatletter
\@for\fkc@env:=proposition,lemma,corollary,definition,example,question,remark\do{%
  \edef\fkc@tmp{\noexpand\AddToHook{env/\fkc@env/begin}{\noexpand\crefalias{theorem}{\fkc@env}}}%
  \fkc@tmp}
\makeatother
\crefname{theorem}{Theorem}{Theorems}
\crefname{proposition}{Proposition}{Propositions}
\crefname{lemma}{Lemma}{Lemmas}
\crefname{corollary}{Corollary}{Corollaries}
\crefname{definition}{Definition}{Definitions}
\crefname{example}{Example}{Examples}
\crefname{question}{Question}{Questions}
\crefname{remark}{Remark}{Remarks}
\crefname{section}{Section}{Sections}
\crefname{subsection}{Section}{Sections}
\crefname{appendix}{Appendix}{Appendices}
\newcommand{\R}{\mathbb R}
\newcommand{\C}{\mathbb C}
\newcommand{\Q}{\mathbb Q}
\newcommand{\N}{\mathbb N}
\newcommand{\No}{\mathbf{No}}
\newcommand{\supp}{\operatorname{supp}}
\newcommand{\otp}{\operatorname{otp}}
\newcommand{\cf}{\operatorname{cf}}
\newcommand{\tp}{\operatorname{tp}}
\newcommand{\om}{\operatorname{om}}
\newcommand{\lc}{\operatorname{lc}}
\newcommand{\RCF}{\mathrm{RCF}}
\newcommand{\ACF}{\mathrm{ACF}}
\newcommand{\ACVF}{\mathrm{ACVF}}
\newcommand{\F}{F_k}
\newcommand{\K}{K_{\kappa,k}}
\newcommand{\KR}{K_{\kappa,\R}}
\newcommand{\KC}{K_{\kappa,\C}}
\newcommand{\FR}{F_{\R}}
\newcommand{\FC}{F_{\C}}
\newcommand{\pre}{T_\kappa}
\newcommand{\dd}{D_x}
\newcommand{\Ball}[2]{B(#1,#2)}
\newcommand{\repoid}{\texttt{465a54b479a1ee842cbf7db1689a7d2f6bfe25e1}}
\newcommand{\status}[1]{\par\smallskip\noindent\textbf{Status.} #1\par\smallskip}
% Added for Part II (batch 32): notation of the class version, see Section 13.4.
\crefname{part}{Part}{Parts}
\newcommand{\mergetag}{\textup{\textsf{\small[merge]}}}% text added in the merge
\newcommand{\Z}{\mathbb Z}
\newcommand{\Oz}{\mathbf{Oz}}
\newcommand{\SC}{\mathbf{No}[i]}
\newcommand{\On}{\mathrm{On}}
\newcommand{\suppw}{\operatorname{supp}_{\omega}}
\newcommand{\Frac}{\operatorname{Frac}}
\newcommand{\Fin}{\operatorname{Fin}}
\newcommand{\ct}{\operatorname{ct}}
\newcommand{\lead}{\operatorname{le}}
\newcommand{\restr}{\mathbin{\upharpoonright}}
\newcommand{\Fix}{\operatorname{Fix}}
\newcommand{\Span}{\operatorname{span}_{\Q}}
\newcommand{\card}[1]{\lvert #1\rvert}
\newcommand{\conj}{\mathrm{conj}}
\newcommand{\KN}[2]{K_{#1,#2}(\No)}% the class field of surreals with < #1 terms
\newcommand{\Ok}{\Oz_{<\kappa}}
\newcommand{\Okl}[1]{\Oz_{<#1}}
\newcommand{\Pk}{\Pi_{<\kappa}}
\newcommand{\lbl}[1]{\nolinkurl{#1}}% a label, printed breakably

\title{\textbf{The First-$\boldsymbol\kappa$ Coefficients}\\[0.35em]
\Large Omitted Types, Singular Compression, and Completion\\in Surreal and Surcomplex Hahn Fields,\\and the Support-Bounded Surreal Fields}
\author{Research manuscripts prepared for Vladimir Reshetnikov\\\small Part I: OpenAI ChatGPT}
\date{Part I 22 September 2026; Part II 23 September 2026}

\begin{document}
\maketitle
\thispagestyle{empty}
\begin{abstract}
Let $\kappa$ be any uncountable cardinal, possibly singular, let $\Gamma\ne0$ be a divisible ordered abelian group, and let
$K_{\kappa,k}\subseteq k((t^\Gamma))$ consist of the Hahn series having fewer than $\kappa$ nonzero coefficients, with $k=\R$ or $\C$.
We give an explicit description of the one-variable types over this field that are realized in its full Hahn extension but omitted in the smaller field. The invariant is exactly the first $\kappa$ nonzero terms of the series: these terms determine its complete ordered-field type in the real case and its complete valued-field type in the complex case. We compute the least number of formulas needed to witness omission of each such type and obtain $\cf(\kappa)$, including at singular cardinals. In the pure complex field language, by contrast, all the missing elements have the same type and the corresponding number is $|K_{\kappa,\C}|$.

The same coefficient invariant gives the exact approximation-value cut, the least cardinality of an empty nest of closed valuation balls, and a concrete description of the valuation completion. Whenever $K_{\kappa,k}$ is proper, it is complete if and only if $\cf(\Gamma)\ne\cf(\kappa)$; it is never spherically complete, and its first empty ball nests have cardinality $\cf(\kappa)$. We also compute the precise valuation loss of every linear retraction from $K+Kx$ to $K$. Explicit Conway normal forms give surreal examples with countable omitted types at $\kappa=\aleph_\omega$, as well as complete fields exhibiting the same countable ball obstruction.

The constructions, general approximation-type theory, and quantifier-elimination inputs are established mathematics. The proposed contribution is the joint first-$\kappa$ classification and its exact cardinal, topological, and linear consequences. Full proofs are supplied. This is an unrefereed research contribution with a limited priority audit, not a claimed solution of a named longstanding conjecture or a proof-assistant-certified result.

Part~II, from a second manuscript, is the class version with every Conway monomial: for $\KN\kappa\R=\{x\in\No:|\suppw x|<\kappa\}$ it classifies the set-presented gaps by first-$\kappa$ prefixes, all of character $(\cf\kappa,\cf\kappa)$, proves $\nu$-saturation exactly for $\nu\le\cf\kappa$ and the thresholds $\cf\kappa$ for strong sums and empty nests, closure under the surreal exponential with no surjective ordered exponential at all, the set-sized quotients of its omnific integer part, and, under Global Choice, a proper class of pairwise nonconjugate real forms of $\No[i]$.

The report is built from two manuscripts: Part~I (Sections 1--12) from one pinned at repository commit \texttt{465a54b}, Part~II (Sections 13--26) from one pinned at \texttt{343dc2c}. \Cref{fkc:subsec:collection,fkc:sb:sec:intro} record their provenance, notation conventions and relation to neighbouring reports; \cref{fkc:subsec:nonclaims,fkc:sb:sub:nonclaims} collect their non-claims.
\end{abstract}
\vfill
\noindent\small\textbf{Conventions.} In Part~I, all Hahn fields and parameter sets are sets in ZFC, and all topological assertions use the intrinsic Hahn valuation topology, not the subspace topology inherited from the full class $\No$. Part~II works in G\"odel--Bernays class theory with proper-class fields and the valuation topology of $\No$ (\cref{fkc:sb:sec:setup}). The support bound is a \emph{cardinality} bound, not a bound on the magnitude of the exponents.
\clearpage
\begingroup
\small
\setlength{\parskip}{0pt}
\tableofcontents
\endgroup
\clearpage

\part{Set-sized Hahn fields with fewer than \texorpdfstring{$\kappa$}{kappa} terms (source 01)}
\label{fkc:sb:part:set}

\section{The problem and the contribution}
\label{fkc:sec:intro}

A series can fail to belong to a subfield for very different reasons. It may require new coefficients, a larger value group, or too many terms. The last obstruction is especially natural for surreal numbers: Conway normal form gives each individual surreal a set-sized, reverse-well-ordered sequence of monomials, but no fixed cardinal bounds the normal forms of all surreals. A support restriction therefore provides a set-theoretic way to organize exact surreal arithmetic.

This paper asks a more precise question than whether the support-restricted field is proper:
\begin{quote}
What can first-order formulas over a cardinally support-restricted field detect about a missing series, and how many formulas are needed to force its absence from the field?
\end{quote}
The answer will retain the coefficients of an initial segment, not merely the order type or cofinality of the support. Its cardinal consequence is that a singular support bound compresses a large amount of missing information into only $\cf(\kappa)$ conditions.

\subsection{Standing setup}
Fix a nonzero divisible ordered abelian group $\Gamma$, an uncountable cardinal $\kappa$, and write $\mu=\cf(\kappa)$. Cardinals are identified with their initial ordinals. For $k\in\{\R,\C\}$, put
\begin{equation}\label{fkc:eq:fields}
 F_k=k((t^\Gamma)),\qquad
 K_{\kappa,k}=\{f\in F_k:|\supp f|<\kappa\}.
\end{equation}
The support is well ordered in the increasing order of $\Gamma$, and
\[
 v(f)=\min\supp f\quad(f\ne0),\qquad v(0)=\infty.
\]
Thus larger values mean smaller elements. For real coefficients, the sign is the sign of the coefficient at the least exponent. We reserve $\cf(\Gamma)$ for the least cardinality of a cofinal subset of $\Gamma$ toward $+\infty$.

If $x\in F_k\setminus K_{\kappa,k}$, enumerate its support increasingly as
\begin{equation}\label{fkc:eq:enumerate}
 x=\sum_{\alpha<\lambda}c_\alpha t^{s_\alpha},\qquad
 c_\alpha\ne0,\quad s_\alpha<s_\beta\ (\alpha<\beta),\quad\lambda\ge\kappa.
\end{equation}
Its \emph{first-$\kappa$ prefix} is
\begin{equation}\label{fkc:eq:prefix-intro}
 T_\kappa(x)=\sum_{\alpha<\kappa}c_\alpha t^{s_\alpha}.
\end{equation}
This series itself is not in $K_{\kappa,k}$. Every proper initial segment of it is.

\subsection{A guide to the main results}
The following statements are proved separately below. The real types use the language of ordered rings; the complex valued types use the language of rings with the binary relation $a\mid_v b$ meaning $v(a)\le v(b)$. All parameters belong to the indicated smaller field.

\begin{theorem}[Main theorem package]\label{fkc:thm:main}
In the standing setup, the following hold.
\begin{enumerate}[label=\textup{(\roman*)}]
\item $K_{\kappa,\R}$ is real closed and $K_{\kappa,\C}$ is algebraically closed, even when $\kappa$ is singular. Their inclusions in the full Hahn fields are elementary in the stated languages.
\item For missing elements $x,y$ in the same full Hahn field,
\[
 \tp(x/K_{\kappa,k})=\tp(y/K_{\kappa,k})
 \quad\Longleftrightarrow\quad T_\kappa(x)=T_\kappa(y).
\]
\item For each such type $p$, the least size of a subset of $p$ with no realization in $K_{\kappa,k}$ is exactly $\mu$.
\item If $K_{\kappa,k}\ne F_k$, every nest of fewer than $\mu$ closed valuation balls meets, and some nest of $\mu$ such balls does not.
\item The completion is the subfield of $F_k$ consisting of those $f$ for which
\[
 |\supp f\cap(-\infty,\gamma)|<\kappa\qquad(\gamma\in\Gamma).
\]
If $K_{\kappa,k}\ne F_k$, this is a proper extension of $K_{\kappa,k}$ exactly when $\cf(\Gamma)=\mu$.
\end{enumerate}
\end{theorem}

The pure complex field language gives a deliberately contrasting answer: there is only one missing one-variable type, and its omission number is $|K_{\kappa,\C}|$ rather than $\mu$; see \cref{fkc:thm:pure}. The linear application in \cref{fkc:thm:projection} provides an exact bound for every retraction $K+Kx\to K$, not merely an existence or nonexistence assertion.

\mergetag{} \Cref{fkc:sb:part:class} proves the class version of parts (i)--(iv) with $\Gamma$ replaced by all of $\No$, for the field $\KN\kappa\R$ of surreals with fewer than $\kappa$ terms, which contains every Conway monomial (\cref{fkc:sb:thm:main}). There the omitted types become the set-presented gaps of a proper-class field, and part~(v) has no analogue: every set-indexed Cauchy net is eventually constant.

\subsection{What is established, and what is proposed here?}
Cardinally bounded Hahn fields are not new. Kuhlmann and Shelah use $\kappa$-bounded Hahn fields for regular uncountable $\kappa$ in their construction of exponential-logarithmic fields \cite{KS}. Related bounded constructions appear in the work of Berarducci, Kuhlmann, Mantova, and Matusinski on omega-maps \cite{BKMM}. We do not claim the construction, Hahn arithmetic, or its familiar regular-cardinal instances as new.

General approximation-type theory already describes immediate extensions by nests of balls. Kuhlmann's account includes their relation to first-order one-types \cite{FK}; the passage from polynomial data to complete types uses classical quantifier elimination \cite{Tarski,Robinson,Yin}. The contribution proposed here is an explicit \emph{cardinal-prefix calculation} for the fields in \eqref{fkc:eq:fields}, together with its sharp singular-cardinal and completion consequences. In particular, the exact omission number is proved for arbitrary formulas, not only for the displayed ball conditions.

The manuscript's repository comparison was made at its pinned revision; see \cref{fkc:sec:audit}. Its report on transcendence over bounded support concerns order-bounded supports and their fraction field, not the cardinal restriction in \eqref{fkc:eq:fields}. Its foundations and Hahn-vector-space reports supply a natural context, but the present proofs do not depend on a new, unaudited theorem from those manuscripts. A limited search did not identify the entire theorem package in the material reviewed. This is a scoped novelty assessment, not a guarantee that no equivalent statement exists. The comparison with the collection as it now stands, including countable-support analogues of \cref{fkc:thm:completion,fkc:thm:dichotomy} in two reports that the manuscript's comparison did not mention, is in \cref{fkc:subsec:collection}.

\subsection{This report in the collection}
\label{fkc:subsec:collection}

\paragraph{Provenance.}
Part~I of this report (Sections 1--12) is built from one manuscript, \emph{The First-$\kappa$ Coefficients: Omitted Types, Singular Compression, and Completion in Surreal and Surcomplex Hahn Fields}, dated 22 September 2026 (21 pages, author line ``OpenAI ChatGPT''), delivered as the archive \texttt{surreal\_first\_kappa\_types} (item~04 of batch~22) and placed as raw files in commit \texttt{7b5f934}. The report was assembled in commit \texttt{68e2960}. Its repository comparison is pinned at commit \texttt{465a54b479a1ee842cbf7db1689a7d2f6bfe25e1}. When assembled, the report was not a merge: no other manuscript proves its results, and the nearest material in the collection is cross-referenced below, not merged. \mergetag{} In batch~32 a second manuscript was added as \cref{fkc:sb:part:class} (source~02; its provenance is in \cref{fkc:sb:sub:provenance}); the manuscript of this part is source~01. The statements, proofs, limitations and priority caveats are the manuscript's. During assembly the following changed.
\begin{itemize}
\item Every label received the prefix \texttt{fkc:}; all 72 labels of the manuscript are kept, and no statement was renumbered or removed.
\item Added: this subsection; \cref{fkc:rem:countable-analogue} (the countable analogues in other reports); \cref{fkc:rem:groups-elsewhere} (the groups $\Gamma_\theta$ elsewhere in the collection); a pointer after the rank-one boundary case; the corrections of \cref{fkc:subsec:corrections}; the collected non-claims of \cref{fkc:subsec:nonclaims}; bibliography entries for collection reports; and the file paths of \cref{fkc:app:checks}, which now name this directory's layout.
\item The program \texttt{code/verify.py}, the script \texttt{code/build.sh}, the recorded run \texttt{data/verification.json}, the build record \texttt{data/build\_audit.json} and the file \texttt{research\_audit.md} are the delivered files, unchanged. The delivered member README and PDF are not shipped; \texttt{data/build\_audit.json} describes the delivered 21-page PDF, not the PDF of this report.
\item Cross-references are now printed with the right names. All environments share one counter, and in the delivered build every reference to a proposition, lemma, corollary or example was printed as ``theorem'', and the reference to the appendix as ``section''. The preamble now declares the names; no statement changed.
\item The family \texttt{surcomplex/} was chosen over \texttt{surreal/} because the results are stated uniformly for $k=\R$ and $k=\C$; this was the only placement choice.
\end{itemize}

\paragraph{Notation.}
No symbol of the manuscript of Part~I was renamed (the renamings in \cref{fkc:sb:part:class} are listed in \cref{fkc:sb:sub:notation}). Its symbols are fully subscripted and local to this report, and the following conventions fix them against the collection's notation guide (\texttt{docs/NOTATION.md}) and neighbouring reports.
\begin{itemize}
\item $F_k=k((t^\Gamma))$ denotes the full Hahn field for \emph{both} $k=\R$ and $k=\C$. The guide writes $F_\Gamma=\R((t^\Gamma))$ for the real workspace and $K_\Gamma=\C((t^\Gamma))$ for the complex one. Thus $F_\R$ here is the guide's $F_\Gamma$, and $F_\C$ here is the guide's $K_\Gamma$. The subscript $\Gamma$ is suppressed because $\Gamma$ is fixed throughout.
\item $K_{\kappa,k}$ is the \emph{cardinally bounded subfield} of \eqref{fkc:eq:fields}, real or complex. It is never the guide's full complex workspace $K_\Gamma$. From \cref{fkc:sec:linear} on, $K$ and $F$ abbreviate $K_{\kappa,k}$ and $F_k$.
\item $L_{\kappa,k}$ is the completion of \eqref{fkc:eq:localfield}; $D_x$ is the approximation cut of \eqref{fkc:eq:D}, not the quaternion algebra $D_\Gamma$ of the surquaternions report; $E_x=K+Kx$ is a two-dimensional valued space, neither the algebraic scalar extension $E_\Gamma(V)$ of the three-duals report nor the ring of entire functions of the entire-functions report.
\item $\kappa$ is an uncountable cardinal and $\mu=\cf(\kappa)$ is a regular cardinal. Neither is a measure, a residue field (the entire-functions report writes $\kappa_n$ for residue fields) or a valuation gain (the expanding-dynamics report writes $\kappa=v(q)$).
\item $I_D$ in \eqref{fkc:eq:ID} is an additive subgroup and not an ideal; it is unrelated to the error ideal $I_\kappa$ of the expanding-dynamics report. $T_\kappa(x)$ is an external invariant, not an operation of any language used here.
\item $\cf(\Gamma)$ is cofinality toward $+\infty$; $v=\min\supp$, larger values mean smaller elements, and $t^\gamma\mapsto\omega^{-j(\gamma)}$. These agree with the guide.
\end{itemize}
Four words carry more than one meaning, and each tempting false reading is excluded here.
\begin{itemize}
\item \emph{Bounded.} A support bound is a bound on the \emph{cardinality} of the support. In the report on transcendence over bounded support, ``bounded support'' means bounded above in the exponent order, and that report says explicitly that it is not a cardinal restriction.
\item \emph{Complete.} A \emph{complete type} is a maximal consistent set of formulas; a \emph{complete} field is complete for Cauchy nets of the valuation uniformity; a \emph{spherically complete} field has nonempty intersection for every nest of closed balls. \Cref{fkc:thm:dichotomy} shows that the last two differ: a proper $K_{\kappa,k}$ is never spherically complete but can be complete.
\item \emph{Closure.} \Cref{fkc:thm:completion} concerns topological closure inside $F_k$; \cref{fkc:prop:closed} concerns real and algebraic closedness.
\item \emph{Omitted.} An omitted type here is a one-type over the set-sized field $K_{\kappa,k}$, realized in $F_k$ and not in $K_{\kappa,k}$. It is not an omitted cut of a proper-class field.
\end{itemize}

\paragraph{Neighbouring reports.}
The comparisons below were checked against the current collection. None of them is used in a proof, and none proves or refutes a statement here.
\begin{itemize}
\item \texttt{foundations-and-computation/surreal-fields-across-universes} (prefix \texttt{univ:}) \cite{CollUniv} studies saturation and omitted cuts of the old proper-class field $\No^M$ inside $\No^N$ for inner models $M\subseteq N$. It shares the vocabulary of omitted types, saturation and cofinality spectra, and it too contrasts languages, but its objects are different: it has no Hahn field, no support bound and no completion, and its topology is the fine order topology of proper classes. Part~I answers none of its questions. It was added after this manuscript's pin. \mergetag{} Part~II answers, in a single universe and without forcing, the clause of \texttt{univ:q:families} on infinite collections of regular thresholds (\cref{fkc:sb:rem:univ}).
\item \texttt{surreal/transcendence-over-bounded-support} (prefix \texttt{bst:}) \cite{CollBST} works over the order-bounded base $\{f:\supp f\text{ bounded above}\}$ and its fraction field. It explicitly excludes cardinal support bounds (its definition of the base, labelled \texttt{bst:eq:base}, and its list of conventions). Its corollary \texttt{bst:cor:optimal} is consistent with this report: its proof uses the fact that a support of cardinality less than $\cf(\Gamma)$ is bounded above. This report answers none of its questions and is not a new version of its independence theorem. \mergetag{} \emph{Related (batch 32).} That report's Galois part asks (\texttt{bst:gr:q:cardinal}) which of its constructions survive in fields of Hahn series with fewer than $\kappa$ terms, at regular and at singular $\kappa$, and says that algebraic closure operations and the supports of all coefficients in a finite cover must be checked. For the algebraic closure step \cref{fkc:prop:closed} and \cref{fkc:rem:singular} give the answer at every uncountable $\kappa$, singular included: a polynomial over $K_{\kappa,\C}$ splits in one Hahn field $\C((t^H))$ with $\card H<\kappa$ (\cref{fkc:lem:workspace}), so $K_{\kappa,\C}$ is algebraically closed and $K_{\kappa,\R}$ real closed. The Galois constructions themselves, over bounded-support fraction fields, are not treated here, and that question stays open.
\item \texttt{surcomplex/three-duals-of-hahn-vector-spaces} (prefix \texttt{duals:}) \cite{CollDuals} proves the vector-space analogue with ``finite coefficient rank'' in place of ``fewer than $\kappa$ support terms''; see \cref{fkc:rem:countable-analogue}. Its full-Hahn spherical completeness, \texttt{duals:prop:spherical}, is the gluing argument that \cref{fkc:lem:glue} refines by a cardinal bound.
\item \texttt{surcomplex/rank-one-berkovich} (bare labels) \cite{CollBerk} proves spherical completeness of $\C((t^\R))$ (its label \texttt{prop:spherical}) and describes the closure of the finite-support series in $\C((t^\Q))$; see \cref{fkc:rem:countable-analogue} and the rank-one boundary case in \cref{fkc:sec:examples}.
\item \texttt{surcomplex/entire-functions-at-arbitrary-rank} (prefix \texttt{ent:}) \cite{CollEnt} uses the group $\Gamma_\infty=\bigoplus_{j\ge0}\Q e_j$ with $e_j=\omega^j$ and the group $\Gamma_{\omega_1}$. These are the groups $\Gamma_\omega$ and $\Gamma_{\omega_1}$ of \eqref{fkc:eq:Gtheta}, with the same embedding in $\No$; see \cref{fkc:rem:groups-elsewhere}.
\item \texttt{foundations-and-computation/foundations} (prefix \texttt{found:}) \cite{CollFound} distinguishes the intrinsic Hahn-workspace topology from the full fine topology (\texttt{found:prop:twotopologies}) and proves that small subsets and small-index nets are discrete in the latter (\texttt{found:thm:discrete}). \Cref{fkc:sec:surreal} relies on that distinction.
\item \texttt{surcomplex/hidden-negative-hermitian-directions} (prefix \texttt{hnd:}) \cite{CollHND} asks, in an unlabelled question, which support conditions on a support-family subfield $E\subseteq F_\Gamma$ give order density and which support invariants bound the degree of approximants to an omitted valuation ball. $K_{\kappa,\R}$ is such a subfield, but this report computes approximation cuts and empty ball nests, not order density or degrees of approximants, so that question is not answered here.
\item \mergetag{} Reports that bear on Part~II, most of them written after this part was assembled, are compared in \cref{fkc:sb:sub:relation}: the large-cardinal report, whose field of short normal forms is the class field of Part~II; the quotient report, whose support-threshold theorem contains Part~II's quotient theorem; the omnific-automorphisms report, whose set-sized field $K_{<\kappa}$ for regular $\kappa$ (\texttt{opa:as:prop:kappa}) is, for $k\in\{\R,\C\}$, the field $\K$ of Part~I; and the surcomplex-automorphisms report.
\end{itemize}

\section{Hahn arithmetic and singular support bounds}
\label{fkc:sec:algebra}

\subsection{The support conventions}
A Hahn series is a function $f:\Gamma\to k$ whose nonzero set is well ordered. Addition is coefficientwise, and multiplication is the Hahn convolution. The well-ordering and finite-contribution properties needed for multiplication are the standard Hahn--Neumann support facts. No ordinary convergence of finite partial sums is implicit in this notation.

For $\gamma\in\Gamma$, define
\[
 f_{<\gamma}=\sum_{g<\gamma}f_g t^g,\qquad
 f_{\le\gamma}=\sum_{g\le\gamma}f_g t^g.
\]
Subseries are Hahn series because a subset of a well-ordered set is well ordered. The essential valuation identity is
\begin{equation}\label{fkc:eq:first-diff}
 v(f-g)\ge\gamma
 \quad\Longleftrightarrow\quad
 f_h=g_h\text{ for every }h<\gamma.
\end{equation}
Strict inequality means agreement at every $h\le\gamma$. When $v(f)\ne v(g)$,
\begin{equation}\label{fkc:eq:ultra-exact}
 v(f+g)=\min\{v(f),v(g)\}.
\end{equation}
Both formulas follow by looking at the least exponent with nonzero coefficient.

\subsection{A small exponent group contains all finite algebraic data}
\begin{lemma}[Small-workspace lemma]\label{fkc:lem:workspace}
Let $A\subseteq F_k$ be finite and suppose each element of $A$ has support of cardinality less than $\kappa$. There is a divisible ordered subgroup $H\subseteq\Gamma$ with $|H|<\kappa$ such that
\[
 A\subseteq k((t^H))\subseteq K_{\kappa,k}.
\]
\end{lemma}
\begin{proof}
Let $S$ be the union of the supports of the finitely many elements of $A$. Then $|S|<\kappa$. Take its rational linear span $H$ in $\Gamma$. Since ordered abelian groups are torsion free and $\Gamma$ is divisible, this span is a well-defined divisible subgroup. Every element of $H$ is a finite rational linear combination of members of $S$, so
\[
 |H|\le\max\{|S|,\aleph_0\}<\kappa.
\]
The last strict inequality uses uncountability, not regularity, of $\kappa$. Every support contained in $H$ has cardinality at most $|H|$, which proves the second inclusion; the first is immediate from the choice of $S$.
\end{proof}

\begin{proposition}[Algebraic closedness and real closedness]\label{fkc:prop:closed}
$\K$ is a subfield of $F_k$. Moreover, $\KC$ is algebraically closed, $\KR$ is real closed, and
\[
 \KR[i]=\KC.
\]
The valuation of $\K$ has value group $\Gamma$ and residue field $k$, so $F_k/\K$ is immediate.
\end{proposition}
\begin{proof}
For two inputs, apply \cref{fkc:lem:workspace}. Their sum, product, and, for a nonzero input, inverse lie in the Hahn field $k((t^H))$. This proves the subfield assertion without any union of $\kappa$ many supports.

We use the classical theorem that a full Hahn field with algebraically closed coefficient field and divisible exponent group is algebraically closed; see \cite[Corollary~4]{Poonen}. Given a polynomial with coefficients in $\KC$, the small-workspace lemma puts all its coefficients in $\C((t^H))$. The polynomial splits there, and every root has support of size at most $|H|<\kappa$. Thus it splits in $\KC$.

Complex coefficients split uniquely into real and imaginary parts. Each resulting support is a subset of the original support, while the union of two supports of size less than $\kappa$ still has size less than $\kappa$. Hence $\KR[i]=\KC$. The field $\KR$ is ordered by its leading coefficient, and its quadratic extension by $i$ is algebraically closed. By the real-closed-field criterion, $\KR$ is real closed. Equivalently, one can apply the usual real Hahn theorem inside each small $H$.

Finally, every monomial $t^\gamma$ and every constant from $k$ has finite support. Thus all values in $\Gamma$ and all residues in $k$ already occur in $\K$. They are also the value group and residue field of $F_k$, proving immediacy.
\end{proof}

\begin{remark}[Why singular cardinals cause no algebraic failure]\label{fkc:rem:singular}
A countable union of sets of size less than a singular $\kappa$ can have size $\kappa$. That fact does not invalidate \cref{fkc:prop:closed}: a fixed polynomial has only finitely many coefficients, and all its roots can be found in \emph{one} exponent group of size less than $\kappa$. The relevant bound is not an uncontrolled union along a root construction. This distinction is central later: algebraic equations stay in one small workspace, whereas an infinite type may use parameters from many increasingly large workspaces.
\end{remark}

\subsection{The imported logical tools}
We use the following standard quantifier-elimination statements. Real closed fields eliminate quantifiers in the ordered-ring language \cite{Tarski}. Algebraically closed fields eliminate quantifiers in the ring language. Algebraically closed nontrivially valued fields eliminate quantifiers in the ring language expanded by valuation divisibility \cite{Robinson,Yin}. The last language has no named coefficient extraction, truncation operation, exponential, or cross-section.

It follows that $\KR\preccurlyeq\FR$ as ordered fields and $\KC\preccurlyeq\FC$ both as pure fields and as valued fields. Here $\preccurlyeq$ means elementary substructure. Each valuation is nontrivial because $\Gamma\ne0$.

We will only need the following concrete consequences. A formula in one variable over a real closed field is equivalent to a Boolean combination of polynomial sign conditions. A valued-field formula in one variable over an algebraically closed valued field is equivalent to a Boolean combination of polynomial equations and comparisons $v(P(X))\le v(Q(X))$. Over an algebraically closed field, a one-variable definable set is finite or cofinite. These are background inputs, not claims of new quantifier elimination.

\section{The first-\texorpdfstring{$\kappa$}{kappa} approximation theorem}
\label{fkc:sec:approx}

For a missing element $x$ as in \eqref{fkc:eq:enumerate}, define
\begin{equation}\label{fkc:eq:D}
 S_\kappa(x)=\{s_\alpha:\alpha<\kappa\},\qquad
 D_x=\{g\in\Gamma:g\le s_\alpha\text{ for some }\alpha<\kappa\}.
\end{equation}
The set $D_x$ is an initial segment of $\Gamma$ with no greatest element. It can be all of $\Gamma$ or a proper initial segment. Do not replace $D_x$ by a presumed supremum in $\Gamma$: that supremum need not exist.

\begin{theorem}[Exact approximation values]\label{fkc:thm:approx}
For every $x\in F_k\setminus\K$,
\begin{equation}\label{fkc:eq:Aexact}
 A(x/\K):=\{v(x-a):a\in\K\}=D_x,
 \qquad \cf(D_x)=\cf(\kappa)=\mu.
\end{equation}
In particular, $x$ has no best approximation in $\K$ in the sense of a largest value of $v(x-a)$.
\end{theorem}
\begin{proof}
Fix $a\in\K$. It cannot have the same nonzero coefficient as $x$ at every $s_\alpha$ for $\alpha<\kappa$, because that would place $\kappa$ distinct elements in $\supp a$. Choose $\alpha<\kappa$ where those coefficients differ. The least exponent where $x$ and $a$ differ is at most $s_\alpha$. Thus $v(x-a)\in D_x$.

Conversely, suppose $g\in D_x$ and choose $\alpha<\kappa$ with $g\le s_\alpha$. The support of $x_{<g}$ is contained in $\{s_\beta:\beta<\alpha\}$, so it has size less than $\kappa$. Choose $b\in k$ different from the coefficient $x_g$, and put
\[
 a=x_{<g}+b t^g.
\]
This lies in $\K$, agrees with $x$ below $g$, and differs at $g$. Hence $v(x-a)=g$, proving the equality.

Choose a strictly increasing cofinal family $(\alpha_i)_{i<\mu}$ in the ordinal $\kappa$. Then $(s_{\alpha_i})_{i<\mu}$ is cofinal in $D_x$, so $\cf(D_x)\le\mu$. For the reverse bound, suppose $E\subseteq D_x$ is cofinal and $|E|<\mu$. For each $e\in E$, choose $\beta_e<\kappa$ with $e\le s_{\beta_e}$. Since $|E|<\cf(\kappa)$, these indices are bounded below $\kappa$. A later $s_\beta$ exceeds every member of $E$, contradicting cofinality. Thus $\cf(D_x)=\mu$.

Finally, for every $s_\alpha$ in the initial segment there is a later $s_\beta$, because the infinite initial ordinal $\kappa$ is a limit ordinal. Therefore $D_x$ has no greatest element, which proves the final assertion.
\end{proof}

\begin{definition}[The invisible tail subgroup]
For an initial segment $D\subseteq\Gamma$, set
\begin{equation}\label{fkc:eq:ID}
 I_D=\{0\}\cup\{h\in F_k\setminus\{0\}:v(h)>d\text{ for all }d\in D\}.
\end{equation}
This is an additive subgroup. No assertion that it is an ideal of the field $F_k$ is intended. If $D=\Gamma$, then $I_D=\{0\}$.
\end{definition}

\begin{lemma}[Prefix fibers]\label{fkc:lem:fibers}
For missing elements $x,y\in F_k\setminus\K$, the following are equivalent:
\[
 T_\kappa(x)=T_\kappa(y),\qquad x-y\in I_{D_x},\qquad
 y\in T_\kappa(x)+I_{D_x}.
\]
Every element of $T_\kappa(x)+I_{D_x}$ is missing from $\K$.
\end{lemma}
\begin{proof}
Equal prefixes mean that the two series agree at every exponent at most $s_\alpha$, for every $\alpha<\kappa$. Any nonzero difference therefore has value larger than all those exponents, hence larger than all of $D_x$. Conversely, such a difference changes none of the first $\kappa$ nonzero coefficients and introduces no additional term before them. It cannot change the first-$\kappa$ prefix. The difference $x-T_\kappa(x)$ itself belongs to $I_{D_x}$, giving the third formulation. Each series in that coset still contains all $\kappa$ nonzero terms of the prefix, so it cannot lie in $\K$.
\end{proof}

\begin{remark}[A cut without coefficients is insufficient]
Two series may have identical sets $D_x$ but different first-$\kappa$ prefixes. For example, changing a single coefficient in the prefix does not change $D_x$. Approximation \emph{values} forget the centers of the balls; the full first-order type will remember the coefficients through those centers. The distinction between $D_x$ and $T_\kappa(x)$ prevents an erroneous classification solely by cofinality or valuation rank.
\end{remark}

\section{Complete one-types and their fibers}
\label{fkc:sec:types}

A complete one-type over $K$ is the collection of all formulas $\varphi(X)$ with parameters in $K$ that hold of an element in an elementary extension. We only classify the types realized in the specified full Hahn extension $F_k$. Other elementary extensions can add new residue elements or values, and their types are outside this classification.

\subsection{The complex valued-field language}
\begin{theorem}[First-$\kappa$ classification: valued case]\label{fkc:thm:valued-type}
For $x,y\in\FC\setminus\KC$,
\begin{equation}\label{fkc:eq:valued-type}
 \tp_{\ACVF}(x/\KC)=\tp_{\ACVF}(y/\KC)
 \quad\Longleftrightarrow\quad T_\kappa(x)=T_\kappa(y).
\end{equation}
The set of realizations in $\FC$ of the type of $x$ is exactly
\[
 T_\kappa(x)+I_{D_x}.
\]
\end{theorem}
\begin{proof}
Suppose first that the prefixes agree. By \cref{fkc:lem:fibers}, either $x=y$ or $v(x-y)$ is greater than every element of $D_x$. For each $a\in\KC$, \cref{fkc:thm:approx} gives $v(x-a)\in D_x$, and therefore
\begin{equation}\label{fkc:eq:linear-data}
 v(y-a)=v(x-a),\qquad \lc(y-a)=\lc(x-a).
\end{equation}
Indeed, the two differences differ only at strictly higher valuation.

Every nonzero polynomial over $\KC$ factors as
\[
 P(X)=b\prod_{j=1}^{n}(X-a_j),\qquad b\in\KC^\times,\quad a_j\in\KC,
\]
with multiplicities retained. Neither $x$ nor $y$ is a root, since $\KC$ is algebraically closed and they do not belong to it. Formula \eqref{fkc:eq:linear-data} implies
\[
 v(P(x))=v(b)+\sum_j v(x-a_j)
         =v(b)+\sum_j v(y-a_j)=v(P(y)).
\]
The zero polynomial is handled separately, with value $\infty$. Thus $x$ and $y$ satisfy the same polynomial equations and the same comparisons of polynomial valuations. Quantifier elimination gives equality of their complete valued-field types.

Conversely, suppose the types agree. For every $\alpha<\kappa$, the element
\[
 b_\alpha=x_{\le s_\alpha}
\]
has fewer than $\kappa$ support terms and lies in $\KC$. The following is a formula over $\KC$, using $t^{s_\alpha}$ as a parameter to name its value:
\begin{equation}\label{fkc:eq:prefix-formula}
 v(X-b_\alpha)>v(t^{s_\alpha}).
\end{equation}
It holds of $x$, hence of $y$. The two series consequently agree at every exponent at most $s_\alpha$, for all $\alpha<\kappa$. Their prefixes agree. The assertion about the entire realization set follows from \cref{fkc:lem:fibers}, which also excludes realizations in $\KC$.
\end{proof}

\begin{corollary}[Rational-function data]\label{fkc:cor:rational}
If $x,y$ have the same first-$\kappa$ prefix, the isomorphism
\[
 \KC(x)\longrightarrow\KC(y),\qquad x\longmapsto y,
\]
fixing $\KC$ preserves both valuation and leading coefficient of every nonzero rational function. In the real case, the corresponding isomorphism of rational function fields is order preserving.
\end{corollary}
\begin{proof}
The missing elements are transcendental over their respective base fields. For complex coefficients, factor numerator and denominator into linear factors and apply \eqref{fkc:eq:linear-data}; leading coefficients multiply and divide. In the real case, the same identity shows that $x-a$ and $y-a$ have the same sign for every $a$ in the base. Factor real polynomials into real linear factors and monic irreducible quadratic factors, the latter positive everywhere. Signs of nonzero rational functions are therefore preserved.
\end{proof}

\subsection{The real ordered-field language}
\begin{theorem}[First-$\kappa$ classification: ordered case]\label{fkc:thm:real-type}
For $x,y\in\FR\setminus\KR$,
\begin{equation}\label{fkc:eq:real-type}
 \tp_{\RCF}(x/\KR)=\tp_{\RCF}(y/\KR)
 \quad\Longleftrightarrow\quad T_\kappa(x)=T_\kappa(y).
\end{equation}
The realizations in $\FR$ of the type of $x$ are exactly $T_\kappa(x)+I_{D_x}$.
\end{theorem}
\begin{proof}
Equal prefixes imply equal signs of $x-a$ and $y-a$ for all $a\in\KR$, by the leading-coefficient argument in \eqref{fkc:eq:linear-data}. A nonzero real polynomial factors into linear factors and positive monic irreducible quadratic factors, together with a nonzero scalar. Thus its values at $x$ and $y$ have the same sign. Quantifier elimination for real closed fields now gives equal types.

For the reverse direction, set $b_\alpha=x_{\le s_\alpha}\in\KR$. The ordered-ring formula
\begin{equation}\label{fkc:eq:ordered-interval}
 b_\alpha-t^{s_\alpha}<X<b_\alpha+t^{s_\alpha}
\end{equation}
holds of $x$: its difference from $b_\alpha$ has valuation strictly greater than $s_\alpha$ and is therefore smaller in absolute value than the positive monomial $t^{s_\alpha}$. Equal types imply that \eqref{fkc:eq:ordered-interval} holds of $y$ as well.

The inequality $|y-b_\alpha|<t^{s_\alpha}$ forces agreement of $y$ and $b_\alpha$ at all exponents strictly below $s_\alpha$. If there were a least disagreement $g<s_\alpha$, its nonzero real leading coefficient times $t^g$ would dominate $t^{s_\alpha}$ in absolute value. For any fixed $\beta<\kappa$, choose $\alpha$ strictly later than $\beta$. Agreement below $s_\alpha$ includes the coefficient at $s_\beta$ and all earlier exponents. Consequently the first-$\kappa$ prefixes agree. The realization-set description follows as before.
\end{proof}

\begin{corollary}[Counting and identifying the missing types]\label{fkc:cor:count}
In either language of \cref{fkc:thm:valued-type,fkc:thm:real-type}, the missing one-types realized in $F_k$ are in bijection with the Hahn series whose support has order type exactly $\kappa$. The bijection sends a type to its first-$\kappa$ prefix. If $K_{\kappa,k}\ne F_k$, there are at least $2^\kappa$ such types.
\end{corollary}
\begin{proof}
Each missing element has exactly one prefix of order type $\kappa$, and each such prefix is itself a missing element. The two type-classification theorems show that this correspondence is well defined and bijective. Fix an increasing support $(s_\alpha)_{\alpha<\kappa}$. Choosing each coefficient independently from $\{1,2\}$ gives $2^\kappa$ different prefixes and therefore $2^\kappa$ types.
\end{proof}

\begin{remark}[Types are not asserted to be automorphism orbits in $F_k$]
Equality of types guarantees conjugacy in a sufficiently homogeneous elementary extension. It does not, by itself, produce an automorphism of the particular full Hahn field $F_k$ taking $x$ to $y$. No homogeneity theorem for $F_k$, and no extension of a partial automorphism to the full class $\No$, is used here.
\end{remark}

\section{The exact number of formulas needed for omission}
\label{fkc:sec:omission}

\begin{definition}[Omission number of a type]
Let $M\preccurlyeq N$ and let $p=\tp(x/M)$ be omitted in $M$. Define
\begin{equation}\label{fkc:eq:omdef}
 \om_M(p)=\min\{\,|\Phi|:\Phi\subseteq p
                  \text{ has no realization in }M\,\}.
\end{equation}
This minimum exists because $p$ itself is a set of formulas omitted in $M$. Each formula is finitary, but different formulas may use different parameters. This invariant is not the saturation cardinal of the entire structure.
\end{definition}

\subsection{Formula-by-formula stabilization}
Choose a strictly increasing cofinal sequence $(\alpha_i)_{i<\mu}$ in $\kappa$, and put
\begin{equation}\label{fkc:eq:approximants}
 a_i=x_{<s_{\alpha_i}}\in\K.
\end{equation}
Then
\begin{equation}\label{fkc:eq:approx-errors}
 v(x-a_i)=s_{\alpha_i},\qquad
 v(a_j-a_i)=s_{\alpha_i}\quad(i<j).
\end{equation}
The second equality uses the nonzero coefficient $c_{\alpha_i}$. Thus $(a_i)$ is a strictly pseudo-Cauchy family. It need not be a Cauchy net: its error values may stay bounded in $\Gamma$.

\begin{lemma}[Eventual elementary agreement]\label{fkc:lem:stabilize}
In the real ordered-field or complex valued-field language, for every formula $\varphi(X)$ over $\K$ there is $i_0<\mu$ such that
\[
 F_k\models\varphi(a_i)\ \Longleftrightarrow\ F_k\models\varphi(x)
 \qquad(i_0\le i<\mu).
\]
\end{lemma}
\begin{proof}
First consider a fixed $b\in\K$. By \cref{fkc:thm:approx}, $r=v(x-b)$ belongs to $D_x$. For all sufficiently large $i$, $s_{\alpha_i}>r$. Hence
\[
 v(a_i-b)=v(x-b),\qquad \lc(a_i-b)=\lc(x-b).
\]
In particular, $a_i\ne b$ eventually.

For a finite collection of nonzero complex polynomials over $\KC$, take their finitely many roots, with multiplicity ignored for this choice of a bound. Choose $i_0$ making the preceding identities hold for every root. For all $i\ge i_0$, each polynomial is nonzero at $a_i$ and has the same value as at $x$. Consequently every Boolean combination of their polynomial equations and valuation comparisons stabilizes. Quantifier elimination reduces any fixed valued-field formula to such a finite Boolean combination.

Over $\KR$, instead factor the finitely many polynomials from a quantifier-free ordered formula into linear and positive irreducible quadratic factors. Their linear factors have stable signs by the first paragraph, and the quadratic signs are already constant. All polynomial sign conditions stabilize, and real quantifier elimination proves the assertion. Constant and identically zero polynomials require no stabilization condition.
\end{proof}

\begin{theorem}[Exact omission number]\label{fkc:thm:omission}
For each $x\in F_k\setminus\K$, in the ordered real or valued complex language,
\begin{equation}\label{fkc:eq:om-main}
 \boxed{\ \om_{\K}\bigl(\tp(x/\K)\bigr)=\cf(\kappa).\ }
\end{equation}
More precisely, every family of fewer than $\mu$ formulas from the type is realized by a sufficiently late truncation $a_i$, whereas an explicitly given family of $\mu$ ball or interval formulas is omitted.
\end{theorem}
\begin{proof}
For the lower bound, let $\Phi$ be a subset of the type of cardinality less than $\mu$. For each $\varphi\in\Phi$, choose a stabilization index $i_\varphi$ from \cref{fkc:lem:stabilize}. The cardinal $\mu=\cf(\kappa)$ is regular, so fewer than $\mu$ such indices are bounded by some $i_*<\mu$. Then $a_i$ realizes every formula in $\Phi$ for all $i\ge i_*$. Elementarity lets us view this realization either in $F_k$ or in $\K$.

For the upper bound in the complex valued case, use the $\mu$ formulas
\begin{equation}\label{fkc:eq:om-ball}
 v(X-a_i)\ge v(t^{s_{\alpha_i}})\qquad(i<\mu).
\end{equation}
They all hold of $x$. A common realization $a\in\KC$ would agree with $x$ at every exponent below each $s_{\alpha_i}$. Cofinality of the indices then forces agreement at all first $\kappa$ nonzero terms, placing at least $\kappa$ elements in $\supp a$. This is impossible.

For the upper bound in the real case, put $b_i=x_{\le s_{\alpha_i}}$ and use
\begin{equation}\label{fkc:eq:om-interval}
 b_i-t^{s_{\alpha_i}}<X<b_i+t^{s_{\alpha_i}}\qquad(i<\mu).
\end{equation}
These formulas belong to the type of $x$. As in the proof of \cref{fkc:thm:real-type}, a common realization must agree with $x$ below every $s_{\alpha_i}$ and hence on all first $\kappa$ terms. Again it cannot lie in $\KR$. This proves the upper bound and the equality.
\end{proof}

\subsection{What forgetting the valuation does}
\begin{theorem}[Pure-field collapse]\label{fkc:thm:pure}
Assume $\KC\ne\FC$ and work in the pure ring language. All elements of $\FC\setminus\KC$ realize the same complete one-type $p_{\mathrm{gen}}$ over $\KC$. Its omission number is
\begin{equation}\label{fkc:eq:pureom}
 \om_{\KC}(p_{\mathrm{gen}})=|\KC|.
\end{equation}
\end{theorem}
\begin{proof}
Since $\KC$ is algebraically closed, every element of $\FC\setminus\KC$ is transcendental over it. Such an element satisfies precisely the generic one-type over the algebraically closed field: all nonzero polynomials over $\KC$ are nonzero there. Quantifier elimination makes this a complete type.

Every one-variable formula from the generic type defines a cofinite subset of $\KC$. Indeed, the definable subset is finite or cofinite, and a finite set definable over an algebraically closed base cannot contain a transcendental element in an elementary extension. Fewer than $|\KC|$ finite exceptional sets cannot cover the infinite field $\KC$. This remains true if $|\KC|$ is singular: for an infinite cardinal $\nu<|\KC|$, a union of $\nu$ finite sets has cardinality at most $\nu$. Therefore fewer than $|\KC|$ formulas from $p_{\mathrm{gen}}$ have a common realization in $\KC$.

Conversely, the family $\{X\ne a:a\in\KC\}$ belongs to the generic type and has no realization in $\KC$. Its cardinality is $|\KC|$, proving the equality.
\end{proof}

In particular, when $\kappa$ is singular and $\KC\ne\FC$, the same individual missing element has omission number $\cf(\kappa)<\kappa\le|\KC|$ in the valued language and omission number $|\KC|$ in the pure field language. This is a quantitative distinction between two descriptions of a surcomplex workspace, not a change in its underlying arithmetic.

\subsection{Singular compression along elementary chains}
\begin{corollary}[Restriction of types is truncation of prefixes]\label{fkc:cor:tower}
Let $\kappa<\lambda$ be uncountable cardinals, with the same $k$ and $\Gamma$. Then
\[
 K_{\kappa,k}\preccurlyeq K_{\lambda,k}\preccurlyeq F_k
\]
in the real ordered or complex valued language. If $x\notin K_{\lambda,k}$, restriction of its type from $K_{\lambda,k}$ to $K_{\kappa,k}$ corresponds to the map
\[
 T_\lambda(x)\longmapsto T_\kappa(x).
\]
The two omission numbers are respectively $\cf(\lambda)$ and $\cf(\kappa)$.
\end{corollary}
\begin{proof}
The inclusions are elementary by \cref{fkc:prop:closed} and quantifier elimination. The assertions about prefixes and omission numbers follow from \cref{fkc:thm:valued-type,fkc:thm:real-type,fkc:thm:omission} applied at the two cardinals.
\end{proof}

For example, moving from the bound $\aleph_1$ to the larger bound $\aleph_\omega$ can \emph{decrease} the omission number from $\aleph_1$ to $\aleph_0$. There is no conflict with compactness: each parameter in the larger field is a single object but may encode a much longer initial segment. Countably many such parameters can encode $\aleph_\omega$ support positions. At the smaller bound, no analogous countable family of parameters can do so for these types.

\section{The precise spherical-completeness threshold}
\label{fkc:sec:balls}

A closed valuation ball in a valued field $K$ is
\[
 B(a,\gamma)=\{z\in K:v(z-a)\ge\gamma\},\qquad a\in K,\quad\gamma\in\Gamma.
\]
A \emph{nest} is a nonempty set of such balls linearly ordered by inclusion. Spherical completeness means that every nest has nonempty intersection. The radii here are values in $\Gamma$, not necessarily real numbers; nests need not be sequences.

\begin{lemma}[Coefficient gluing with a cardinal bound]\label{fkc:lem:glue}
Let $\{B(a_j,\gamma_j):j\in J\}$ be a nest of balls in $\K$. If $|J|<\cf(\kappa)$, its intersection contains an element of $\K$.
\end{lemma}
\begin{proof}
For any $i,j\in J$, nesting implies
\[
 v(a_i-a_j)\ge\min\{\gamma_i,\gamma_j\}.
\]
Consequently, for a fixed $g\in\Gamma$, all the coefficients $(a_j)_g$ with $g<\gamma_j$ agree. Define $f_g$ to be this common coefficient when such a $j$ exists, and zero otherwise.

We first check that $\supp f$ is well ordered. Let $U\subseteq\supp f$ be nonempty, and choose $g_0\in U$. There is $j_0$ with $g_0<\gamma_{j_0}$, and on all exponents below $\gamma_{j_0}$ the coefficients of $f$ agree with those of $a_{j_0}$. Thus the nonempty set $U\cap(-\infty,g_0]$ is a subset of the well-ordered support of $a_{j_0}$ and has a least element. That element is also least in $U$. Hence $f$ is a Hahn series.

Moreover,
\[
 \supp f\subseteq\bigcup_{j\in J}\supp a_j.
\]
The union on the right has size less than $\kappa$: a union of fewer than $\cf(\kappa)$ sets, each of cardinality less than $\kappa$, still has cardinality less than $\kappa$. Therefore $f\in\K$. For each $j$, it agrees with $a_j$ at every exponent below $\gamma_j$, so \eqref{fkc:eq:first-diff} gives $f\in B(a_j,\gamma_j)$.
\end{proof}

The gluing argument without the final cardinal bound is the standard proof of spherical completeness for full Hahn fields; compare \cite[Lemma~4]{Poonen}. The extra point here is exactly where the union of supports remains below $\kappa$.

\begin{theorem}[First empty nests]\label{fkc:thm:balls}
If $\K\ne F_k$, the least cardinality of an empty nest of closed valuation balls in $\K$ is $\mu=\cf(\kappa)$. In particular, $\K$ is not spherically complete.
\end{theorem}
\begin{proof}
The lower bound is \cref{fkc:lem:glue}. For the upper bound, choose a missing $x$ and use the approximants from \eqref{fkc:eq:approximants}. Let
\[
 B_i=B(a_i,s_{\alpha_i})\quad(i<\mu).
\]
For $i<j$, the center difference has value $s_{\alpha_i}$, and the radius value $s_{\alpha_j}$ is larger. The ultrametric inequality therefore gives $B_j\subseteq B_i$. The inclusion is strict: $a_i\notin B_j$. A common point would satisfy all formulas \eqref{fkc:eq:om-ball}, whose common realization in $\K$ was shown impossible. Thus this is an empty nest of cardinality exactly $\mu$.
\end{proof}

\begin{corollary}[Its intersection in the full Hahn field]\label{fkc:cor:ballfiber}
The balls just constructed, interpreted in $F_k$, have intersection
\[
 \bigcap_{i<\mu}B_{F_k}(a_i,s_{\alpha_i})
   =T_\kappa(x)+I_{D_x}.
\]
Thus the same coset is a fiber of the type map and a full-field intersection of a canonical omitted ball nest.
\end{corollary}
\begin{proof}
The ball conditions say precisely that all coefficients below the cofinal thresholds $s_{\alpha_i}$ agree with those of $x$. Their union forces exactly the first-$\kappa$ prefix and leaves precisely the tails in $I_{D_x}$ free.
\end{proof}

\begin{remark}[Not a full saturation theorem]
The ball theorem does not say that every type over fewer than $\mu$ parameters is realized. Our lower bound in \cref{fkc:thm:omission} applies to types realized in the fixed full Hahn extension. Other types may demand new values or residues. A field can have all the small ball intersections established here while failing a different saturation condition. Likewise, the pure-field omission number in \cref{fkc:thm:pure} is not contradicted by the smaller valued-field ball obstruction. \mergetag{} For the proper-class field $\KN\kappa\R$ of Part~II, which contains every monomial, saturation is proved: it is $\nu$-saturated exactly for $\nu\le\cf\kappa$ (\cref{fkc:sb:cor:saturation}). Nothing there changes the set-field statement of this remark.
\end{remark}

\section{The valuation completion and its cofinality criterion}
\label{fkc:sec:completion}

Completeness and spherical completeness differ because a Cauchy net must reach every valuation scale, whereas the radii of a nest can remain confined to a proper initial segment of the value group. In the present fields this distinction is determined exactly by the first-$\kappa$ support.

\subsection{An explicit completion inside the full Hahn field}
Define
\begin{equation}\label{fkc:eq:localfield}
 L_{\kappa,k}=\left\{f\in F_k:
 |\supp f\cap(-\infty,\gamma)|<\kappa
 \text{ for every }\gamma\in\Gamma\right\}.
\end{equation}
The condition is local in the \emph{exponent order}, not a condition of local finiteness in an ordinary real variable.

\begin{theorem}[Completion by locally small support]\label{fkc:thm:completion}
The closure of $\K$ in $F_k$ is $L_{\kappa,k}$. This is a complete valued subfield of $F_k$ and, with the given inclusion of $\K$, is its Hausdorff valuation completion.
\end{theorem}
\begin{proof}
Suppose $f$ is in the closure. Given $\gamma\in\Gamma$, there is $a\in\K$ with $v(f-a)\ge\gamma$. By \eqref{fkc:eq:first-diff}, the support of $f$ below $\gamma$ is contained in the support of $a$, so it has size less than $\kappa$. Thus $f\in L_{\kappa,k}$.

Conversely, if $f\in L_{\kappa,k}$, then $f_{<\gamma}\in\K$ for every $\gamma$. Moreover,
\[
 v(f-f_{<\gamma})\ge\gamma.
\]
Direct $\Gamma$ by its increasing order. These truncations converge to $f$: given a target value $\beta$, choose $\gamma>\beta$. Hence $f$ lies in the closure of $\K$.

The full Hahn field is complete; this follows from its classical spherical completeness, or directly from coefficient stabilization in a Cauchy net. For clarity, in a Cauchy net all coefficients up to any fixed exponent eventually agree simultaneously. Their eventual values define a Hahn series: any nonempty part of its proposed support has a least element by comparison below one selected exponent with a sufficiently late member of the net. The same stabilization proves convergence to that series.

The closure of a subfield in a topological field is a subfield: addition and multiplication are continuous, and inversion is continuous at each nonzero point. A closed subspace of the complete Hahn uniform space is complete. Since $\K$ is dense in its closure, this closure is its Hausdorff uniform completion, with the induced field and valuation structures.
\end{proof}

\begin{proposition}[Exactly which supports are added]\label{fkc:prop:new-supports}
For $f\in F_k\setminus\K$, the following are equivalent:
\begin{enumerate}[label=\textup{(\roman*)}]
\item $f\in L_{\kappa,k}$;
\item $\supp f$ has order type exactly $\kappa$ and is cofinal in $\Gamma$;
\item $D_f=\Gamma$.
\end{enumerate}
In particular, every new completion element satisfies $T_\kappa(f)=f$.
\end{proposition}
\begin{proof}
Write the support increasingly as $(s_\alpha)_{\alpha<\lambda}$ with $\lambda\ge\kappa$. If $\lambda>\kappa$, the exponent $s_\kappa$ exists, and the support strictly below it already has cardinality $\kappa$. This violates \eqref{fkc:eq:localfield}. Hence membership in $L_{\kappa,k}$ forces $\lambda=\kappa$.

If that support were not cofinal in $\Gamma$, choose $\gamma$ strictly above it. Then the support below $\gamma$ has size $\kappa$, again a contradiction. This proves (i)$\Rightarrow$(ii).

If the support has order type $\kappa$ and is cofinal, then for each $\gamma$ there is $\alpha<\kappa$ with $\gamma\le s_\alpha$. Its support below $\gamma$ is contained in the earlier $\alpha$ positions, so it has size less than $\kappa$. Thus (ii)$\Rightarrow$(i). Condition (ii) clearly implies (iii). Conversely, (iii) says that the first $\kappa$ support positions are already cofinal in $\Gamma$; there can be no later exponent $s_\kappa$ above them all. The full support therefore has exactly order type $\kappa$, proving (iii)$\Rightarrow$(ii).
\end{proof}

\subsection{The group-theoretic existence calculation}
\begin{lemma}[Long cofinal supports]\label{fkc:lem:long-cofinal}
Suppose $\Gamma$ contains a well-ordered subset of cardinality $\kappa$. Then it contains a cofinal well-ordered subset of order type $\kappa$ if and only if
\begin{equation}\label{fkc:eq:cofinal-condition}
 \cf(\Gamma)=\cf(\kappa).
\end{equation}
\end{lemma}
\begin{proof}
If a cofinal increasing sequence of order type $\kappa$ exists, its cofinality is $\cf(\kappa)$, and cofinal subsets have the same cofinality as the ambient linear order. More explicitly, a cofinal subsequence of $\kappa$ indices gives the upper bound, and a smaller cofinal subset of $\Gamma$ would, by choosing sequence entries above its members, give a smaller cofinal subset of $\kappa$. This proves necessity.

For sufficiency, put $\mu=\cf(\Gamma)=\cf(\kappa)$. Take an increasing sequence $(s_\alpha)_{\alpha<\kappa}$ in $\Gamma$. If it is already cofinal, there is nothing to prove. Otherwise, translating it by $-s_0$ gives a sequence $(u_\alpha)_{\alpha<\kappa}$ in an interval $[0,h)$ for some $h>0$.

If $\kappa$ is regular, then $\mu=\kappa$, and the definition of $\cf(\Gamma)$ gives a strictly increasing cofinal sequence of length $\kappa$. Such a sequence can be built by choosing, at each stage, a point above the fewer than $\cf(\Gamma)$ previous choices and above a prescribed cofinal seed.

Suppose now that $\kappa$ is singular. Choose increasing infinite cardinals $\kappa_i<\kappa$ for $i<\mu$ with supremum $\kappa$. Construct a cofinal increasing sequence $(q_i)_{i<\mu}$ in $\Gamma$ such that
\[
 q_j>q_i+h\qquad(i<j<\mu).
\]
At stage $j$, fewer than $\cf(\Gamma)$ previously selected values $q_i+h$ have been specified, so they are bounded. Choose the next value above them and above the next member of a fixed cofinal seed.

Now concatenate the blocks
\[
 \{q_i+u_\alpha:\alpha<\kappa_i\}\qquad(i<\mu).
\]
Different blocks are strictly separated, and their union is cofinal because it contains every $q_i$. Its order type is the ordinal sum $\sum_{i<\mu}\kappa_i$. This sum equals the initial ordinal $\kappa$. Indeed, every proper initial collection of blocks has cardinality less than $\kappa$, since it uses fewer than $\cf(\kappa)$ cardinals below $\kappa$; the same is true after adding a proper initial part of the next block. Thus every proper initial segment of the concatenation has cardinality less than $\kappa$, while the whole concatenation has cardinality $\kappa$. Its order type is therefore exactly $\kappa$.
\end{proof}

\begin{theorem}[Completeness dichotomy]\label{fkc:thm:dichotomy}
The following exhaustive alternatives hold.
\begin{enumerate}[label=\textup{(\roman*)}]
\item If $\Gamma$ has no well-ordered subset of cardinality $\kappa$, then $\K=F_k$, so $\K$ is complete and spherically complete.
\item If $\Gamma$ has such a subset, then $\K\ne F_k$ and
\begin{equation}\label{fkc:eq:complete-dichotomy}
 \boxed{\ \K\text{ is complete}\quad\Longleftrightarrow\quad
          \cf(\Gamma)\ne\cf(\kappa).\ }
\end{equation}
In this second alternative, the least size of an empty closed-ball nest is always $\cf(\kappa)$, regardless of completeness.
\end{enumerate}
\end{theorem}
\begin{proof}
A missing Hahn series yields a well-ordered support containing its first $\kappa$ positions. Conversely, an increasing support of order type $\kappa$ supports a Hahn series with all coefficients equal to one, which is missing from $\K$. This proves the criterion for properness.

By \cref{fkc:thm:completion,fkc:prop:new-supports}, a proper $\K$ fails to be complete exactly when $\Gamma$ admits a cofinal increasing sequence of order type $\kappa$. By \cref{fkc:lem:long-cofinal}, under the properness hypothesis this is equivalent to equality of the two cofinalities. The final assertion is \cref{fkc:thm:balls}.
\end{proof}

\begin{remark}[The completion need not be the full Hahn field]
A new completion element is required to have cofinal support of order type exactly $\kappa$. A series with $\kappa$ support positions confined below a fixed exponent is never added by completion. Thus the completion is not obtained simply by replacing $<\kappa$ by $\le\kappa$, nor by replacing $\kappa$ with $\kappa^+$. Example~\ref{fkc:ex:partialcompletion} gives an explicit strict chain $K_{\kappa,k}\subsetneq L_{\kappa,k}\subsetneq F_k$.
\end{remark}

\begin{remark}[The countable analogue elsewhere in the collection; added during assembly]\label{fkc:rem:countable-analogue}
The case $\kappa=\aleph_0$ is excluded here, since finite-support series do not form a field (\cref{fkc:app:checks}). Two other reports prove the countable analogue of \cref{fkc:thm:completion,fkc:thm:dichotomy} for other objects. In the report on three duals of Hahn vector spaces, the completion of the algebraic scalar extension $E_\Gamma(V)$ of an infinite-dimensional $k$-vector space $V$ consists of the vectors with finite coefficient rank below every $\gamma\in\Gamma$ (\texttt{duals:thm:completion}), and it is strictly larger than $E_\Gamma(V)$ exactly when $\cf(\Gamma)=\aleph_0$ (\texttt{duals:thm:cofinality-completion}). This is the shape of \eqref{fkc:eq:localfield} and \eqref{fkc:eq:complete-dichotomy} with ``finite coefficient rank'' in place of ``fewer than $\kappa$ support terms'' and $\aleph_0=\cf(\aleph_0)$ in place of $\cf(\kappa)$. In the rank-one Berkovich report, the metric closure of the finite-support series in $\C((t^\Q))$ consists of the series whose support meets every $(-\infty,B]$ in a finite set, which is \eqref{fkc:eq:localfield} with $\kappa$ replaced by $\aleph_0$. Neither result is used here, and neither implies nor is implied by the theorems above: the objects there are a vector space over the full Hahn field and the ring of finite-support series, and the bound is countable.
\end{remark}

\section{An exact two-dimensional linear obstruction}
\label{fkc:sec:linear}

The same approximation cut has a direct functional interpretation. Fix either coefficient field and abbreviate $K=\K$, $F=F_k$. For $x\in F\setminus K$, let
\[
 E_x=K+Kx\subseteq F
\]
with the valuation inherited from $F$. The vectors $1,x$ are linearly independent over $K$, so $E_x$ is two dimensional. This is a $\Gamma$-valued linear space, not an asserted real-normed Banach space.

Every $K$-linear retraction $P:E_x\to K$ of the inclusion $K\subseteq E_x$ has the form
\begin{equation}\label{fkc:eq:Pc}
 P_c(a+bx)=a+bc\qquad(a,b\in K)
\end{equation}
for a unique $c\in K$. For $\delta\in\Gamma_{\ge0}$, say that it has \emph{valuation loss at most $\delta$} when
\begin{equation}\label{fkc:eq:loss}
 v(P_cz)\ge v(z)-\delta\qquad(z\in E_x).
\end{equation}
Loss zero is the nonexpanding, or contractive, case in valuation notation.

\begin{theorem}[Exact retraction-loss formula]\label{fkc:thm:projection}
Put $r_c=v(x-c)\in D_x$. Then
\begin{equation}\label{fkc:eq:projectionexact}
 \boxed{\ P_c\text{ satisfies \eqref{fkc:eq:loss}}
 \quad\Longleftrightarrow\quad
 d\le r_c+\delta\text{ for every }d\in D_x.\ }
\end{equation}
Consequently no $P_c$ has loss zero. A continuous retraction exists if and only if $D_x$ is bounded above in $\Gamma$, equivalently $x\notin L_{\kappa,k}$.
\end{theorem}
\begin{proof}
Vectors in $K$ satisfy \eqref{fkc:eq:loss} automatically. A vector $a+bx$ with $b\ne0$ is $b(x-d)$ for $d=-a/b\in K$. By linearity and multiplicativity of the valuation, it suffices to check
\begin{equation}\label{fkc:eq:reducedloss}
 v(c-d)\ge v(x-d)-\delta\qquad(d\in K).
\end{equation}
Write $r_d=v(x-d)$. If $r_d<r_c$, then \eqref{fkc:eq:ultra-exact} gives $v(c-d)=r_d$, so \eqref{fkc:eq:reducedloss} holds. If $r_d=r_c$, the ultrametric inequality gives $v(c-d)\ge r_c$, so it also holds. If $r_d>r_c$, the exact value is $v(c-d)=r_c$, and \eqref{fkc:eq:reducedloss} is equivalent to $r_d\le r_c+\delta$. As the set of all $r_d$ is $D_x$, this is exactly \eqref{fkc:eq:projectionexact}.

There is no zero-loss retraction because $D_x$ has no greatest element. If $D_x$ is bounded above by $b$, then for any $c$ the nonnegative value $\delta=b-r_c$ is an admissible loss. Such a bound implies continuity directly.

Conversely, any continuous $K$-linear map $P:E_x\to K$ has some uniform valuation-loss bound. Indeed, continuity at zero supplies $\beta\in\Gamma$ such that $v(z)>\beta$ implies $v(Pz)>0$. Choose $\eta>0$. For $z\ne0$, multiplication by the scalar
\[
 q=t^{\beta+\eta-v(z)}\in K
\]
makes $v(qz)=\beta+\eta>\beta$. If $Pz\ne0$, linearity yields
\[
 v(Pz)>v(z)-\beta-\eta.
\]
The same lower bound is automatic if $Pz=0$. Thus $\delta=\max\{0,\beta+\eta\}$ is a loss bound. Applying \eqref{fkc:eq:projectionexact} to a continuous retraction forces $D_x$ to be bounded above. Finally, an initial segment of a linear order is unbounded exactly when it is the whole order, and \cref{fkc:prop:new-supports} identifies $D_x=\Gamma$ with $x\in L_{\kappa,k}$.
\end{proof}

\begin{corollary}[A zero-or-full continuous-dual dichotomy]\label{fkc:cor:dual}
Let $E_x'=\operatorname{Hom}_{K,\mathrm{cont}}(E_x,K)$. Then
\[
 \dim_K E_x'=
 \begin{cases}
 0,&x\in L_{\kappa,k}\setminus K,\\
 2,&x\in F_k\setminus L_{\kappa,k}.
 \end{cases}
\]
In the second case every algebraic linear functional is continuous, although no retraction has valuation loss zero.
\end{corollary}
\begin{proof}
If $x\in L_{\kappa,k}\setminus K$, then $K$ is dense in $E_x$: approximate $x$ by elements of $K$ and multiply and translate to approximate $a+bx$. If a continuous functional $\ell$ had $\ell(1)\ne0$, dividing by $\ell(1)$ would produce a continuous retraction, contrary to \cref{fkc:thm:projection}. Hence $\ell$ vanishes on the dense subspace $K$ and therefore vanishes everywhere, since the target is Hausdorff.

Suppose instead that $D_x$ is bounded. Fix $c\in K$ and a loss bound $\delta$ for $P_c$. The coefficient functional $Q(a+bx)=b$ is continuous as well. For $z=a+bx$, the identity $z-P_cz=b(x-c)$ gives, whenever $b\ne0$,
\[
 v(b)=v(z-P_cz)-r_c
 \ge\min\{v(z),v(P_cz)\}-r_c
 \ge v(z)-\delta-r_c.
\]
Replacing the right-hand loss by a larger nonnegative value if necessary gives a uniform continuity bound. Since $P_c$ and $Q$ span the algebraic dual of $E_x$, every linear functional is continuous and the dual dimension is two.
\end{proof}

This application does not claim a new general Hahn--Banach theorem. The mechanism is the elementary best-approximation obstruction: contractivity of $P_c$ would require $c$ to approximate $x$ at least as well as every other base-field element. The contribution within the present package is that the full set of permitted losses, and the continuous-dual dichotomy, are explicitly calculated from the first-$\kappa$ support cut.

\section{Explicit examples at regular and singular cardinals}
\label{fkc:sec:examples}

\subsection{A concrete family of exponent groups}
For an infinite limit ordinal $\theta$, let
\begin{equation}\label{fkc:eq:Gtheta}
 \Gamma_\theta=\bigoplus_{\alpha<\theta}\Q e_\alpha
\end{equation}
with finite supports and the order determined by the \emph{largest} nonzero coordinate. Thus $0<e_\alpha\ll e_\beta$ for $\alpha<\beta$: every positive integral multiple of $e_\alpha$ is less than $e_\beta$. This is a divisible ordered abelian group.

\begin{lemma}\label{fkc:lem:groupcf}
For an infinite limit ordinal $\theta$,
\[
 \cf(\Gamma_\theta)=\cf(\theta).
\]
Moreover, $(e_\alpha)_{\alpha<\theta}$ is increasing and cofinal in $\Gamma_\theta$.
\end{lemma}
\begin{proof}
Every group element has finite support. An $e_\beta$ with index larger than all indices in that support exceeds the group element, proving cofinality of the displayed sequence. Taking a cofinal subset of indices gives $\cf(\Gamma_\theta)\le\cf(\theta)$. Conversely, a set of fewer than $\cf(\theta)$ group elements uses a bounded set of coordinate indices: it is a union of fewer than $\cf(\theta)$ finite sets. Some later $e_\beta$ exceeds them all, so the set cannot be cofinal. This proves the reverse inequality.
\end{proof}

\subsection{A singular bound with a countable missing Cauchy limit}
\begin{example}\label{fkc:ex:singular}
Let $\kappa=\aleph_\omega$, $\Gamma=\Gamma_\kappa$, and take either $k=\R$ or $\C$. Define
\begin{equation}\label{fkc:eq:xkappa}
 x_\kappa=\sum_{\alpha<\kappa}t^{e_\alpha},\qquad
 a_n=\sum_{\alpha<\aleph_n}t^{e_\alpha}\quad(n<\omega).
\end{equation}
Each $a_n$ belongs to $K_{\kappa,k}$, whereas $x_\kappa$ does not. Furthermore,
\[
 v(x_\kappa-a_n)=e_{\aleph_n},\qquad
 v(a_m-a_n)=e_{\aleph_n}\quad(n<m).
\]
These values are cofinal in $\Gamma_\kappa$, so $(a_n)$ is a Cauchy sequence with limit $x_\kappa$ in the full Hahn field and no limit in $K_{\kappa,k}$.
\end{example}

All algebraic equations appropriate to real closedness or algebraic closedness can still be solved in the smaller field, by \cref{fkc:prop:closed}. The omitted type of $x_\kappa$ in the ordered real or valued complex language needs exactly countably many formulas. Its realization set in this particular full Hahn field is a singleton, because $D_{x_\kappa}=\Gamma_\kappa$ and $I_{D_{x_\kappa}}=\{0\}$. This does not mean the type is isolated by a single formula: \cref{fkc:thm:omission} proves the opposite.

For the two-dimensional valued space $K+Kx_\kappa$, the continuous $K$-linear dual is zero. Thus even two algebraic coordinates need not supply continuous coordinate functionals over this incomplete base field with its inherited valuation.

In fact, the completion in \cref{fkc:ex:singular} is the entire full Hahn field. Every bounded interval in $\Gamma_\kappa$ has cardinality less than $\kappa$. To see this, choose an index $\beta<\kappa$ at least as large as all coordinate indices in its two endpoints. A group element with a nonzero coordinate beyond $\beta$ lies strictly above both endpoints or strictly below both, according to the sign of its last coordinate. The interval is consequently contained in the rational span of $\{e_\alpha:\alpha\le\beta\}$, a set of size less than $\kappa$. Since a nonempty Hahn support has a least exponent, each initial slice of it is contained in such a bounded interval. Thus every Hahn series has locally small support, and \cref{fkc:thm:completion} gives $L_{\kappa,k}=F_k$.

\subsection{The same countable obstruction in a complete field}
\begin{example}\label{fkc:ex:complete}
Keep $\kappa=\aleph_\omega$ but enlarge the exponent group to $\Gamma_{\kappa^+}$. The same series $x_\kappa$ in \eqref{fkc:eq:xkappa} now has support bounded above by $e_\kappa$. Since
\[
 \cf(\Gamma_{\kappa^+})=\kappa^+\ne\omega=\cf(\kappa),
\]
\cref{fkc:thm:dichotomy} says that $K_{\kappa,k}$ is complete. Nevertheless, the countable nest
\[
 B(a_n,e_{\aleph_n})\qquad(n<\omega)
\]
has empty intersection in $K_{\kappa,k}$. Its radii never reach the scale $e_\kappa$, so it is not a Cauchy obstruction. In particular, the sequence $(a_n)$ itself is not Cauchy in this larger valuation topology.
\end{example}

For this example, all retractions $K+Kx_\kappa\to K$ are continuous but none is contractive. Moreover,
\[
 x_\kappa\quad\text{and}\quad x_\kappa+t^{e_\kappa}
\]
are distinct elements with exactly the same ordered real type or complex valued type over $K_{\kappa,k}$. The coefficient at $e_\kappa$ is part of the invisible tail. Its addition changes the full support order type from $\kappa$ to $\kappa+1$ but leaves the first-$\kappa$ prefix unchanged.

\subsection{A completion strictly between the two support fields}
\begin{example}\label{fkc:ex:partialcompletion}
Let $\kappa=\aleph_\omega$ and take the ordered rational vector space
\[
 \Gamma^*=\Gamma_\kappa\oplus\Q u,
\]
where the $u$ coordinate dominates every coordinate of $\Gamma_\kappa$. Its cofinality is $\omega$, since $(nu)_{n\ge1}$ is cofinal. In $F_k=k((t^{\Gamma^*}))$, the series $x_\kappa=\sum_{\alpha<\kappa}t^{e_\alpha}$ has its entire support below $u$, so it does not belong to the completion of $K_{\kappa,k}$. On the other hand,
\[
 y=\sum_{n<\omega}\ \sum_{\alpha<\aleph_n}t^{nu+e_\alpha}
\]
has a cofinal well-ordered support of order type $\sum_{n<\omega}\aleph_n=\aleph_\omega$. The blocks are strictly separated because every $e_\alpha$ is less than $u$. Therefore $y\in L_{\kappa,k}\setminus K_{\kappa,k}$ by \cref{fkc:prop:new-supports}, whereas $x_\kappa\in F_k\setminus L_{\kappa,k}$. This proves the strict chain
\[
 K_{\kappa,k}\subsetneq L_{\kappa,k}\subsetneq F_k.
\]
\end{example}

\subsection{A regular bound with a genuinely uncountable obstruction}
\begin{example}\label{fkc:ex:regular}
Take $\kappa=\aleph_1$ and $\Gamma=\Gamma_{\omega_1}$. Every countable nest of closed balls in $K_{\aleph_1,k}$ meets, but a nest of cardinality $\aleph_1$ can be empty. The field is not complete, because $\cf(\Gamma)=\aleph_1=\cf(\kappa)$, and
\[
 x=\sum_{\alpha<\omega_1}t^{e_\alpha}
\]
is a missing limit of its $\omega_1$-indexed truncation net. Its ordered or valued type has omission number $\aleph_1$.
\end{example}

Here every Cauchy \emph{sequence} is eventually constant. To prove this directly, consider the countable set of finite values $v(z_m-z_n)$ of all nonzero pairwise differences in a given sequence. Since $\cf(\Gamma)>\omega$, this set is bounded above by some $\beta\in\Gamma$. The Cauchy condition at scale $\beta$ forces every sufficiently late pairwise difference to be zero. Thus sequential completeness does not imply completeness for the valuation uniformity in this example. The missing net cannot be replaced by a countable Cauchy sequence.

\subsection{The rank-one boundary case}
If $\Gamma$ is a subgroup of $(\R,+)$, every well-ordered subset of $\Gamma$ is countable. One proof is to assign to each nonfinal element of a well-ordered subset a rational number strictly between it and its successor; the chosen intervals are disjoint. There are only countably many rationals. Consequently, for every uncountable $\kappa$,
\[
 K_{\kappa,k}=k((t^\Gamma)).
\]
The nontrivial cardinal-support phenomena in this paper therefore require value groups beyond the Archimedean case. Merely replacing the real coefficients by complex ones in a rank-one Hahn field does not create these missing types. This is consistent with the rank-one Berkovich report, where the proper closure of the finite-support series in $\C((t^\Q))$ occurs only at the excluded countable bound (\cref{fkc:rem:countable-analogue}).

\begin{remark}[The same groups elsewhere in the collection; added during assembly]\label{fkc:rem:groups-elsewhere}
The construction \eqref{fkc:eq:Gtheta}, together with the embedding $e_\alpha\mapsto\omega^\alpha$ of \cref{fkc:sec:surreal}, is the one used in the report on entire functions at arbitrary rank. Its group $\Gamma_\infty=\bigoplus_{j\ge0}\Q e_j$ with $e_j=\omega^j$ (\texttt{ent:eq:concrete-group}) is $\Gamma_\omega$ here, and its group $\Gamma_{\omega_1}$ (\texttt{ent:eq:Gamma-omegaone}) is the group of \cref{fkc:ex:regular}. Both reports compute the same cofinalities, $\aleph_0$ and $\aleph_1$. The questions asked about these groups are different: there, entire functions and their rings; here, cardinally bounded subfields.
\end{remark}

\section{Transfer to actual surreal and surcomplex numbers}
\label{fkc:sec:surreal}

\subsection{The normal-form realization}
Let $j:\Gamma\to(\No,+)$ be a specified increasing additive embedding of the set-sized exponent group. Conway normal-form arithmetic gives the field embedding
\begin{equation}\label{fkc:eq:conway}
 \iota_j:\R((t^\Gamma))\longrightarrow\No,\qquad
 \sum_{\gamma\in S}r_\gamma t^\gamma
 \longmapsto
 \sum_{\gamma\in S}r_\gamma\omega^{-j(\gamma)}.
\end{equation}
The exponent set on the right is reverse well ordered, since $S$ is well ordered and $j$ is increasing. The minus sign is essential: the increasing-support convention for $t^\Gamma$ corresponds to decreasing exponents of $\omega$. Existence, uniqueness, and arithmetic of these Conway normal forms are the standard surreal normal-form theorem \cite{Gonshor}.

The map is injective, order preserving, and preserves the support cardinality and its indexing order type after reversing the exponent convention. For complex coefficients, use the coordinatewise embedding
\begin{equation}\label{fkc:eq:complexconway}
 \C((t^\Gamma))\cong\R((t^\Gamma))[i]
 \longrightarrow\No[i].
\end{equation}
Its two real coordinates are the real and imaginary coefficient series. Their union is exactly the support of the original complex series, so the cardinal restriction transfers without change.

Consequently every theorem in this paper applies to the actual subfields $\iota_j(K_{\kappa,\R})\subseteq\No$ and their quadratic complexifications, when furnished with the transported valued or ordered structures. No proper-class quantification is required for the one-types: their base fields and ambient full Hahn workspaces remain sets.

\subsection{The explicit groups already live in the surreal field}
For the groups $\Gamma_\theta$ in \eqref{fkc:eq:Gtheta}, define
\[
 j\left(\sum_{\alpha\in A}q_\alpha e_\alpha\right)
    =\sum_{\alpha\in A}q_\alpha\omega^\alpha,
 \qquad A\subseteq\theta\text{ finite}.
\]
Finite Conway normal forms make this an additive embedding, and the largest exponent controls the sign, so it has exactly the required order. Thus $e_\alpha$ can be represented concretely by the surreal monomial $\omega^\alpha$.

Under this embedding, the missing series in \eqref{fkc:eq:xkappa} is the actual surreal
\begin{equation}\label{fkc:eq:actualx}
 \boxed{\ X_\kappa=\sum_{\alpha<\kappa}\omega^{-\omega^\alpha}.\ }
\end{equation}
For $\kappa=\aleph_\omega$, its approximants are
\[
 A_n=\sum_{\alpha<\aleph_n}\omega^{-\omega^\alpha}.
\]
In the workspace with exponent group $j(\Gamma_\kappa)$, these are a Cauchy sequence in the intrinsic valuation topology and have no limit in the support-restricted subfield. In the workspace with exponent group $j(\Gamma_{\kappa^+})$, the same displayed approximants are not Cauchy, and the support-restricted field is complete. The smaller scale $\omega^{-\omega^\kappa}$ now exists in the workspace and detects the failure of the Cauchy condition.

This is a precise example of why a surreal limit assertion must specify which scales are permitted in its neighborhoods. The formulas for the numbers alone do not determine a topology on an arbitrarily chosen subfield.

\subsection{The topology and language boundaries}
The topological statements here concern the valuation taking its finite values in the chosen group $\Gamma$. They are not statements about the topology induced from the order topology of the full class $\No$. In the latter setting, a set of nonzero differences has a positive surreal lower bound for its absolute values, so every set-sized subset is discrete. The repository's foundations treatment explicitly distinguishes this full fine topology from an intrinsic Hahn-workspace topology \cite{Repo} (\texttt{found:prop:twotopologies} in \cite{CollFound}). In particular, convergence of the displayed sequence in one workspace does not contradict the eventual-constancy results for small-index nets in the full fine topology.

Similarly, the support-restricted subfields are not asserted to be birthday-initial segments. No closure under the full surreal exponential, no compatibility with the Berarducci--Mantova derivation, and no definability of the prefix operator in the field language is claimed. The prefix is an external invariant describing the types; the individual tests used in the proofs are ordinary finitary formulas with allowed parameters.

\emph{Status (batch 32).} These disclaimers stand for the set-sized subfields $\iota_j(K_{\kappa,\R})$, which in general contain only the monomials $\omega^{-j(\gamma)}$ and are not closed under the exponential. For the proper-class field $\KN\kappa\R$ of all surreals with fewer than $\kappa$ terms, which contains every monomial, \cref{fkc:sb:part:class} proves closure under the surreal exponential but not under the logarithm (\cref{fkc:sb:thm:expclosure,fkc:sb:prop:expimage}), and the nonexistence of any surjective ordered exponential (\cref{fkc:sb:thm:noexp}). Birthday-initiality, the derivation (\cref{fkc:sb:q:derivations}) and definability of the prefix operator remain unaddressed there too.

Naming complex conjugation would create another language. Its fixed field is the real axis, and positivity on that fixed field is definable by being a nonzero square. Thus the pure-field collapse from \cref{fkc:thm:pure} cannot be transferred without qualification to a surcomplex field with named conjugation. We have not classified all one-types in that expanded language: an arbitrary complex element corresponds to a pair of real coordinates, whose joint algebraic relations require a tuple analysis.

\emph{Status (batch 32).} For the class fields of \cref{fkc:sb:part:class}, the structures $(\KN\kappa\C,\conj_\kappa)$ with named conjugation are all elementarily equivalent to $(\No[i],\conj)$, by interpretation in the real closed fixed field, and are $\nu$-saturated exactly for $\nu\le\cf\kappa$ (\cref{fkc:sb:prop:sametheory}). This does not classify the one-types in that language over the set fields $\KC$, which remain unclassified.

\section{Scope, provenance, and a formalization route}
\label{fkc:sec:audit}

\subsection{Repository boundary}
The repository was examined at the pinned commit
\begin{center}\small\repoid.\end{center}
The root overview, documentation catalogue, the \texttt{docs/new} directory, and the research audit of the closest bounded-support transcendence report were read through the GitHub connector. The catalogue at that revision describes 44 reports and explicitly distinguishes mathematical arguments, finite checks, and Lean coverage \cite{Repo}. No complete audit of every article, historical version, or Lean declaration is claimed.

Three comparisons matter. The report on \emph{transcendence over bounded support} uses supports bounded in exponent order and the fraction field they generate; our base is instead the already-real-closed or algebraically closed field of all supports of cardinality less than $\kappa$. The foundations report treats universes, localization, and the full fine topology; our theorems concern exact one-types and intrinsic completion in a fixed set-sized localization. The Hahn-vector-space report studies strong maps and several duals; \cref{fkc:cor:dual} instead uses the inherited valuation on a two-dimensional space over a potentially incomplete cardinally bounded base field. These differences are substantive, not merely changes of notation.

A targeted repository search returned no indexed matches for one relevant term, but the index was not treated as an exhaustive content audit. The comparisons above rely on the retrieved catalogue and the closest report's stated base field and scope, not on that negative search.

\subsection{Literature and novelty ledger}
The mathematical dependency boundary can be summarized without claiming publication priority for standard ingredients.

\begin{center}
\small
\begin{tabular}{@{}p{0.34\textwidth}p{0.60\textwidth}@{}}
\toprule
Ingredient or conclusion & Status in this manuscript \\
\midrule
Hahn fields; normal-form realization & Classical input; not new \cite{Poonen,Gonshor}.\\[0.4em]
Regular $\kappa$-bounded constructions & Prior literature; not new \cite{KS,BKMM}.\\[0.4em]
Quantifier elimination & Classical logical input; not new \cite{Tarski,Robinson,Yin}.\\[0.4em]
Approximation types and ball conditions & General framework already established \cite{FK}.\\[0.4em]
First-$\kappa$ type fibers; exact omission number & Explicit calculation proved here, including singular $\kappa$; proposed contribution.\\[0.4em]
Completion and first empty nests & Sharp cardinal-support classification proved here; gluing mechanism is classical.\\[0.4em]
Retraction loss and dual dichotomy & Direct consequences of the exact approximation cut; no general new Hahn--Banach theorem claimed.\\
\bottomrule
\end{tabular}
\end{center}

Kuhlmann's general approximation-type account is the closest logical precedent, especially its discussion of one-types and its ball formulas \cite[Section~4.2]{FK}. Merely using such ball conditions would not establish novelty. Here the cardinal computation supplies the explicit prefix parameterization, proves the lower bound for \emph{arbitrary} formulas, and connects it to a completion criterion for both regular and singular bounds. These are the statements proposed for further expert review.

The search was targeted rather than exhaustive. It included cardinally bounded Hahn fields, approximation types, singular-cardinal support bounds, completion, and spherical completeness. Some broad queries returned irrelevant results; those were not counted as evidence of absence. No exhaustive MathSciNet or zbMATH priority check, citation-network review, or survey of unpublished manuscripts was performed. Accordingly, the warranted claim is that the combined statements were not identified in the reviewed material. A later identification of an earlier equivalent theorem would change the novelty assessment, not replace the proofs provided here.

\subsection{A feasible proof-assistant decomposition}
No Lean file or Lean build accompanies this manuscript. A formalization would be most naturally divided into the following interfaces rather than beginning with proper-class model theory.

\paragraph{Coefficient layer.}
Work with a set-sized ordered exponent type and a Hahn-series type. Formalize strict and non-strict truncations, first disagreement, the cardinality of finite unions of supports, and the small-workspace lemma. The argument at singular $\kappa$ needs the countable cardinal bound on the rational span, not regularity.

\paragraph{Cardinal and support layer.}
Represent the first-$\kappa$ prefix by an order isomorphism from a well-ordered support to its ordinal order type. Prove the exact approximation set, prefix fibers, coefficient gluing, and the long-cofinal-support lemma. These conclusions do not depend on quantifier elimination. A useful interface theorem is that fewer than $\cf(\kappa)$ subsets of size less than $\kappa$ have union of size less than $\kappa$.

\paragraph{Algebra and topology layer.}
Import full-Hahn real or algebraic closedness with its hypotheses explicitly checked. Then formalize completion as closure inside the full Hahn field, the classification of newly added supports, and the retraction-loss trichotomy. This already certifies the main valuation and linear conclusions without a theory-of-types library.

\paragraph{Logical layer.}
State the quantifier-elimination interfaces for RCF and ACVF and prove eventual agreement first for atomic polynomial data, then for Boolean combinations, and finally for formulas. This isolates the only model-theoretic import in the omission-number lower bound. An implementation that assumes these interfaces proves a conditional theorem; it must not be described as an unconditional Lean verification until those imports are discharged.

\subsection{Further questions, not conclusions}
The one-variable invariant is complete precisely because over an algebraically closed base every polynomial splits into linear factors, while real polynomials add only positive quadratic factors. In several variables there is no analogous factorization into conditions on single coordinates. A natural next question is whether finite tuples in a full Hahn extension admit an effective prefix description that also records cancellations in polynomial combinations. The one-type theorem alone does not answer it.

A second question concerns richer languages: named truncation, named conjugation, a surreal exponential, or a derivation. Each can make additional relations visible, and closure of the base field under the new operations must be established before even formulating an elementary-substructure comparison. The present language-sensitive examples supply test cases, not a claim that the same classification survives those expansions.

\emph{Status (batch 32): partly addressed in the class setting, open here.} For the proper-class field $\KN\kappa\R$ of \cref{fkc:sb:part:class}, closure under the surreal exponential is established (\cref{fkc:sb:thm:expclosure}), and the exponential comparison with $\No$ is then shown not to be elementary (\cref{fkc:sb:cor:expsep}); with named conjugation, all the class structures $(\KN\kappa\C,\conj_\kappa)$ have one complete theory (\cref{fkc:sb:prop:sametheory}). No type classification in these expanded languages is obtained, for the class fields or for the set fields $\K$; named truncation and derivations are not treated. The question stays open as asked.

\subsection{Corrections to the manuscript's repository statements}
\label{fkc:subsec:corrections}
The statements above describe the repository at the manuscript's pin \texttt{465a54b}, which is kept as provenance; the pinned URLs in \cite{Repo} are also kept. During assembly they were checked against the current collection, with the following results.
\begin{enumerate}[label=\textup{(\arabic*)}]
\item \emph{Catalogue size.} At the pin the catalogue described 44 reports, as stated. Since then batch~21 added \texttt{surreal-fields-across-universes} and \texttt{birthday-cutoffs-and-hereditary-sets}, and batch~22 added this report and \texttt{single-dilation-hahn-support}; the collection had 48 report directories when this report was written.
\item \emph{Universes.} At the pin, universes were treated in the foundations report. The collection now also has a report on surreal fields across set-theoretic universes, which studies saturation and omitted cuts of proper-class surreal fields of inner models. Its objects differ from the set-sized fields here; see \cref{fkc:subsec:collection}.
\item \emph{The negative search.} The indexed search that returned no matches was for the word ``spherical'' (\texttt{research\_audit.md}). The negative result was wrong: at the pin the collection already stated spherical completeness of full Hahn fields, in the rank-one Berkovich report (its \texttt{prop:spherical}) and in the three-duals report (\texttt{duals:prop:spherical}). The manuscript did not rely on that search, as it says, and both statements are the classical full-Hahn result that it credits to Poonen.
\item \emph{The Hahn-vector-space comparison was incomplete.} Besides strong maps and duals, the three-duals report already contained, at the pin, the completion formula \texttt{duals:thm:completion} and the cofinality criterion \texttt{duals:thm:cofinality-completion}, and the rank-one Berkovich report the closure of the finite-support series in $\C((t^\Q))$. These are countable analogues of \cref{fkc:thm:completion,fkc:thm:dichotomy} (\cref{fkc:rem:countable-analogue}). They are not cardinally bounded fields and do not anticipate the uncountable, and in particular the singular, case treated here, but they belong in the priority context of the completion results.
\item \emph{Statements found accurate.} The description of the bounded-support report (order-bounded supports and their fraction field, explicitly not a cardinal restriction), the foundations report's separation of the fine and intrinsic topologies (\texttt{found:prop:twotopologies}), and the absence from the collection of any study of cardinally bounded Hahn fields, their types or their completions were confirmed.
\item \mergetag{} \emph{Later additions (batch 32).} The last confirmation describes the collection when this part was assembled. Since then the large-cardinal report has gained the field of short normal forms $\mathscr S_{<\kappa}$, which is the class field $\KN\kappa\R$ of Part~II, with its real closedness for every uncountable $\kappa$ (\texttt{lce:sg:thm:shortfield}); the quotient report has gained the bounded-support rings $\mathcal A^{<\lambda}_{D,k}$ and their support threshold (\texttt{osq:thm:support}); and the omnific-automorphisms report has gained the set-sized field $K_{<\kappa}$ for regular $\kappa$ with a summability criterion (\texttt{opa:as:prop:kappa}). None of them studies the types or completions of this part. Part~II adds the class version; see \cref{fkc:sb:sub:relation}.
\end{enumerate}

\subsection{What is not claimed}
\label{fkc:subsec:nonclaims}
The manuscript states its limitations where they arise. They are collected here, with the one limitation added during assembly; none is withdrawn.
\begin{enumerate}[label=\textup{N\arabic*.}, leftmargin=2.6em]
\item $\kappa$ is uncountable; the finite-support case $\kappa=\aleph_0$ is excluded (\cref{fkc:app:checks}).
\item Established mathematics is credited, not claimed: cardinally bounded Hahn fields (Kuhlmann--Shelah for regular $\kappa$; Berarducci--Kuhlmann--Mantova--Matusinski), Hahn arithmetic and closedness, Conway normal forms, general approximation-type theory including its ball formulas (F.-V.~Kuhlmann, Proposition~4.4 and Section~4.2), full-Hahn ball gluing (Poonen), and quantifier elimination for real closed and algebraically closed valued fields (Tarski, Robinson, Yin).
\item The contribution is proposed only: the joint first-$\kappa$ classification, the exact omission number and the completion and linear consequences. The combined statements were not identified in the reviewed material; the priority search was targeted, not exhaustive (no MathSciNet or zbMATH search), and expert review may find equivalents or a simpler general formulation.
\item No named longstanding conjecture is settled. The repository comparison was not an exhaustive audit, and a negative index search was not treated as evidence (it was in fact a false negative; \cref{fkc:subsec:corrections}). There is no Lean build or certification.
\item Only types realized in the fixed full Hahn extension $F_k$ are classified. No saturation theorem and no automorphism homogeneity is proved, and types are not asserted to be automorphism orbits of $F_k$. \emph{Status (batch 32):} this stands for the set fields $\K$; \cref{fkc:sb:cor:saturation} proves the saturation theorem for the class field $\KN\kappa\R$ of Part~II.
\item All topology is intrinsic to the set-sized workspace, not the subspace topology from $\No$; support bounds are cardinalities, not order bounds. \emph{Status (batch 32):} this stands for Part~I; Part~II treats the topology relative to $\No$ for the class fields (\cref{fkc:sb:prop:cauchy,fkc:sb:prop:closed}).
\item The subfields are not asserted to be birthday-initial. No closure under the surreal exponential, no compatibility with the Berarducci--Mantova derivation, and no definability of the prefix operator is claimed. \emph{Status (batch 32):} this stands for the set-sized subfields; for the class field $\KN\kappa\R$, Part~II proves closure under the exponential and not under the logarithm (\cref{fkc:sb:thm:expclosure}); the other three disclaimers stand there too.
\item Languages with named conjugation, truncation, exponential or derivation are not classified, nor are types of several-variable tuples; these remain questions. \emph{Status (batch 32):} still not classified; for the class fields Part~II computes the complete theory with named conjugation and shows the exponential inclusion is not elementary (\cref{fkc:sb:prop:sametheory,fkc:sb:cor:expsep}).
\item No new general Hahn--Banach theorem is claimed. The zero continuous dual is over an incomplete base in the inherited topology and does not contradict finite-dimensional functional analysis.
\item This is not a new version of the independence theorem of the report on transcendence over bounded support.
\item The 9,622 finite assertions of \texttt{code/verify.py} test finite local identities only. They do not verify the transfinite theorems, quantifier elimination, novelty or any Lean implementation.
\item Added during assembly: the comparisons with other reports in \cref{fkc:subsec:collection,fkc:rem:countable-analogue,fkc:rem:groups-elsewhere} are editorial. They prove no new theorem, and no statement here is derived from those reports.
\end{enumerate}

\section{Conclusion}
The obstruction created by a cardinal support bound has a precise shape. A missing Hahn series is first-order visible through its first $\kappa$ nonzero terms, but the least number of formulas needed to force that missing prefix is only $\cf(\kappa)$. Its approximation values are the downward closure of those first support positions. Whether this cut reaches all of $\Gamma$ decides both membership in the completion and the existence of continuous retractions from the associated two-dimensional valued space.

This gives a unified explanation of three superficially different effects: singular-cardinal compression of omitted types, complete fields with small empty ball nests, and failure of continuous coordinate functionals over an incomplete base. All three occur in explicitly specified surreal and surcomplex subfields. Their formulation requires neither a proper-class first-order structure nor an unspecified notion of surreal convergence.

\part{Support-bounded surreal fields: the class version (source 02)}
\label{fkc:sb:part:class}

\section{Source 02, its relation to Part I, and the conventions of this part \mergetag}
\label{fkc:sb:sec:intro}

This section, written for the merge, introduces the second manuscript of the report. Everything in \cref{fkc:sb:sec:setup,fkc:sb:sec:fields,fkc:sb:sec:gaps,fkc:sb:sec:saturation,fkc:sb:sec:spherical,fkc:sb:sec:topology,fkc:sb:sec:exponential,fkc:sb:sec:omnific,fkc:sb:sec:realforms,fkc:sb:sec:examples,fkc:sb:sec:questions,fkc:sb:sec:provenance,fkc:sb:sec:conclusion} and \cref{fkc:sb:app:audit,fkc:sb:app:realization} is source~02's, in the notation fixed in \cref{fkc:sb:sub:notation}; text added in the merge is marked \mergetag.

\subsection{Provenance of Part II}
\label{fkc:sb:sub:provenance}
Part~II is built from the manuscript \emph{Support-Bounded Surreal Fields: Sharp saturation and spherical thresholds, an exponential obstruction, and indistinguishable omnific quotients}, dated 23 September 2026 (25 pages, author line ``Research manuscript prepared for Vladimir Reshetnikov''). It was delivered as the archive \texttt{support\_bounded\_surreal\_fields} (manuscript~04 of batch~32, delivered in commit \texttt{aa268a4}) and placed in commit \texttt{7d04483}. Its repository comparison is pinned at commit \texttt{343dc2c471212bb9b53ff4623bace2e1943f255b}. In this report it is \emph{source~02}; the manuscript of Part~I is source~01. Its program, build script, recorded run, build record and source audit are shipped with the prefix \texttt{02-support-bounded-fields-}, unchanged; the manuscript, its README and its PDF are not shipped, and its text is printed here.

Source~02 does not cite Part~I, although Part~I was in the collection at its pin. It is the \emph{class} version of Part~I: the exponent group is the whole class $\No$ rather than a set $\Gamma$, so the field contains every Conway monomial. It contributes the classification and the common character of the set-presented gaps, the exact saturation threshold, the strong-sum threshold, the class form of the empty-nest threshold, the topology relative to $\No$, closure under the surreal exponential together with the nonexistence of any surjective ordered exponential, the arithmetic of the omnific integer part and its set-sized quotients, and a proper class of pairwise nonconjugate real forms of $\SC$. It also asks eight questions (\cref{fkc:sb:sec:questions}).

The merge made the following choices.
\begin{itemize}
\item \emph{Placement.} Part~II comes after the conclusion of Part~I and before the appendices, so that no number of Part~I changes. Source~02's appendices are \cref{fkc:sb:app:audit,fkc:sb:app:realization}, after \cref{fkc:app:checks}.
\item \emph{Notation.} Source~02's $F_\kappa$ collides with Part~I's full Hahn field $F_k$, and several further symbols collide with Part~I or with neighbouring reports. They are renamed as listed in \cref{fkc:sb:sub:notation}.
\item \emph{Printed once.} Where source~02 proves a statement that Part~I proves in substance, the argument is printed once, with a reference to Part~I and a credit to source~02 as an independent derivation: the small-exponent-field lemma (\cref{fkc:sb:lem:hahn}), the ambient gluing lemma (\cref{fkc:sb:lem:glue}) and the lower bound and sharpness of the empty-nest threshold (\cref{fkc:sb:thm:spherical}). All other results, proofs, examples, questions and non-claims of source~02 are printed in full.
\item \emph{Stale statements.} Source~02's novelty claim for the set-sized quotients of the omnific integer part is stale against the collection as it was at its pin; it is corrected in \cref{fkc:sb:sub:corrections}. Statements of Part~I that Part~II makes incomplete carry notes marked ``Status (batch~32)''.
\end{itemize}

\subsection{The main results of source 02}
\label{fkc:sb:sub:main}

\paragraph{Status and scope (source 02).}
Source~02 is an AI-assisted, unrefereed research manuscript with proposed contributions. Its proofs are mathematical arguments, not certified Lean proofs. The specific hierarchy and its combined consequences are the candidate contributions; bounded Hahn fields, exponential obstructions, normal-form methods, constant-term quotient arguments, and nonconjugate surcomplex real forms have substantial antecedents. No named published conjecture is represented as settled, and bibliographic priority has not been independently certified. \Cref{fkc:sb:sec:provenance} records the precise comparison with the supplied repository and the literature.

\paragraph{The questions.}
Surreal numbers package three structures in the same object: an ordered field, a hierarchy of monomials, and a set-supported normal form. An individual number may involve enormously complicated exponents while having just one term. This makes it important to distinguish a bound on \emph{support cardinality} from a bound on birthday, on support order type, or on the exponents themselves.

We ask what happens when one keeps \emph{every} monomial but imposes a uniform cardinal bound on the number of nonzero coefficients. The answer is not simply a smaller version of the full surreal field. Algebraic closedness properties survive; certain completeness properties fail at an exactly computable cardinal; exponentiation survives only in one direction; and all set-sized arithmetic quotients become unable to see the support bound.

Fix an uncountable cardinal $\kappa$ and put $\mu=\cf(\kappa)$. Throughout Part~II,
\begin{equation}\label{fkc:sb:eq:family}
 \KN\kappa\R=\{x\in\No:\card{\suppw(x)}<\kappa\},\qquad
 \KN\kappa\C=\KN\kappa\R[i],\qquad \Ok=\KN\kappa\R\cap\Oz.
\end{equation}
The meaning of the proper classes and of set-sized tests is specified in \cref{fkc:sb:sec:setup}.

\begin{theorem}[The support hierarchy]\label{fkc:sb:thm:main}
For every uncountable cardinal $\kappa$, the following assertions hold.
\begin{enumerate}[label=\textup{(\roman*)}]
\item $\KN\kappa\R$ is a truncation-closed real closed proper-class subfield of $\No$, and $\KN\kappa\C$ is algebraically closed. They have the full surreal value group and residue fields $\R$ and $\C$, respectively.
\item The set-presented omitted cuts of $\KN\kappa\R$ are in bijection with normal forms of order type exactly $\kappa$. Each has cofinality pair $(\mu,\mu)$. The field is $\mu$-saturated and is not $\mu^+$-saturated in the ordered-field language.
\item Every ambient strongly summable family of fewer than $\mu$ members of $\KN\kappa\R$ has its sum in $\KN\kappa\R$, and this fails for a family of size $\mu$. Every nest of fewer than $\mu$ closed valuation balls has a common point, and an explicit nest of size $\mu$ has no common point.
\item The inherited surreal exponential maps $\KN\kappa\R$ into itself. There is, however, no ordered-group isomorphism $(\KN\kappa\R,+)\cong(\KN\kappa\R^{>0},\cdot)$.
\item $\Ok$ is an integer part of $\KN\kappa\R$, with fraction field $\KN\kappa\R$. Every unital homomorphism from $\Ok$ to a set-sized ring factors uniquely through $\ct:\Ok\to\Z$.
\end{enumerate}
\end{theorem}

The field and quotient assertions are proved separately rather than inferred from properties of the whole surreal field: a homomorphism on a proper subring need not extend to the larger ring. The singular-cardinal case is also treated directly, rather than hidden behind a regularity assumption. Part~(i) is \cref{fkc:sb:thm:fields,fkc:sb:prop:hierarchy}; part~(ii) is \cref{fkc:sb:thm:prefix,fkc:sb:thm:gapcharacter,fkc:sb:cor:saturation}; part~(iii) is \cref{fkc:sb:thm:sums,fkc:sb:thm:spherical}; part~(iv) is \cref{fkc:sb:thm:expclosure,fkc:sb:thm:noexp}; part~(v) is \cref{fkc:sb:prop:integerpart,fkc:sb:thm:smallquotient}.

\begin{theorem}[Intrinsic real forms]\label{fkc:sb:thm:mainrealforms}
Assume Global Choice. For each uncountable regular cardinal $\kappa$ there is an involution $\tau_\kappa$ of the same pure field $\SC$ such that $\Fix(\tau_\kappa)$ is isomorphic to $\KN\kappa\R$. These involutions can all agree with ordinary complex conjugation on $\C$, but no two are conjugate by a field automorphism of $\SC$. Their fixed fields, and their corresponding rings $\Ok$, are pairwise nonisomorphic. Nevertheless, all the rings $\Ok$ have the same set-sized quotient theory described in \cref{fkc:sb:thm:main}\textup{(v)}.
\end{theorem}

\paragraph{What is and is not classified.}
We classify all \emph{set-presented gaps of the specified fields}, not all surreal gaps, not all real forms, and not all bounded Hahn fields. Different cofinalities distinguish these fields. Whether equal cofinalities suffice for an abstract field isomorphism remains a question in this part (\cref{fkc:sb:q:equalcf}). The class-wide isomorphisms in \cref{fkc:sb:thm:mainrealforms} are not asserted to preserve monomials, valuation, exponentiation, or the distinguished omnific ring.

\subsection{Relation to Part I and to the collection \mergetag}
\label{fkc:sb:sub:relation}

Part~I fixes a \emph{set} $\Gamma$ and studies the set-sized fields $\K$; its surreal transfer (\cref{fkc:sec:surreal}) realizes them inside $\No$ through an embedding $j:\Gamma\to\No$, so their images contain only the monomials $\omega^{-j(\gamma)}$. Part~II takes the exponents to be all of $\No$. For each divisible subgroup $\Gamma\subseteq\No$ that is a set, with $j$ the inclusion, the image $\iota_j(K_{\kappa,\R})$ of \eqref{fkc:eq:conway} is exactly $\KN\kappa\R\cap\R((\omega^{\Gamma}))$ (the exponents $-j(\gamma)$ range over $-\Gamma=\Gamma$), and $\KN\kappa\R$ is the union of these images; but no single set-sized workspace carries $\KN\kappa\R$, and the statements of Part~II concern the proper class. The table records, for each result of Part~II, the result of Part~I it extends or parallels and the antecedents in other reports; all are credited again where the result is printed.

\begin{center}
\small
\begin{longtable}{@{}>{\raggedright\arraybackslash}p{0.27\textwidth}>{\raggedright\arraybackslash}p{0.30\textwidth}>{\raggedright\arraybackslash}p{0.35\textwidth}@{}}
\toprule
Part II & Part I (set-sized $\Gamma$) & Elsewhere in the collection \\
\midrule\endhead
\Cref{fkc:sb:lem:card} (cardinal estimates) & the counts in \cref{fkc:lem:workspace,fkc:lem:glue} & \\
\Cref{fkc:sb:lem:hahn} (small exponent field) & \cref{fkc:lem:workspace,fkc:prop:closed} with $\Gamma=H$; printed once & the proof of \lbl{lce:sg:thm:shortfield} \\
\Cref{fkc:sb:thm:fields,fkc:sb:prop:hierarchy} (fields, hierarchy) & \cref{fkc:prop:closed}, \cref{fkc:cor:tower} & \lbl{lce:sg:thm:shortfield} (real closed, monomials, truncation, every uncountable $\kappa$) \\
\Cref{fkc:sb:thm:prefix,fkc:sb:thm:gapcharacter} (gaps, character $(\mu,\mu)$) & \cref{fkc:thm:real-type,fkc:lem:fibers,fkc:cor:count}; \cref{fkc:thm:approx,fkc:thm:omission} & \lbl{univ:thm:firstgap}, \lbl{univ:cor:recover} (other fields) \\
\Cref{fkc:sb:thm:supporttypes,fkc:sb:cor:saturation} (saturation) & fills non-claim N5 for the class field & \lbl{univ:cor:spectrum} (other fields) \\
\Cref{fkc:sb:thm:sums} (strong sums) & none & \lbl{opa:as:prop:kappa} (regular $\kappa$, set $\Gamma$); \lbl{lce:sg:thm:threshold} \\
\Cref{fkc:sb:lem:glue,fkc:sb:thm:spherical} (nests) & \cref{fkc:lem:glue,fkc:thm:balls}; printed once & \lbl{duals:prop:spherical}; \lbl{prop:spherical} of the rank-one report \\
\Cref{fkc:sb:prop:cauchy,fkc:sb:prop:closed} (topology) & fills N6; compare \cref{fkc:thm:dichotomy} & \lbl{found:thm:discrete}, \lbl{univ:prop:topology}; \lbl{lce:sg:prop:closed-support} \\
\Cref{fkc:sb:thm:expclosure,fkc:sb:prop:expimage,fkc:sb:thm:noexp,fkc:sb:thm:loghull} (exponential) & fills N7 for the class field & \lbl{lce:cor:expclosure}, \lbl{lce:thm:twoclosures} \\
\Cref{fkc:sb:prop:integerpart} (integer part) & none & \lbl{lce:sg:prop:integerpart}, \lbl{lce:eq:Ozfrac} \\
\Cref{fkc:sb:lem:reservoir,fkc:sb:thm:smallquotient,fkc:sb:cor:modules}, \cref{fkc:sb:sub:gaussian} (quotients) & none & \lbl{osq:thm:support}~(i), \lbl{osq:lem:collision} \\
\Cref{fkc:sb:cor:rings} (rings) & none & \\
\Cref{fkc:sb:lem:backforth,fkc:sb:prop:classiso} (back-and-forth) & compare \cref{fkc:thm:pure} & \lbl{saut:thm:acfhomogeneity} \\
\Cref{fkc:sb:thm:involutions,fkc:sb:prop:sametheory} (real forms) & the named-conjugation language of \cref{fkc:sec:surreal} & \lbl{univ:q:families}, \lbl{univ:thm:manygrounds}, \lbl{saut:prop:conjugacycriterion} \\
\bottomrule
\end{longtable}
\end{center}

The reports cited are \nolinkurl{large-cardinal-embeddings-and-normal-forms} (prefix \lbl{lce:}) \cite{CollLCE}, \nolinkurl{surreal-fields-across-universes} (\lbl{univ:}) \cite{CollUniv}, \nolinkurl{omnific-preserving-automorphisms} (\lbl{opa:}) \cite{CollOPA}, \nolinkurl{set-sized-quotients-of-omnific-integers} (\lbl{osq:}) \cite{CollOSQ}, \nolinkurl{surcomplex-field-automorphisms} (\lbl{saut:}) \cite{CollSAUT}, \nolinkurl{foundations} (\lbl{found:}) \cite{CollFound}, \nolinkurl{three-duals-of-hahn-vector-spaces} (\lbl{duals:}) \cite{CollDuals} and \nolinkurl{rank-one-berkovich} (bare labels) \cite{CollBerk}. None of them is used as a lemma in Part~II, except where a proof below says so explicitly. In particular, the field $\KN\kappa\R$ is the large-cardinal report's field of short normal forms $\mathscr S_{<\kappa}$, and $\Ok$ is its ring $\mathscr I_{<\kappa}$; that report proves real closedness for every uncountable $\kappa$ and the facts cited in the table, but not the gap, saturation, strong-sum, nest or quotient results of Part~II. Its exponential results are stated at the critical point of an elementary embedding.

\subsection{Notation of this part}
\label{fkc:sb:sub:notation}

Part~II writes surreal normal forms with \emph{decreasing} exponents, as source~02 does, and keeps source~02's symbols except those in the table. Every renaming is a change of name only; no normalization changed.

\begin{center}
\small
\begin{longtable}{@{}>{\raggedright\arraybackslash}p{0.2\textwidth}>{\raggedright\arraybackslash}p{0.22\textwidth}>{\raggedright\arraybackslash}p{0.5\textwidth}@{}}
\toprule
Source 02 & Here & Reason \\
\midrule\endhead
$F_\kappa$, $F_\lambda$ & $\KN\kappa\R$, $\KN\lambda\R$ & $F_k$ is Part~I's \emph{full} Hahn field; Part~I's $\K$ is the bounded field, so the class version is written $K_{\kappa,k}(\No)$. It is the large-cardinal report's $\mathscr S_{<\kappa}$. \\
$C_\kappa$ & $\KN\kappa\C$ & The class version of $\KC$; not the large-cardinal report's common field $C_j$. \\
$O_\kappa$ & $\Ok$ & The large-cardinal report's $\mathscr I_{<\kappa}$; the quotient report's $\mathcal A^{<\kappa}_{\Z,\R}$. \\
$J_\kappa$ & $\Pk$ & The large-cardinal report's $J$ is an elementary-embedding lift; $\Pk$ is the quotient report's $\{f\in\Pi_\R:\card{\supp f}<\kappa\}$. \\
$G_\kappa$ & $\Ok[i]$ & No new letter. \\
$\supp$ & $\suppw$ & Part~I's $\supp f\subseteq\Gamma$ is increasing; $\suppw(\iota_j f)=-j(\supp f)$. \\
$s_\kappa$, $s_\beta$ & $h_\kappa$, $h_\beta$ & Part~I's $s_\alpha$ are the support exponents of a missing $x$. $h_\kappa$ is the large-cardinal report's $h_\kappa$. \\
$p_\kappa$ & $z_\kappa$ & Part~I's $p$ is a type; $z_\kappa$ is the large-cardinal report's $z_\kappa$ (its $p_\kappa$ is a different series). \\
prefix $p\in\mathcal P_\kappa$; $t_\beta$ & $T\in\mathcal P_\kappa$; $T\restr\beta$ & Part~I's $p$ is a type and $t$ its Hahn variable. $T_\kappa(y)=y\restr\kappa$. \\
$L=\R((\omega^H))$ & $\R((\omega^H))$ & Part~I's $L_{\kappa,k}$ is the completion. \\
$(L,R)$, $L_y$, $R_y$ & $(\mathcal L,\mathcal R)$, $\mathcal L_y$, $\mathcal R_y$ & Gap sides; $L_{\kappa,k}$ and $\R$ are taken. \\
$E$, $E_2$ (exponential) & $\mathcal E$, $\mathcal E_2$ & Part~I's $E_x=K+Kx$. \\
$E$ (reservoir field) & $W$ & The same. \\
$h=g^{-1}$ & $g^{-1}$ & $h_\kappa$ is now the series above. \\
$h$ (tail), $h_\alpha$, $h\in H$ (reservoir) & $w$, $d_\alpha$, $u\in H$ & The same. \\
$\tau$ (reservoir cardinal) & $\rho$ & $\tau_\kappa$ are the involutions. \\
$\phi$ (ring map), $\bar\phi$ & $\chi$, $\bar\chi$ & Part~I's $\varphi$ is a formula. \\
$\phi_\kappa$ & $\psi_\kappa$ & Part~I's $\varphi$ and $\Phi$ are formulas and sets of formulas. \\
$c_\kappa$, $c$ & $\conj_\kappa$, $\conj$ & In Part~I, $c_\kappa$ would be the $\kappa$-th coefficient of \eqref{fkc:eq:enumerate}. \\
$j_\kappa$ & $\tau_\kappa$ & Part~I's $j:\Gamma\to\No$ is an exponent embedding; the large-cardinal report's $j$ is an elementary embedding. \\
$S$ (Lemma 4.1) & $\Sigma$ & $S$ is kept for the set-sized target ring, as in the quotient report. \\
$D$ (proof of Thm.~4.4) & $D'$ & Part~I's $D_x$ is an approximation cut. \\
\bottomrule
\end{longtable}
\end{center}

The following symbols agree with Part~I and are not renamed: $\kappa$, $\lambda$, $\mu=\cf(\kappa)$; the valuation, since $v(x)=-\lead(x)$ on $\No$ is Part~I's $v=\min\supp$ transported by $t^\gamma\mapsto\omega^{-\gamma}$; the closed ball $B(a,\gamma)=\{x:v(x-a)\ge\gamma\}$; and the small divisible group $H$ of \cref{fkc:lem:workspace}. The letter $g$ is reserved in Part~II for Gonshor's increasing bijection $g:\No^{>0}\to\No$ of \eqref{fkc:sb:eq:gonshor-reindex}, as in the large-cardinal report; exponent variables are written $e$. Local letters of Part~I's proofs ($a_i$, $b_\alpha$, $u_\alpha$, $\delta$, $\ell$, $u$) do not occur in Part~II with those meanings; in Part~II, $\ell_\beta,u_\beta$ are brackets (\eqref{fkc:sb:eq:brackets}), $\delta=\omega^d$ is a monomial (\cref{fkc:sb:lem:reservoir}), and $a_\alpha$ are normal-form exponents.

Five tempting readings are false.
\begin{itemize}
\item $\KN\kappa\R$ is \emph{not} $\iota_j(K_{\kappa,\R})$ for any set-sized $\Gamma$: it is a proper class containing every monomial. Part~I's theorems apply to each set-sized piece $\KN\kappa\R\cap\R((\omega^H))$, not to $\KN\kappa\R$.
\item The gap character $(\mu,\mu)$ of \cref{fkc:sb:thm:gapcharacter} and the omission number $\mu$ of \cref{fkc:thm:omission} are equal cardinals attached to different objects: a cut of a proper-class field, and a type over a set-sized field realized in a fixed full Hahn field.
\item The saturation of \cref{fkc:sb:cor:saturation} is saturation of the class field over set-sized parameter sets. It does not contradict non-claim N5, which concerns the set fields $\K$.
\item ``Complete'' in \cref{fkc:thm:dichotomy} refers to the intrinsic topology of a set workspace. In Part~II every set-indexed Cauchy net is eventually constant (\cref{fkc:sb:prop:cauchy}), for every $\kappa$, because the value group $\No$ has no set-sized cofinal subset.
\item A strong sum is not a limit in any topology used here (\cref{fkc:sb:sec:spherical}).
\end{itemize}

\section{Foundations, normal forms, and cardinal bookkeeping}\label{fkc:sb:sec:setup}

\subsection{Set and class conventions}
We work in G\"odel--Bernays class theory with Global Choice. All supports, polynomials, parameter sets, index sets of sums, nests of balls, and cofinal presentations of cuts are sets. A class function is required to have a set image on each set domain, as supplied by class Replacement. Assertions about satisfaction use each fixed first-order formula; no unrestricted truth predicate for a class structure is used.

Most proofs require only ordinary set constructions with the surreal class as an available parameter. Global Choice is used explicitly for the class back-and-forth in \cref{fkc:sb:sec:realforms}. No class version of Zorn's lemma is used. A ``proper-class-indexed family'' there is a single class relation coding the indicated maps, not a set of all those maps.

For an infinite cardinal $\nu$, a class structure is $\nu$-saturated if every complete type in finitely many variables over a parameter set of cardinality less than $\nu$, consistent with its elementary diagram on those parameters, is realized. Types themselves are sets. The usual reduction to one-variable types applies by successive realization, since only finitely many parameters are added at a time.

\subsection{Normal forms and the natural valuation}
We use the normal-form theorem of Conway and Gonshor in the form reviewed in \cite[Section~2]{MM}:
\begin{equation}\label{fkc:sb:eq:nf}
 x=\sum_{\alpha<\lambda}r_\alpha\omega^{a_\alpha},\qquad
 r_\alpha\in\R\setminus\{0\},\qquad
 a_\alpha>a_\beta\quad(\alpha<\beta<\lambda).
\end{equation}
Here $\lambda$ is a set ordinal. Conversely, every such reverse well-ordered, set-supported expression defines a surreal number. The support $\suppw(x)$ is the set of its \emph{growth exponents} $a_\alpha$. In particular, $\suppw(\omega^a)=\{a\}$, regardless of the complexity of $a$.

For nonzero $x$ define
\[
 \lead(x)=a_0,\qquad v(x)=-a_0,
\]
and put $v(0)=\infty$. Thus $v(\omega^{-a})=a$, and
\[
 v(xy)=v(x)+v(y),\qquad v(x+y)\ge\min\{v(x),v(y)\}.
\]
For $z=x+iy\in\SC$, combine the two normal forms coefficientwise over $\C$ and define $v(z)=\min\{v(x),v(y)\}$. Its support is the union of the two real supports. This is the Hahn extension of the valuation to $\SC$; no order on the complex field is intended.

Comparison in $\No$ is lexicographic at the greatest exponent where coefficients differ. A truncation is obtained by retaining an initial segment of the normal form. We write $x\restr\beta$ for the first $\beta$ terms, when $\beta\le\lambda$. \mergetag{} Under \eqref{fkc:eq:conway}, $x\restr\kappa$ is Part~I's first-$\kappa$ prefix $T_\kappa$ of \eqref{fkc:eq:prefix-intro}, and we also write $T_\kappa(x)=x\restr\kappa$.

Every set of surreal numbers has upper and lower bounds in $\No$. Every set of positive surreal numbers has a strictly positive lower bound. More generally, if $A,B$ are sets and $A<B$, there is $y\in\No$ with $A<y<B$, including when one side is empty. These are immediate consequences of the surreal cut construction. All arguments below using ``a fresh smaller scale'' refer to this set-cut property.

Every surreal decomposes uniquely as
\begin{equation}\label{fkc:sb:eq:split}
 x=x^\uparrow+r+\epsilon,
\end{equation}
where the three supports lie, respectively, in $\No^{>0}$, $\{0\}$, and $\No^{<0}$. Thus $r\in\R$, $\epsilon$ is infinitesimal, and $x^\uparrow$ is purely infinite; ``purely infinite'' allows negative numbers and includes zero. The omnific integers have the normal-form description
\begin{equation}\label{fkc:sb:eq:oz}
 \Oz=\Z+\{x\in\No:\suppw(x)\subseteq\No^{>0}\}.
\end{equation}

\subsection{The support estimates that matter}
\begin{lemma}[Two different cardinal estimates]\label{fkc:sb:lem:card}
Let $\kappa>\aleph_0$ and $\mu=\cf\kappa$.
\begin{enumerate}[label=\textup{(\alph*)}]
\item If $S$ is a set of cardinality less than $\kappa$, its $\Q$-linear span in any rational vector space has cardinality at most $\max(\card S,\aleph_0)<\kappa$.
\item A union of fewer than $\mu$ sets, each of cardinality less than $\kappa$, has cardinality less than $\kappa$.
\item The additive monoid of finite sums from a fixed set $S$ has cardinality at most $\max(\card S,\aleph_0)$.
\end{enumerate}
\end{lemma}
\begin{proof}
Finite rational linear combinations are indexed by finite sequences from $S\times\Q$, proving (a). For (b), fewer than $\cf\kappa$ cardinals below $\kappa$ are bounded by one cardinal below $\kappa$; the product with the number of sets is still below $\kappa$. Part (c) is the same finite-sequence count without rational coefficients.
\end{proof}

Part (c), not a countable-union regularity argument, is the reason that Taylor exponentials remain inside $\KN\kappa\R$ even when $\cf\kappa=\aleph_0$. By contrast, arbitrary strong sums and arbitrary nests are governed by part (b). \mergetag{} Part~(a) is the count in the proof of \cref{fkc:lem:workspace}, part~(b) the count at the end of the proof of \cref{fkc:lem:glue}, and part~(c) the monoid count behind the countable boundary of \cref{fkc:app:checks}.

\section{Real closed fields with unrestricted monomials}\label{fkc:sb:sec:fields}

\begin{lemma}[A small exponent field containing any finite tuple]\label{fkc:sb:lem:hahn}
Every finite tuple from $\KN\kappa\R$ lies in a set-sized real closed Hahn subfield $\R((\omega^H))\subseteq\KN\kappa\R$. One can take
\[
 H=\Span\bigcup_j\suppw(x_j).
\]
More generally, this conclusion holds for any set of elements whose union of supports has cardinality less than $\kappa$.
\end{lemma}
\begin{proof}[Proof \textup{(printed once)}]
By \cref{fkc:sb:lem:card}, $H$ is a divisible ordered set group of cardinality less than $\kappa$. Every series supported in $H$ is a surreal normal form with support of cardinality less than $\kappa$, and normal-form addition and multiplication agree with Hahn addition and multiplication, so the natural realization is a field embedding of the set Hahn field into $\No$: it is the map $\iota_j$ of \eqref{fkc:eq:conway} for $\Gamma=H$ and $j$ the inclusion, whose image is supported in $-H=H$. The Hahn field $\R((\omega^H))$ is real closed because $H$ is divisible \cite[Section~2]{KS}, \cite[Section~2]{MM}. This is \cref{fkc:lem:workspace} together with the real case of \cref{fkc:prop:closed}, applied with $\Gamma=H$, where $K_{\kappa,\R}=\R((t^H))$ because $\card H<\kappa$; source~02 gives the same argument independently. All coefficients and all possible supports here range over sets, so this is an ordinary set-field theorem, not a new class-field closure principle.
\end{proof}

\begin{theorem}[Algebraic structure]\label{fkc:sb:thm:fields}
The class $\KN\kappa\R$ is a real closed subfield of $\No$, contains $\R$ and every monomial $\omega^a$, and is closed under truncation. Its complexification $\KN\kappa\C$ is algebraically closed. Both inclusions
\[
 \KN\kappa\R\subseteq\No,\qquad \KN\kappa\C\subseteq\SC
\]
are elementary in the respective pure ordered-field and pure-field languages.
\end{theorem}
\begin{proof}
Apply \cref{fkc:sb:lem:hahn} to the finitely many elements needed for each field operation, positive square root, or odd-degree polynomial. The usual characterization of real closed fields then proves the assertion. Alternatively, $\KN\kappa\R$ is a directed union of the real closed fields from that lemma. Truncation can only decrease support, and monomials have one term. Algebraic closedness of $\KN\kappa\R[i]$ follows from real closedness. The two elementarity assertions are the standard quantifier elimination consequences for real closed and algebraically closed fields.
\end{proof}
\mergetag{} This is the class analogue of \cref{fkc:prop:closed} and of \cref{fkc:thm:main}(i). The large-cardinal report proves the same real closedness, with the monomials and truncation, for every uncountable $\kappa$ by the same directed-union argument (\lbl{lce:sg:thm:shortfield}); source~02 does not cite it.

\begin{proposition}[Strict hierarchy and unchanged local invariants]\label{fkc:sb:prop:hierarchy}
If $\kappa<\lambda$ are uncountable cardinals, then $\KN\kappa\R\subsetneq\KN\lambda\R$. The union of the fields $\KN\kappa\R$ is $\No$. Each $\KN\kappa\R$ is a proper class, with natural value group $(\No,+)$ and residue field $\R$; $\KN\kappa\C$ has residue field $\C$.
\end{proposition}
\begin{proof}
The explicit series
\begin{equation}\label{fkc:sb:eq:hkappa}
 h_\kappa=\sum_{\alpha<\kappa}\omega^{-\alpha}
\end{equation}
has exactly $\kappa$ nonzero terms, so it belongs to $\KN\lambda\R$ but not to $\KN\kappa\R$. Every surreal has a set support, proving the union claim. Every $\omega^a$ belongs to every $\KN\kappa\R$; hence these fields are proper classes and realize every value. The finite elements modulo infinitesimals give exactly their real constant terms. The complex case uses the union of the real and imaginary supports and the same argument with complex coefficients.
\end{proof}

\begin{remark}[Why the countable boundary is excluded]\label{fkc:sb:rem:countable}
The definition with $\kappa=\aleph_0$ gives finite-support series, not a field. For example,
\[
 (1-\omega^{-1})^{-1}=\sum_{n\ge0}\omega^{-n}
\]
has infinite support. No cardinal arithmetic such as GCH or $\kappa^{<\kappa}=\kappa$ is assumed anywhere in \cref{fkc:sb:thm:fields}. Singular uncountable cardinals cause no problem because every individual finite algebraic calculation lives in one small exponent field. \mergetag{} This is the class form of \cref{fkc:rem:singular} and of the first paragraph of \cref{fkc:app:checks}.
\end{remark}

\section{Exact classification of set-presented gaps}\label{fkc:sb:sec:gaps}

\subsection{An economical separator}
\begin{lemma}[First exponent outside a prescribed support set]\label{fkc:sb:lem:outside}
Let $\Sigma\subseteq\No$ be a set and let $y\in\No$ with $\suppw(y)\nsubseteq\Sigma$. Let $\gamma$ be the greatest exponent of $y$ not in $\Sigma$, and let $z$ be the truncation of $y$ through its $\gamma$-term. Then
\[
 \suppw(z)\subseteq\Sigma\cup\{\gamma\}.
\]
For every $a\in\No$ supported in $\Sigma$,
\[
 a<y\ \Longleftrightarrow\ a<z,
 \qquad
 a>y\ \Longleftrightarrow\ a>z.
\]
\end{lemma}
\begin{proof}
The nonempty subset $\suppw(y)\setminus\Sigma$ of a reverse well-ordered set has a greatest member. The coefficients of $y$ and $z$ agree at every exponent at least $\gamma$, whereas $a$ has coefficient zero at $\gamma$ and $y$ has a nonzero coefficient there. Thus the first coefficient difference between $a$ and $y$ occurs at or above $\gamma$, where truncation has changed nothing. All exponents of $z$ above $\gamma$ belong to $\Sigma$ by the choice of $\gamma$.
\end{proof}

The point of the lemma is that $\Sigma$ need not itself be reverse well-ordered. It may be a large union of supports or a divisible group. Only the portion actually used by $y$ is retained.

\subsection{Gap conventions and the prefix classification}
\begin{definition}\label{fkc:sb:def:gap}
A gap of an ordered class $F$ is a partition $F=\mathcal L\sqcup\mathcal R$ into two nonempty classes with $\mathcal L<\mathcal R$, with no greatest element of $\mathcal L$ and no least element of $\mathcal R$. It is \emph{set-presented} if $\mathcal L$ has a set-sized cofinal subset and $\mathcal R$ a set-sized coinitial subset. Its cofinality pair is $(\cf\mathcal L,\operatorname{ci}\mathcal R)$, computed using such set subsets.
\end{definition}

For $y\in\No\setminus\KN\kappa\R$, let
\[
 \mathcal L_y=\{x\in\KN\kappa\R:x<y\},\qquad
 \mathcal R_y=\{x\in\KN\kappa\R:x>y\}.
\]
These do form a gap. Indeed, if $x<y$, choose a positive monomial smaller than $y-x$ and add it to $x$; the result remains in $\KN\kappa\R$ and lies strictly between them. The upper side is similar. Both sides are nonempty because every surreal is bounded in magnitude by a monomial.

\begin{theorem}[First-$\kappa$-term classification]\label{fkc:sb:thm:prefix}
Let $\mathcal P_\kappa$ be the class of normal forms whose support has reverse order type exactly the initial ordinal $\kappa$.
\begin{enumerate}[label=\textup{(\alph*)}]
\item If $y\notin\KN\kappa\R$, its first $\kappa$ terms $T=y\restr\kappa$ determine its cut on $\KN\kappa\R$.
\item Two elements of $\No\setminus\KN\kappa\R$ determine the same cut on $\KN\kappa\R$ if and only if their first $\kappa$ terms agree.
\item The map $T\mapsto(\mathcal L_T,\mathcal R_T)$ is a bijection from $\mathcal P_\kappa$ to all set-presented gaps of $\KN\kappa\R$.
\item If $T=\sum_{\alpha<\kappa}r_\alpha\omega^{a_\alpha}$, the class of its fillers in $\No$ is exactly
\begin{equation}\label{fkc:sb:eq:fillercone}
 T+\{\epsilon\in\No:v(\epsilon)>-a_\alpha
                         \text{ for every }\alpha<\kappa\}.
\end{equation}
\end{enumerate}
\end{theorem}
\begin{proof}
Fix $x\in\KN\kappa\R$. If $x$ agreed with $y$ at every exponent at or above each $a_\alpha$ for $\alpha<\kappa$, its support would contain all these $\kappa$ exponents. That is impossible. Thus some first difference between $x$ and $y$ occurs before the end of the first $\kappa$ terms. Discarding later terms does not change its sign. This proves (a).

If two first-$\kappa$ prefixes are distinct, consider their greatest coefficient difference. Their common prefix above that exponent has fewer than $\kappa$ terms: otherwise both first-$\kappa$ prefixes would already have been exhausted in a common initial segment. Retain this common prefix and use, at the first difference, a real coefficient strictly between the two coefficients, regarding a missing coefficient as zero. This constructs an element of $\KN\kappa\R$ strictly between the two numbers. It separates their cuts. The converse is (a), proving (b).

To see directly that each $T\in\mathcal P_\kappa$ gives a set-presented gap, put
\begin{equation}\label{fkc:sb:eq:brackets}
 \ell_\beta=T\restr\beta-(|r_\beta|+1)\omega^{a_\beta},\qquad
 u_\beta=T\restr\beta+(|r_\beta|+1)\omega^{a_\beta}
       \qquad(\beta<\kappa).
\end{equation}
They belong to $\KN\kappa\R$ and satisfy $\ell_\beta<T<u_\beta$. For every $x<T$ in $\KN\kappa\R$, a sufficiently late first-difference index is passed by $\beta$, at which point $x<\ell_\beta$. Similarly, the $u_\beta$ are coinitial above $T$.

Conversely, let $(\mathcal L,\mathcal R)$ be a set-presented gap, with sets $A\subseteq\mathcal L$ and $B\subseteq\mathcal R$ respectively cofinal and coinitial. The set-cut property of $\No$ gives $A<y<B$. Such $y$ cannot belong to $\KN\kappa\R$: if it belonged to $\mathcal L$, cofinality of $A$ would give $y\le a$ for some $a\in A$, and the case $y\in\mathcal R$ is dual. Cofinality also shows that $y$ induces precisely $(\mathcal L,\mathcal R)$. Part (a) supplies its prefix, completing (c).

Finally, equality of the first $\kappa$ terms means equality of every coefficient down through each $a_\alpha$. Equivalently, the difference has valuation strictly greater than every $-a_\alpha$. This proves (d). The right-hand class is nonempty and has nonzero members: the set of values $\{-a_\alpha:\alpha<\kappa\}$ has a strict upper bound in $\No$.
\end{proof}
\mergetag{} Part~(b) and the filler class \eqref{fkc:sb:eq:fillercone} are the class counterparts of \cref{fkc:thm:real-type} and \cref{fkc:lem:fibers}: under $t^\gamma\mapsto\omega^{-\gamma}$, the condition $v(\epsilon)>-a_\alpha$ for all $\alpha<\kappa$ is membership in Part~I's tail group $I_{D_x}$. Part~I classifies types over a set field realized in one full Hahn field; here the objects are the cuts of a proper-class field, and the set-cut property of $\No$ supplies every filler.

\subsection{The single gap character}
\begin{theorem}[Exact cofinalities, including singular cardinals]\label{fkc:sb:thm:gapcharacter}
Every set-presented gap of $\KN\kappa\R$ has cofinality pair $(\mu,\mu)$, where $\mu=\cf\kappa$. Such gaps exist. Consequently, if $A<B$ are sets of elements of $\KN\kappa\R$ and $\min(\card A,\card B)<\mu$, there is $x\in\KN\kappa\R$ with $A<x<B$.
\end{theorem}
\begin{proof}
Use \cref{fkc:sb:thm:prefix} and its brackets. Choose an increasing cofinal sequence $(\beta_\xi)_{\xi<\mu}$ in $\kappa$. The brackets at these indices are still cofinal and coinitial, so both characters are at most $\mu$.

Let $D'\subseteq\mathcal L_T$ have cardinality less than $\mu$. For each $d\in D'$, choose an index $\alpha_d<\kappa$ by which its first difference from $T$ has occurred. The set of these indices is bounded in $\kappa$. A bracket $\ell_\beta$ at a later index exceeds every $d\in D'$ and still lies below $T$. Thus $D'$ is not cofinal. The proof for the upper side is the same. Existence follows from $T=h_\kappa$ in \eqref{fkc:sb:eq:hkappa}.

For the last assertion, first fill $A<B$ in $\No$. If the filler lies in $\KN\kappa\R$, we are done. Otherwise it induces one of the gaps above. If no element of $\KN\kappa\R$ lies strictly between $A$ and $B$, then $A$ must be cofinal on the lower side of that gap: a lower-side element strictly greater than all of $A$ would itself be a filler. Similarly, $B$ must be coinitial. This contradicts the smaller-side hypothesis. If a side is empty, a monomial bound for the other set gives a filler directly.
\end{proof}
\mergetag{} The upper bound is the class form of $\cf(D_x)=\mu$ in \cref{fkc:thm:approx}, and the lower bound parallels the lower bound of \cref{fkc:thm:omission}. The across-universes report proves the analogous statement for the old surreal field of an inner model, with its threshold $\delta(M,N)$ in place of $\mu$ (\lbl{univ:thm:firstgap}, \lbl{univ:cor:recover}).

\begin{corollary}[An intrinsic isomorphism obstruction]\label{fkc:sb:cor:noniso}
If $\cf\kappa\ne\cf\lambda$, then $\KN\kappa\R$ and $\KN\lambda\R$ are not isomorphic as ordered fields, as ordered sets, or as abstract fields.
\end{corollary}
\begin{proof}
An order isomorphism sends a set cofinal presentation to a set cofinal presentation and preserves its two least cardinalities. \Cref{fkc:sb:thm:gapcharacter} distinguishes the spectra. An abstract isomorphism of real closed fields preserves order because positivity is exactly being a nonzero square.
\end{proof}

This theorem does not bound the support of the exponents. For example, $\KN{\aleph_1}\R$ contains all $\omega^a$, even for exponents $a$ of arbitrarily large birthday, while missing the single normal form $h_{\aleph_1}$. The omitted cut is a support phenomenon, not a failure to possess a sufficiently large exponent.

\section{Sharp ordered-field saturation}\label{fkc:sb:sec:saturation}

\begin{theorem}[Types with a small collective support]\label{fkc:sb:thm:supporttypes}
Let $A\subseteq\KN\kappa\R$ be a set such that
\[
 \card{\bigcup_{a\in A}\suppw(a)}<\kappa.
\]
Every complete one-variable ordered-field type over $A$ is realized in $\KN\kappa\R$.
\end{theorem}
\begin{proof}
Let $H$ be the rational span of the displayed union. The field $\R((\omega^H))\subseteq\KN\kappa\R$ is a set real closed field containing $A$. By set-sized compactness, the type extends to a type over $\R((\omega^H))$ realized by some $b$ in a real closed extension of $\R((\omega^H))$. To justify consistency explicitly, every finite part of the original type is realized in $\KN\kappa\R$; the elementary diagram of $\R((\omega^H))$ is true there by quantifier elimination. Their union is therefore finitely satisfiable.

If $b$ is algebraic over $\R((\omega^H))$, then $b\in\R((\omega^H))$, since a real closed field has no proper ordered algebraic extension. Otherwise $b$ defines a cut in the set $\R((\omega^H))$, allowing either side to be empty. Fill that cut by $y\in\No$. It is strict at every element of $\R((\omega^H))$, so $y\notin\R((\omega^H))$ and its support is not contained in $H$. \Cref{fkc:sb:lem:outside}, with $\Sigma=H$, produces a number $z$ supported in $H$ and one additional exponent, having the same cut as $y$ on all of $\R((\omega^H))$. Its support has cardinality less than $\kappa$, so $z\in\KN\kappa\R$.

Over a real closed field, a transcendental one-variable type is determined by its cut: signs of polynomials are determined by their ordered real roots, and quantifier elimination then determines every formula. Thus $z$ realizes the chosen extension over $\R((\omega^H))$, and hence the original type over $A$.
\end{proof}

\begin{corollary}[Exact saturation threshold]\label{fkc:sb:cor:saturation}
For every infinite cardinal $\nu$, the field $\KN\kappa\R$ is $\nu$-saturated if and only if $\nu\le\cf\kappa$.
\end{corollary}
\begin{proof}
If $\card A<\mu$, \cref{fkc:sb:lem:card}(b) bounds the union of its supports by a cardinal below $\kappa$, and \cref{fkc:sb:thm:supporttypes} applies. Successive realization gives finite-variable saturation.

For failure, take the $\mu$ lower and $\mu$ upper brackets of a missing prefix from \cref{fkc:sb:thm:gapcharacter}. The partial type requiring an element to lie above all lower brackets and below all upper brackets is finitely satisfiable in $\KN\kappa\R$ but omitted. It extends to a complete type over this parameter set of size $\mu$. Hence the field is not $\mu^+$-saturated, and cannot be $\nu$-saturated for any larger cardinal $\nu$.
\end{proof}
\mergetag{} This is the saturation theorem that non-claim N5 of Part~I disclaims, proved here for the class field $\KN\kappa\R$ and not for the set fields $\K$. It has the shape of the across-universes report's spectrum theorem (\lbl{univ:cor:spectrum}) for different fields.

\begin{remark}[The countable-cofinality case]\label{fkc:sb:rem:countablecf}
When $\cf\kappa=\aleph_0$, the conclusion includes genuine $\aleph_0$-saturation over \emph{finite parameter sets}, not merely the ability to fill finite order cuts. A finite parameter set can generate a countable field and determine infinitely many polynomial conditions. The small exponent field and the first-outside-exponent lemma handle all these conditions together. The omitted countable cut uses \emph{countably many parameters}; it does not contradict this assertion.
\end{remark}

\section{Strong sums and spherical completeness}\label{fkc:sb:sec:spherical}

\subsection{Ambient strong summation}
\begin{definition}\label{fkc:sb:def:strong}
A set family $(x_i)_{i\in I}$ of surreal normal forms is \emph{strongly summable} if the union of its supports is reverse well-ordered and each exponent occurs in only finitely many of the supports. Its strong sum is obtained by adding the finitely many coefficients at each exponent. Cancellation is allowed.
\end{definition}

This is a support operation. It is not defined as the limit of partial sums in the order or valuation topology. The distinction is indispensable at proper-class rank.

\begin{theorem}[Exact strong-sum threshold]\label{fkc:sb:thm:sums}
Every strongly summable family of fewer than $\mu$ elements of $\KN\kappa\R$ has its sum in $\KN\kappa\R$. There is a strongly summable family of $\mu$ elements of $\KN\kappa\R$ whose sum is not in $\KN\kappa\R$. The same assertions hold for $\KN\kappa\C$ and for families in $\Ok$ whose ambient sums are omnific.
\end{theorem}
\begin{proof}
The support of a sum is contained in the union of the supports of its summands. \Cref{fkc:sb:lem:card}(b) proves the positive assertion.

Choose a continuous increasing cofinal sequence $(\beta_\xi)_{\xi<\mu}$ in $\kappa$, starting with $\beta_0=0$. Partition $\kappa$ into the successive intervals $[\beta_\xi,\beta_{\xi+1})$, and put
\[
 b_\xi=\sum_{\beta_\xi\le\alpha<\beta_{\xi+1}}\omega^{-\alpha}.
\]
Each $b_\xi$ has fewer than $\kappa$ terms. Their supports are disjoint, the union is reverse well-ordered, and their strong sum is $h_\kappa$. This gives the negative assertion for the fields. Multiplying every summand by $\omega^\kappa$ gives purely infinite omnific summands, with sum
\begin{equation}\label{fkc:sb:eq:zkappa}
 z_\kappa=\sum_{\alpha<\kappa}\omega^{\kappa-\alpha}\notin\Ok.
\end{equation}
All the exponents $\kappa-\alpha$ are strictly positive in the surreal ordered field. The complex case contains the real example.
\end{proof}
\mergetag{} Part~I has no strong sums. For a singular $\kappa$ the negative half answers, in the class setting, the case that the omnific-automorphisms report leaves open after \lbl{opa:as:prop:kappa}: that proposition characterizes summability in a set-sized $K_{<\kappa}$ for regular $\kappa$ by fewer than $\kappa$ nonzero members plus pointwise finiteness, and says that for singular $\kappa$ ``the union argument fails, and no extension is claimed''. The family $(b_\xi)_{\xi<\mu}$ has $\mu<\kappa$ members, pairwise disjoint supports and a well-ordered union of supports, so it is Hahn summable in the full field, but its sum has $\kappa$ terms; the same construction works in any set-sized $K_{<\kappa}$ whose exponent group contains a strictly monotone sequence of length $\kappa$. So for singular $\kappa$, fewer than $\kappa$ nonzero members and ambient summability do not imply summability in $K_{<\kappa}$. This does not answer \lbl{opa:as:q:singular}, which asks about automorphisms. The large-cardinal report's \lbl{lce:sg:thm:threshold} uses the same series $h_\kappa$ and $z_\kappa$ (its $Y_\kappa$ is another omnific witness) for a different purpose: the failure of its embedding lift to preserve strong sums of size $\kappa$ at a measurable critical point.

\mergetag{} \emph{Related (batch 32).} The large-cardinal report's canonical strong part (\lbl{lce:mf:thm:atomic}) gives a linear companion of the positive assertion. For an uncountable $\nu$ and a $k$-linear map $\Psi$ between Hahn fields that preserves all Hahn sums of fewer than $\nu$ members, the unique strong map with the same monomial values agrees with $\Psi$ on every series with fewer than $\nu$ terms. Taking the source field $\No$ and $\nu=\kappa$, a map on $\No$ that preserves all Hahn sums of fewer than $\kappa$ members is determined on $\KN\kappa\R$ by its monomial values, as that report records. No large cardinal is used in that theorem, and nothing of it is used here.

\subsection{Gluing a nest without taking a topological limit}
For $a\in\KN\kappa\R$ and $\gamma\in\No$, the closed valuation ball is
\[
 B_{\KN\kappa\R}(a,\gamma)
   =\{x\in\KN\kappa\R:v(x-a)\ge\gamma\}.
\]
A nest is a nonempty set of such balls totally ordered by inclusion. Radii are arbitrary surreal values, not restricted to the support of a center. The support bound places no restriction on available radii.

\begin{lemma}[Support-controlled intersection in the ambient field]\label{fkc:sb:lem:glue}
Let $(B_{\No}(a_i,\gamma_i))_{i\in I}$ be a set-indexed nest of nonempty closed valuation balls in $\No$. There is a common point $a$ with
\[
 \suppw(a)\subseteq\bigcup_{i\in I}\suppw(a_i).
\]
The identical assertion holds for complex coefficients in $\SC$.
\end{lemma}
\begin{proof}[Proof \textup{(printed once)}]
This is the first two paragraphs of the proof of \cref{fkc:lem:glue}, without its final cardinal count and read in the decreasing convention; source~02 gives the same argument. Membership in $B(a_i,\gamma_i)$ prescribes precisely the coefficients at growth exponents $e>-\gamma_i$: they must agree with $a_i$. Inclusion of two balls implies compatibility of their prescriptions, since the center of the smaller ball belongs to the larger one; because every value is realized, inclusion also orders the radii in the expected reverse order of ball size. The prescribed coefficients, and zero elsewhere, define a series $a$. Its support is reverse well-ordered: given a nonempty set $U$ of its exponents and $e\in U$, some one center $a_i$ with $e>-\gamma_i$ agrees with $a$ at every exponent above $-\gamma_i$, so $U\cap[e,\infty)$ is a nonempty subset of the reverse well-ordered set $\suppw(a_i)$ and its greatest member is greatest in $U$. The series agrees with every center at all its prescribed exponents, so it belongs to every ball, and no exponent outside the union of center supports was introduced. The proof is unchanged over $\C$.
\end{proof}

\begin{theorem}[Exact spherical threshold]\label{fkc:sb:thm:spherical}
Every nest of fewer than $\mu$ closed valuation balls in $\KN\kappa\R$ has a common point in $\KN\kappa\R$. The least cardinality of a nest with empty intersection is exactly $\mu$. The same is true in $\KN\kappa\C$.
\end{theorem}
\begin{proof}
A nest of balls in $\KN\kappa\R$ gives a nest of the corresponding ambient balls in $\No$: inequalities between the radii and centers express the inclusion, and all values are already realized in $\KN\kappa\R$. Apply \cref{fkc:sb:lem:glue}, then \cref{fkc:sb:lem:card}(b), to obtain a common point with support cardinality less than $\kappa$.

For sharpness choose an increasing cofinal sequence $(\beta_\xi)_{\xi<\mu}$ in $\kappa$ and write
\[
 h_\beta=\sum_{\alpha<\beta}\omega^{-\alpha},\qquad
 B_\xi=B_{\KN\kappa\R}(h_{\beta_\xi},\beta_\xi).
\]
For $\beta<\delta<\kappa$, $v(h_\delta-h_\beta)=\beta$. Hence these balls are nested. If $x$ belonged to every $B_\xi$, its coefficient at $-\alpha$ would be $1$ for every $\alpha<\kappa$: choose $\xi$ with $\alpha<\beta_\xi$, so that this coefficient is prescribed by that ball. This would require $\kappa$ support terms. Thus the intersection in $\KN\kappa\R$ is empty. The ambient point $h_\kappa$ belongs to every ambient ball. The same coefficient argument excludes any common point in $\KN\kappa\C$.
\end{proof}
\mergetag{} This is the class form of \cref{fkc:thm:balls} and of \cref{fkc:thm:main}(iv), printed once: the lower bound is the argument of \cref{fkc:lem:glue}, and the nest $(B_\xi)$ is Part~I's canonical nest $B(a_i,s_{\alpha_i})$ of \eqref{fkc:eq:approximants} for the missing element $h_\kappa$, transported by $t^\gamma\mapsto\omega^{-\gamma}$ (its approximant $a_i=x_{<s_{\alpha_i}}$ is $h_{\beta_\xi}$). As in \cref{fkc:cor:ballfiber}, the intersection of the ambient balls in $\No$ is the filler class \eqref{fkc:sb:eq:fillercone} of $h_\kappa$: the balls prescribe exactly the coefficients at the exponents $-\alpha$, $\alpha<\kappa$.

\begin{remark}[The singular distinction]\label{fkc:sb:rem:singular}
For a regular $\kappa$, the first failing nest has size $\kappa$. For $\kappa=\aleph_\omega$, it already has size $\aleph_0$: each center may use an increasingly large, but still sub-$\kappa$, support. Thus ``each element uses fewer than $\kappa$ terms'' does not imply that one can glue $\kappa$-small collections of such elements. The correct threshold is the additivity cardinal $\cf\kappa$. \mergetag{} Part~I makes the same point for omission numbers after \cref{fkc:cor:tower} and in \cref{fkc:ex:singular}.
\end{remark}

\section{Why ordinary Cauchy completeness sees none of this}\label{fkc:sb:sec:topology}

\begin{proposition}[All set Cauchy nets stabilize]\label{fkc:sb:prop:cauchy}
Every set-indexed valuation-Cauchy net in $\KN\kappa\R$ or $\KN\kappa\C$ is eventually constant. Every set of distinct elements is uniformly discrete for the valuation neighborhoods and is closed.
\end{proposition}
\begin{proof}
Let $X$ be the set of values taken by the net. The nonzero differences of pairs of elements of $X$ have a set of valuations. Choose a surreal $\gamma$ strictly greater than all these valuations. The Cauchy condition at $\gamma$ says that sufficiently late pairwise differences have valuation at least $\gamma$; by its choice, those differences must be zero.

For uniform discreteness apply the same bound to a given set $X$. A ball of that radius around any member contains no other member. To prove closedness at a point $y\notin X$, bound the set of valuations $\{v(y-x):x\in X\}$ and choose a strictly larger radius. The resulting ball at $y$ misses $X$.
\end{proof}

This is the familiar proper-class-rank phenomenon already appearing in the repository's foundations and across-universes reports \cite{CollUniv}. Its application here is useful, but it is not claimed as a new observation. \mergetag{} The statements are \lbl{found:thm:discrete} \cite{CollFound} (for the fine topology, which on $\KN\kappa\R$ has the same open sets as the valuation topology with surreal radii, since $\KN\kappa\R$ contains every monomial) and \lbl{univ:prop:topology} \cite{CollUniv}. In Part~I's terms, the class value group $\No$ has no set-sized cofinal subset, so the analogue for $\Gamma=\No$ of the condition $\cf(\Gamma)=\cf(\kappa)$ of \cref{fkc:thm:dichotomy} fails for every $\kappa$: there are no missing set-Cauchy limits, while the empty nests of \cref{fkc:sb:thm:spherical} remain.

\begin{proposition}[Closed and nowhere dense inclusions]\label{fkc:sb:prop:closed}
For uncountable $\kappa<\lambda$, $\KN\kappa\R$ is closed and nowhere dense in $\KN\lambda\R$ for valuation neighborhoods. It is also closed and nowhere dense in $\No$. The analogous assertions hold for the complexifications.
\end{proposition}
\begin{proof}
Let $y\notin\KN\kappa\R$ and consider its first $\kappa$ growth exponents $a_\alpha$. Choose $\gamma>-a_\alpha$ for all $\alpha<\kappa$. The ball $B(y,\gamma)$ forces agreement at all these exponents, so it misses $\KN\kappa\R$. This proves closedness.

For empty interior, it suffices to treat a ball centered at $x\in\KN\kappa\R$; a ball with no such center already misses $\KN\kappa\R$. Given its radius $\gamma$, choose a growth exponent $b$ strictly below every exponent in $\suppw(x)$ and below $-\gamma$. The tail
\[
 w=\sum_{\alpha<\kappa}\omega^{b-\alpha}
\]
has exactly $\kappa$ terms, is disjointly supported below $x$, and satisfies $v(w)>\gamma$. Thus $x+w$ belongs to the ball but not to $\KN\kappa\R$. It lies in $\KN\lambda\R$ because its support cardinality is $\kappa<\lambda$. Closedness plus empty interior proves nowhere denseness. Use complex coefficients for the same argument in $\KN\kappa\C$.
\end{proof}
\mergetag{} The assertion in $\No$ is \lbl{lce:sg:prop:closed-support} of the large-cardinal report (order topology, every uncountable $\kappa$), proved there by dividing a missing element by a large ordinal; the proof above is a second route, and the assertion in $\KN\lambda\R$ is not in that report. Part~I's non-claim N6 concerns the set fields; for the class fields this proposition is the topology relative to $\No$.

\paragraph{The apparent paradox resolved.}
The empty nest of \cref{fkc:sb:thm:spherical} is not produced by a Cauchy net in the full valuation. Its radii $\beta_\xi<\kappa$ are bounded by the surreal ordinal $\kappa$. The centers never become arbitrarily close at all surreal scales. In particular, the strong sum $h_\kappa$ is not the valuation limit of these partial sums. A proper immediate extension with the same residue field and value group can be nondense. Completeness, spherical completeness, and closure under strong summation are genuinely different requirements here. \mergetag{} Part~I's \cref{fkc:ex:complete} shows the same separation of completeness and spherical completeness for a set-sized value group.

\section{Exponential closure without any surjective exponential}\label{fkc:sb:sec:exponential}

\subsection{The inherited exponential stays inside}
The following classical facts about Gonshor's surreal exponential are used, with provenance in \cite[Theorem~2.16 and Section~2.3]{MM}:
\begin{equation}\label{fkc:sb:eq:classicalexp}
 \exp(\No^\uparrow)=\omega^{\No},\qquad
 \exp(x^\uparrow+r+\epsilon)
   =\exp(x^\uparrow)e^r\sum_{n=0}^{\infty}\frac{\epsilon^n}{n!}.
\end{equation}
Here $\No^\uparrow$ is the additive group of purely infinite normal forms, and the Taylor sum is strongly summable for infinitesimal $\epsilon$. These facts are imported, not claimed as contributions.

\begin{theorem}[One-sided closure and explicit logarithmic escape]\label{fkc:sb:thm:expclosure}
For every uncountable $\kappa$,
\[
 \exp(\KN\kappa\R)\subseteq\KN\kappa\R^{>0},\qquad
 \log(\KN\kappa\R^{>0})\nsubseteq\KN\kappa\R.
\]
More precisely, the monomial $m_\kappa=\exp(z_\kappa)$ belongs to every $\KN\lambda\R$ with $\lambda>\aleph_0$, whereas $\log(m_\kappa)=z_\kappa\notin\KN\kappa\R$.
\end{theorem}
\begin{proof}
Split $x\in\KN\kappa\R$ as in \eqref{fkc:sb:eq:split}. The first exponential factor in \eqref{fkc:sb:eq:classicalexp} is a single monomial, the second is real, and the support of the Taylor sum is contained in the additive monoid generated by $\suppw(\epsilon)$. That monoid has cardinality at most $\max(\card{\suppw(\epsilon)},\aleph_0)<\kappa$. Multiplication by a monomial translates the support and does not increase its cardinality. This proves exponential closure, including when $\cf\kappa=\aleph_0$.

The element $z_\kappa$ of \eqref{fkc:sb:eq:zkappa} is purely infinite with exactly $\kappa$ terms. By \eqref{fkc:sb:eq:classicalexp}, its exponential is a single Conway monomial. The logarithm is its inverse on the full surreal field, proving the asserted escape.
\end{proof}
\mergetag{} Part~I's non-claim N7 disclaims closure under the surreal exponential for its set fields; this theorem is that closure for the class field. The large-cardinal report proves the same one-sided closure, with the same witness $z_\kappa$, only when $\kappa$ is the critical point of an elementary embedding, so measurable (\lbl{lce:cor:expclosure}); its proof goes through the embedding. The proof above needs no large cardinal and works for every uncountable $\kappa$, regular or singular.

Gonshor's more precise normal-form formula makes the image explicit. The increasing bijection $g:\No^{>0}\to\No$ satisfies
\begin{equation}\label{fkc:sb:eq:gonshor-reindex}
 \exp\left(\sum_\alpha r_\alpha\omega^{a_\alpha}\right)
 =\omega^{\sum_\alpha r_\alpha\omega^{g(a_\alpha)}}
 \qquad(a_\alpha>0).
\end{equation}
This is another classical part of the same theorem \cite[Theorem~2.16]{MM}. In particular, the coefficient-preserving reindexing in the exponent does not change support cardinality.

\begin{proposition}[Exact inherited image]\label{fkc:sb:prop:expimage}
For every $y\in\KN\kappa\R^{>0}$,
\[
 \log y\in\KN\kappa\R\quad\Longleftrightarrow\quad
 \lead(y)\in\KN\kappa\R.
\]
Consequently,
\begin{equation}\label{fkc:sb:eq:expimage}
 \exp(\KN\kappa\R)=
 \{y\in\KN\kappa\R^{>0}:\lead(y)\in\KN\kappa\R\},\qquad
 \exp(\Pk)=\omega^{\KN\kappa\R}.
\end{equation}
The reindexing in \eqref{fkc:sb:eq:gonshor-reindex} restricts to an ordered additive-group isomorphism $\Pk\cong\KN\kappa\R$.
\end{proposition}
\begin{proof}
Write $y=r\omega^b(1+\epsilon)$ with $r>0$ real, $b=\lead(y)$, and $\epsilon$ infinitesimal. Dividing by the leading monomial and removing the constant term preserves the support bound, so $\epsilon\in\KN\kappa\R$. Its logarithmic Taylor series is supported in the finite-sum monoid of its support, hence belongs to $\KN\kappa\R$. Now
\[
 \log y=\log(\omega^b)+\log r+\log(1+\epsilon).
\]
The first summand is purely infinite and the last two are finite, so there can be no support cancellation between them. By the inverse of \eqref{fkc:sb:eq:gonshor-reindex}, $\log(\omega^b)$ has exactly as many nonzero terms as $b$: replace each growth exponent $b_\alpha$ of $b$ by $g^{-1}(b_\alpha)>0$ and keep its coefficient. Thus the displayed logarithm belongs to $\KN\kappa\R$ exactly when $b$ does. This proves the image assertions. The same reindexing is additive and order-preserving, and its inverse preserves the support bound, proving the last claim. It need not be multiplicative.
\end{proof}

These formulas already rule out surjectivity of the \emph{inherited} exponential. The next argument is stronger: no other choice of exponential can repair surjectivity on this ordered field.

\subsection{An intrinsic obstruction}
\begin{definition}\label{fkc:sb:def:fin}
For an ordered field $F$, write
\[
 \Fin(F)=\{x\in F:|x|<n\text{ for some }n\in\N_{>0}\},
\]
and let
\[
 U(F)=\{y\in F^{>0}:n^{-1}<y<n
                       \text{ for some }n\in\N_{>0}\}.
\]
These are respectively a convex additive subgroup and a convex subgroup of the positive multiplicative group. The latter is the group of positive units of the natural valuation ring.
\end{definition}

\begin{lemma}[Normalization of any hypothetical exponential]\label{fkc:sb:lem:normalize}
If an ordered field $F$ admits an ordered-group isomorphism $\mathcal E:(F,+)\to(F^{>0},\cdot)$, it admits one $\mathcal E_2$ with $\mathcal E_2(1)=2$. Every such $\mathcal E_2$ maps $\Fin(F)$ onto $U(F)$ and induces an ordered-group isomorphism
\[
 F/\Fin(F)\ \cong\ F^{>0}/U(F).
\]
\end{lemma}
\begin{proof}
Let $a=\mathcal E^{-1}(2)>0$ and put $\mathcal E_2(x)=\mathcal E(ax)$. Multiplication by $a$ is an ordered additive automorphism, so this is the required normalization. If $|x|<n$, then $2^{-n}<\mathcal E_2(x)<2^n$, and $\mathcal E_2(x)$ is a positive finite unit. Conversely, if $y$ and $y^{-1}$ are bounded by an integer, choose $n$ with $2^{-n}<y<2^n$. Surjectivity gives $y=\mathcal E_2(x)$; monotonicity then gives $-n<x<n$. Convexity makes the induced quotient isomorphism an ordered one.
\end{proof}

\begin{lemma}[The two quotients do not match]\label{fkc:sb:lem:quotientmismatch}
For $F=\KN\kappa\R$, the additive quotient $F/\Fin(F)$ is isomorphic as an ordered group to $\Pk$, while $F^{>0}/U(F)$ is isomorphic as an ordered group to $(\No,+)$. The first has a set-presented omitted cut and the second has none.
\end{lemma}
\begin{proof}
The purely infinite truncation $x\mapsto x^\uparrow$ is an additive surjection $\KN\kappa\R\to\Pk$, with kernel precisely the finite elements. It identifies the ordered quotient with $\Pk$. For the multiplicative quotient, the greatest growth exponent gives a surjective homomorphism
\[
 \lead:\KN\kappa\R^{>0}\longrightarrow(\No,+).
\]
Its kernel consists exactly of elements with leading exponent zero, namely $U(\KN\kappa\R)$. Its order on cosets is the quotient order: a strictly positive growth-exponent difference means that the ratio is positive infinite. Hence it induces the second isomorphism.

Use the prefix $z_\kappa$ of \eqref{fkc:sb:eq:zkappa}. The lower and upper brackets of \eqref{fkc:sb:eq:brackets} now have only positive growth exponents, so they belong to $\Pk$. Exactly the first-difference argument of \cref{fkc:sb:sec:gaps} shows that they present an omitted cut of $\Pk$, of character $(\mu,\mu)$. By contrast, the surreal set-cut property precludes any set-presented gap in $\No$.
\end{proof}

\begin{theorem}[No surjective ordered exponential]\label{fkc:sb:thm:noexp}
There is no ordered-group isomorphism
\[
 (\KN\kappa\R,+,0,<)\longrightarrow(\KN\kappa\R^{>0},\cdot,1,<).
\]
No compatibility with the real exponential, monomials, valuation, or normal-form summation is assumed in this nonexistence assertion.
\end{theorem}
\begin{proof}
Normalize a hypothetical isomorphism by \cref{fkc:sb:lem:normalize}. It would identify the two quotients of \cref{fkc:sb:lem:quotientmismatch}. An ordered-group isomorphism preserves set-presented gaps, contradicting that lemma.
\end{proof}

\begin{corollary}[A first-order separation of exponential expansions]\label{fkc:sb:cor:expsep}
The inclusion of exponential structures
\[
 (\KN\kappa\R,\exp|_{\KN\kappa\R})\subseteq(\No,\exp)
\]
is not elementary, although the ordered-field inclusion is elementary. The sentence
\[
 \forall y\,(y>0\ \longrightarrow\ \exists x\ \exp(x)=y)
\]
holds in the latter and fails in the former.
\end{corollary}
\mergetag{} The large-cardinal report proves this corollary at a measurable critical point with the same sentence (\lbl{lce:cor:expclosure}). With \cref{fkc:sb:thm:expclosure} it bears on the second further question of Part~I (\cref{fkc:sec:audit}, richer languages) in the class setting: the base field is closed under the named exponential, and the comparison with $\No$ is then not elementary.

The obstruction has the same general additive-versus-multiplicative quotient philosophy as the classical theory of exponential Hahn fields \cite{KKS,KS}. The specific mismatch here is an exact set-cut invariant inside an exp-closed, proper-class subfield with unrestricted surreal monomials. The theorem does not contradict the $\kappa$-bounded exponential constructions in \cite{KS}: those arrange their exponent groups differently. Nor is it a counterexample to the positive results for omega-fields in \cite{BKMM}; our ordered value group is $\No$, whereas the additive group of $\KN\kappa\R$ has omitted set cuts.

\subsection{One logarithm recovers any surreal after a monomial shift}
\begin{theorem}[Logarithmic hull collapse]\label{fkc:sb:thm:loghull}
Let $L\subseteq\No$ be a subfield containing every Conway monomial. If $\log(y)\in L$ for every $y\in L^{>0}$, then $L=\No$. In fact every $x\in\No$ has a representation
\begin{equation}\label{fkc:sb:eq:onelog}
 x=\omega^{-A}\log(\omega^B)
\end{equation}
for suitable surreal exponents $A,B$.
\end{theorem}
\begin{proof}
For $x\ne0$, its support is a set, so choose $A>0$ with $A+a>0$ for all $a\in\suppw(x)$. Then $b=\omega^A x$ is purely infinite. Its exponential is a monomial $\omega^B$ by \eqref{fkc:sb:eq:classicalexp}, and taking logarithms gives \eqref{fkc:sb:eq:onelog}. For $x=0$, take $A=B=0$. A field $L$ as in the statement contains both monomials in this expression and the required logarithm, hence every $x$.
\end{proof}
\mergetag{} The letter $L$ here is source~02's, a subfield of $\No$; it is not Part~I's completion $L_{\kappa,k}$. The large-cardinal report proves the same collapse for its short field: the least logarithmically closed field containing $\mathscr S_{<\kappa}$ is all of $\No$, by the same monomial shift (the last assertion of \lbl{lce:thm:twoclosures}, with the representation \lbl{lce:eq:singlelog}). It is stated there at a measurable critical point, but the argument uses only that the field contains every monomial.

\begin{remark}[The all-monomials hypothesis cannot be dropped]\label{fkc:sb:rem:allmonomials}
There are many proper subfields of the surreal field closed under both exponential and logarithm. Their constructions restrict the exponents or use birthday and transseries hierarchies; see \cite{BG}. They do not contain every $\omega^a$ for arbitrary $a\in\No$. \Cref{fkc:sb:thm:loghull} addresses precisely this stronger all-monomials requirement, not the general existence of exp--log closed subfields.
\end{remark}

\section{Omnific integer parts with identical set-sized quotients}\label{fkc:sb:sec:omnific}

\subsection{Integer parts, fraction fields, and the constant term}
Put
\[
 \Pk=\{x\in\KN\kappa\R:\suppw(x)\subseteq\No^{>0}\}.
\]
It is an additive group, an $\R$-vector space, and an ideal of $\Ok=\Z+\Pk$. The vector-space assertion does not mean that $\Ok$ contains all real constants: its constant coefficients are integers.

\begin{proposition}[The elementary arithmetic package]\label{fkc:sb:prop:integerpart}
The ring $\Ok$ has the following properties.
\begin{enumerate}[label=\textup{(\alph*)}]
\item It is a discretely ordered integer part of $\KN\kappa\R$: every $x\in\KN\kappa\R$ has a unique $a\in\Ok$ with $a\le x<a+1$.
\item $\Frac(\Ok)=\KN\kappa\R$ and its only units are $1,-1$.
\item The integer constant term is a split surjective ring map $\ct:\Ok\to\Z$, with kernel $\Pk$.
\item For each nonzero integer $n$, $\Pk=n\Pk$ and $\Ok/n\Ok\cong\Z/n\Z$.
\end{enumerate}
\end{proposition}
\begin{proof}
Write $x=x^\uparrow+r+\epsilon$ as in \eqref{fkc:sb:eq:split}, and put $n=\lfloor r\rfloor\in\Z$. If $r$ is not an integer, the infinitesimal $\epsilon$ is smaller in magnitude than both positive real margins to $n$ and $n+1$, so $x^\uparrow+n$ is a floor. If $r=n$ is an integer, the floor is $x^\uparrow+n$ when $\epsilon\ge0$, and $x^\uparrow+n-1$ when $\epsilon<0$. All these elements lie in $\Ok$.

An element of $\Ok$ bounded by a real number is an ordinary integer: any nonzero purely infinite part has infinite magnitude. Thus $1$ is the least positive element, which proves discreteness and uniqueness of the floor. The negative-infinitesimal boundary case is essential; for example, $\lfloor-\omega^{-1}\rfloor=-1$.

For the fraction assertion, choose $A>0$ so large that $A+a>0$ for every $a\in\suppw(x)$. Then both $\omega^A x$ and $\omega^A$ belong to $\Ok$, and $x=(\omega^A x)/\omega^A$. A nonconstant element of $\Ok$ has strictly positive greatest growth exponent. Greatest exponents add under multiplication, while all nonzero elements of $\Ok$ have nonnegative greatest exponent. A unit must therefore be an ordinary integer, and hence must be $\pm1$.

Positive growth exponents cannot add to zero, so the constant term of a product in $\Ok$ is the product of the two integer constant terms. Inclusion of $\Z$ splits this map. Finally $\Pk$ is a rational vector space, so division by $n$ stays in $\Pk$. This proves (d) and the claimed quotient.
\end{proof}

These statements follow the standard normal-form arithmetic already present in the supplied project. The following uniform theorem for the smaller rings requires an additional argument. \mergetag{} Parts~(a) and~(b) are, for every uncountable $\kappa$, the large-cardinal report's \lbl{lce:sg:prop:integerpart} and the fraction identity \lbl{lce:eq:Ozfrac} (the latter stated there at a measurable critical point, with the same shift argument).

\subsection{A scaled reservoir of countably supported elements}
We write $|u|\ll d$ for $n|u|<d$ for every positive integer $n$.

\begin{lemma}[Arbitrarily large set fields at an arbitrarily small scale]
\label{fkc:sb:lem:reservoir}
Let $d\in\No^{>0}$ and let $\rho$ be an infinite set cardinal. There is a set subfield $W\subseteq\No$ of cardinality at least $\rho$ with the following properties: every member has at most countable normal-form support; all its growth exponents lie in a divisible set group $H$ such that $|u|\ll d$ for every $u\in H$; and, on setting $\delta=\omega^d$,
\[
 \delta e\in\Pk\quad(e\in W),\qquad
 \frac{\delta}{e}\in\Pk\quad(e\in W\setminus\{0\}).
\]
These conclusions hold simultaneously for every uncountable $\kappa$.
\end{lemma}
\begin{proof}
For $\alpha<\rho$, choose the distinct positive surreal numbers
\[
 d_\alpha=\frac{d}{\omega^{\alpha+1}},\qquad
 H=\Span\{d_\alpha:\alpha<\rho\}.
\]
Each $d_\alpha$ is infinitesimal relative to $d$, and every finite rational linear combination retains this property. In particular, $d+u>0$ for every $u\in H$. Let
\[
 W=\R(\omega^u:u\in H),
\]
the field of rational expressions in these monomials, \emph{not} the full Hahn field on $H$. It is a set field, and its distinct monomials give cardinality at least $\rho$.

An individual rational expression uses only finitely many monomials. To invert a nonzero finite sum, factor out its leading term and write what remains as $1+\eta$, where $\eta$ is a finite-support infinitesimal. The geometric expansion $(1+\eta)^{-1}=\sum_{n\ge0}(-\eta)^n$ is strongly summable by the usual Hahn support lemma. Its support lies in the additive monoid of a finite set, followed by a translate, and is therefore at most countable and contained in $H$. Finite sums and products of such expressions preserve this countable bound. This proves the asserted support property for every member of $W$, including reciprocals.

Multiplication by $\delta$ shifts these supports into $d+H$, which consists of positive exponents. The shifted supports remain countable, hence have cardinality less than every uncountable $\kappa$. This proves the final assertions.
\end{proof}

The rational-expression restriction in this lemma is important. The full Hahn field on a very large $H$ need not lie in $\KN\kappa\R$. The reservoir needs large cardinality as a \emph{set of elements}, but each of those elements individually needs only a countable support.

\begin{theorem}[Uniform universal set-sized quotient]\label{fkc:sb:thm:smallquotient}
Let $S$ be any set-sized unital ring, not necessarily commutative. Every unital ring homomorphism $\chi:\Ok\to S$ has the unique factorization
\[
 \chi=\bar\chi\circ\ct,
 \qquad \bar\chi:\Z\longrightarrow S.
\]
In particular it kills all of $\Pk$.
\end{theorem}
\begin{proof}
Fix a nonzero $x\in\Pk$. Its growth exponents form a set of positive surreal numbers. Choose $d>0$ with $2d<a$ for every $a\in\suppw(x)$. Put $\delta=\omega^d$; then $x/\delta^2\in\Pk$.

Apply \cref{fkc:sb:lem:reservoir} with a cardinal $\rho>\card S$. The function
\[
 W\longrightarrow S,\qquad e\longmapsto\chi(\delta e)
\]
cannot be injective. Choose $e\ne f$ with the same image. Subtraction gives $\chi(\delta(e-f))=0$. Both factors in
\begin{equation}\label{fkc:sb:eq:collision}
 \bigl(\delta(e-f)\bigr)\frac{\delta}{e-f}=\delta^2
\end{equation}
belong to $\Pk$, so $\chi(\delta^2)=0$. Finally,
\[
 \chi(x)=\chi(\delta^2)\chi(x/\delta^2)=0.
\]
This proves that the whole ideal $\Pk$ is killed. Since $\Ok=\Z\oplus\Pk$ additively, the required factor is the restriction of $\chi$ to $\Z$; it is unique and a ring map. No cancellation, reducedness, commutativity, or characteristic assumption on $S$ was used.
\end{proof}

The scaled collision in \eqref{fkc:sb:eq:collision} is the method credited to the repository's universal-quotient report \cite{CollOSQ}. The point of the present proof is that a countably supported reservoir suffices, so the argument already works in \emph{every} $\Ok$. It would be invalid simply to cite the theorem for $\Oz$ and assume that maps on these subrings extend to $\Oz$.

\begin{remark}[The theorem was in the collection at the pin; a third route \mergetag]\label{fkc:sb:rem:osq}
The quotient report's support-threshold theorem (\lbl{osq:thm:support}~(i), present at source~02's pin) states, for every uncountable $\lambda$, every unital $D\subseteq k\in\{\R,\C\}$ and every set-sized unital ring $S$, that restriction gives $\operatorname{Hom}_1(\mathcal A^{<\lambda}_{D,k},S)\cong\operatorname{Hom}_1(D,S)$, where $\mathcal A^{<\lambda}_{D,k}=D+\{f\in\Pi_k:\card{\supp f}<\lambda\}$. For $D=\Z$, $k=\R$ and $\lambda=\kappa$ this ring is $\Ok$, so \cref{fkc:sb:thm:smallquotient} is that theorem, and source~02's description of it as a candidate extension (\cref{fkc:sb:sub:claims}) is stale; see \cref{fkc:sb:sub:corrections}. The quotient report proves it by packing and separation, or by explicit telescopes, and records that its field-collision proof does not apply to the bounded rings, because the Hahn field over its scale subgroup is not contained in them. Source~02's reservoir removes exactly that obstacle: the scaled-field collision \lbl{osq:lem:collision}, applied with the countably supported field $W$ of \cref{fkc:sb:lem:reservoir} in place of that Hahn field, gives a third route to the same theorem.
\end{remark}

\begin{corollary}[All small quotients and all set module actions]\label{fkc:sb:cor:modules}
The nonzero set-sized quotient rings of $\Ok$, up to ring isomorphism, are exactly
\[
 \Z\quad\text{and}\quad\Z/n\Z\quad(n\ge2).
\]
The kernel for the quotient $\Z$ is $\Pk$; for $\Z/n\Z$ it is $n\Ok$. Every unital $\Ok$-module structure on a set is the ordinary abelian-group structure acted on through $\ct$. Its module homomorphisms are exactly the abelian-group homomorphisms.
\end{corollary}
\begin{proof}
A quotient map with set-sized target factors through $\Z$, and its surjectivity makes the target a quotient of $\Z$. Its kernel is the preimage of an ideal of $\Z$. For $n\ne0$, this preimage is $n\Z+\Pk=n\Ok$. For module actions, the action map has set-sized target $\operatorname{End}_{\Z}(M)$, so \cref{fkc:sb:thm:smallquotient} applies. Conversely, any abelian group has exactly this action. Linearity then reduces to additivity.
\end{proof}

\begin{corollary}[Many rings, one set-sized arithmetic image]\label{fkc:sb:cor:rings}
If $\cf\kappa\ne\cf\lambda$, then $\Ok$ and $\Okl\lambda$ are not isomorphic as abstract rings, even though they have the same set-sized quotient rings and the same description of all set module actions.
\end{corollary}
\begin{proof}
A ring isomorphism would extend to an isomorphism of their fraction fields, namely $\KN\kappa\R$ and $\KN\lambda\R$. This contradicts \cref{fkc:sb:cor:noniso}. The identical small-target descriptions are \cref{fkc:sb:thm:smallquotient,fkc:sb:cor:modules}.
\end{proof}

\mergetag{} \emph{Related (batch 32).} The omnific groups report (\nolinkurl{surreal/omnific-groups-and-lattices}) lifts \cref{fkc:sb:thm:smallquotient} from rings to elementary Chevalley groups. By its \lbl{ogl:ch:cor:support} with $\lambda=\kappa$, $D=\Z$ and $k=\R$, for every root system $\Phi$ without an $A_1$ component every homomorphism from $\ker\bigl(E_\Phi(\Ok)\to E_\Phi(\Z)\bigr)$, or from the corresponding Steinberg kernel, to a set-sized group is trivial. So the constant term makes $E_\Phi(\Z)$ the universal set-sized quotient of $E_\Phi(\Ok)$, as that report notes, and likewise for the Steinberg groups. Its proof uses \lbl{osq:thm:support}~(i), which contains \cref{fkc:sb:thm:smallquotient} (\cref{fkc:sb:rem:osq}).

\subsection{Gaussian omnific rings}\label{fkc:sb:sub:gaussian}
Let $\Ok[i]\subseteq\KN\kappa\C$. Coefficientwise constant extraction gives a split ring epimorphism
\[
 \ct_\C:\Ok[i]\longrightarrow\Z[i].
\]
Every set-sized unital ring map on $\Ok[i]$ factors uniquely through $\ct_\C$: its restriction to $\Ok$ kills $\Pk$, and then it kills $i\Pk$ as well. Thus its set-sized quotients are exactly the quotients of $\Z[i]$, with their evident existence obtained from constant extraction. Moreover, $\Frac(\Ok[i])=\KN\kappa\C$. \mergetag{} The factorization is also \lbl{osq:thm:support}~(i) with $D=\Z[i]$ and $k=\C$.

We do \emph{not} infer that the pure rings $\Ok[i]$ are pairwise nonisomorphic from their complex fraction fields: the next section shows that those pure fraction fields are all isomorphic. The real forms, or conjugation as additional structure, retain the distinction.

\section{An intrinsic proper-class family of surcomplex real forms}\label{fkc:sb:sec:realforms}

\subsection{Pure algebraically closed class fields}
\begin{lemma}[Class back-and-forth over a prescribed set]\label{fkc:sb:lem:backforth}
Under Global Choice, two proper-class algebraically closed fields of characteristic zero are isomorphic over any prescribed isomorphism between set-sized algebraically closed subfields.
\end{lemma}
\begin{proof}
Global Choice provides set-like class enumerations of the two fields: choose a rank-respecting global well-order and enumerate its restrictions. Construct, by recursion along $\On$, compatible isomorphisms between set-sized algebraically closed subfields. Start with the prescribed isomorphism.

At a forth step, take the first element of the left field not yet in the domain. It is transcendental over the current domain, because that domain is algebraically closed. The right field has an element transcendental over the current range: the elements algebraic over a set field form a set, whereas the field is a proper class. Map the new transcendental to such an element, extend the rational-function-field isomorphism, and then extend to the relative algebraic closures in the two algebraically closed fields. This is the usual set-sized uniqueness and extension theorem for algebraic closures. The back step is symmetric. Global Choice selects one of the permitted choices at each step.

At a limit set ordinal take the union; it is again a set algebraically closed field because every polynomial uses finitely many coefficients. Every stage and every partial map is a set. Ordinary set recursion with class parameters supplies every set initial segment; their uniquely specified compatible union is a class map. The back-and-forth enumerations make that map a total surjective field isomorphism.
\end{proof}

This is the standard class-homogeneity mechanism used in the repository's surcomplex automorphism and across-universes reports \cite{CollUniv}. It is included to display the exact choice and size requirements, not as a claim of a new back-and-forth theorem. \mergetag{} In the automorphism report it is \lbl{saut:thm:acfhomogeneity} \cite{CollSAUT}, which extends every isomorphism between set-sized subfields, algebraically closed or not. \emph{Related (batch 34).} The set-sized-quotients report prints the case of prime subfields as \lbl{osq:cr:cor:acfiso} \cite{CollOSQ}: every proper-class algebraically closed field of characteristic zero is isomorphic to $\SC$ as a pure field. It credits this lemma and the automorphism report's theorem, and uses the one-sided form \lbl{osq:cr:thm:universalfield} to embed class residue fields of omnific rings in $\SC$.

\begin{proposition}[Pure-field isomorphism with the surcomplex field]\label{fkc:sb:prop:classiso}
For every uncountable $\kappa$ there is a field isomorphism
\[
 \psi_\kappa:\KN\kappa\C\overset{\sim}{\longrightarrow}\SC
\]
fixing every member of $\C$. The pure field $\KN\kappa\C$ realizes every set-parameter type in finitely many variables.
\end{proposition}
\begin{proof}
Apply \cref{fkc:sb:lem:backforth} starting with the identity of $\C$. For saturation in one variable, over the algebraic closure of any set of parameters, a type in an algebraically closed field is either an algebraic type or the transcendental type. Algebraic roots are available and a transcendental element exists outside that set algebraic closure. Iterating proves the finite-variable statement. These are the standard quantifier-elimination consequences for algebraically closed fields.
\end{proof}
\mergetag{} This is the class counterpart of Part~I's pure-field collapse (\cref{fkc:thm:pure}): without the valuation, or here without conjugation, the support bound is invisible.

\subsection{Transported involutions and their invariants}
Let $\conj_\kappa(a+ib)=a-ib$ on $\KN\kappa\C$, and define
\begin{equation}\label{fkc:sb:eq:tau}
 \tau_\kappa=\psi_\kappa\conj_\kappa\psi_\kappa^{-1}.
\end{equation}
Then $\tau_\kappa^2=1$ and $\Fix(\tau_\kappa)=\psi_\kappa(\KN\kappa\R)$. Since $\psi_\kappa$ fixes $\C$, the involution restricts to ordinary conjugation there.

\begin{theorem}[Proper-class many nonconjugate involutions]\label{fkc:sb:thm:involutions}
If $\cf\kappa\ne\cf\lambda$, the involutions $\tau_\kappa$ and $\tau_\lambda$ are not conjugate by an automorphism of the pure field $\SC$. None is conjugate to standard surreal conjugation. In particular, restricting $\kappa$ to uncountable regular cardinals gives a proper-class-indexed family of pairwise nonconjugate involutions.
\end{theorem}
\begin{proof}
If $\sigma\tau_\kappa\sigma^{-1}=\tau_\lambda$, then $\sigma$ restricts to a field isomorphism of their fixed fields. Positivity in a real closed field is definable as nonzero-squarehood, so that restriction is an order isomorphism. \Cref{fkc:sb:cor:noniso} excludes it. Standard conjugation has fixed field $\No$, which has no set-presented omitted cuts, whereas every $\KN\kappa\R$ has such a cut. There are proper-class many uncountable regular cardinals, and the construction can be made uniformly using the fixed global choices.
\end{proof}

\begin{remark}[Credits and the across-universes question \mergetag]\label{fkc:sb:rem:univ}
(1) The non-conjugacy to standard conjugation also follows from the automorphism report's criterion \lbl{saut:prop:conjugacycriterion}: an involution of $\SC$ is conjugate to $\conj$ if and only if its fixed field has the set-cut interpolation property, and $\KN\kappa\R$ fails it by \cref{fkc:sb:thm:gapcharacter}.

(2) The across-universes report asks (\lbl{univ:q:families}) which external saturation spectra, and which infinite collections of regular thresholds or entire gap spectra, can be realized simultaneously by fixed fields of involutions of one $\SC$ that agree on $\C$ and have no set-sized cofinal subset; its finite-threshold version is answered there by a finite forcing tower (\lbl{univ:cor:finitetower}). Read in a single universe, with the threshold of a real form taken to be the least cofinality of a set-presented gap of its fixed field (the invariant of \lbl{univ:cor:recover}), \cref{fkc:sb:thm:involutions} answers the clause on infinite collections of regular thresholds, without forcing: for every class $\mathcal C$ of infinite regular cardinals, the involutions $\tau_\nu$ ($\nu\in\mathcal C$ uncountable) together with $\tau_{\aleph_\omega}$ (if $\aleph_0\in\mathcal C$) are pairwise nonconjugate, agree with $\conj$ on $\C$, have fixed fields isomorphic to fields $\KN\kappa\R$ that contain every monomial and hence have no set-sized cofinal subset, and have thresholds exactly the members of $\mathcal C$, by \cref{fkc:sb:thm:gapcharacter}. By \cref{fkc:sb:cor:saturation} their saturation spectra are the initial segments $\{\nu':\nu'\le\nu\}$, and by \cref{fkc:sb:thm:gapcharacter} each gap spectrum is the single pair $\{(\nu,\nu)\}$. The other clauses (arbitrary saturation spectra, gap spectra with more than one character, and families of old fields $\No^M$ of inner models) are not answered, and no real form is classified. This needs Global Choice, as does the rest of this section.

(3) The involutions $\tau_\kappa$ are pure-field involutions, not claimed to preserve the valuation, the Gaussian omnific integers or Hahn sums (non-claim N18). The batch-32 addition to the automorphism report classifies instead the strong valued involutions preserving $\Oz[i]$; its fixed fields are full Hahn fields over real closed coefficient fields. The two families are complementary, and neither result implies the other. \emph{Related (batch 32).} In that report, two nonidentity involutions in the group of valued automorphisms preserving $\Oz[i]$ are conjugate there exactly when their restrictions to $\C$ are conjugate in $\operatorname{Aut}(\C)$ (\lbl{saut:fs:prop:conjugacy}); its $2^{\mathfrak c}$ conjugacy types (\lbl{saut:fs:thm:main-count}) therefore all come from the coefficient involution. Every $\tau_\kappa$ restricts to ordinary conjugation on $\C$, as $\conj$ does, yet none is conjugate to $\conj$ (\cref{fkc:sb:thm:involutions}); since $\conj$ lies in that group, no $\tau_\kappa$ is both valued and $\Oz[i]$-preserving, as that report also records. The pairwise nonconjugacy of the $\tau_\kappa$ is thus invisible on $\C$, whereas the $2^{\mathfrak c}$ types there are detected on $\C$.
\end{remark}

\begin{proposition}[Same elementary theory, different saturation]\label{fkc:sb:prop:sametheory}
All the structures $(\KN\kappa\C,\conj_\kappa)$, all their transported copies $(\SC,\tau_\kappa)$, and $(\SC,\conj)$ satisfy the same first-order sentences in the field language with an involution. Nevertheless, $(\KN\kappa\C,\conj_\kappa)$ is $\nu$-saturated exactly for infinite $\nu\le\cf\kappa$.
\end{proposition}
\begin{proof}
The fixed field is definable by $\conj(x)=x$, and its order is definable by squares. After choosing an element $i$ with $i^2=-1$ and $\conj(i)=-i$, every element is uniquely $a+ib$ with $a,b$ in that fixed field. Field operations and involution become polynomial operations on pairs. Thus a sentence translates into a sentence about a real closed field; completeness of the theory of real closed fields gives elementary equivalence. The two choices of $i$ are interchanged by conjugation, so the unmarked structures have the same conclusion.

A parameter set of cardinality less than $\nu$ supplies fewer than $\nu$ real and imaginary coordinates (with finitely many coordinates when the parameter set is finite). Types in the involutive field reduce to finite tuples of types over these coordinates in $\KN\kappa\R$. \Cref{fkc:sb:cor:saturation} gives the positive assertion. For failure, use the omitted cut type there and require its variable to lie in the fixed field. It is finitely satisfiable and remains omitted.
\end{proof}
\mergetag{} This bears on the named-conjugation paragraph of \cref{fkc:sec:surreal} and on Part~I's second further question: for the class fields, the expansion by conjugation is interpreted in the real closed fixed field, so all these structures have one complete theory, while the support bound reappears in saturation. It does not classify one-types in that language over the set fields $\KC$. The across-universes report proves the analogous elementary equivalence for its transported forms (\lbl{univ:thm:involutions}).

\paragraph{How this extends, rather than repeats, the repository.}
The across-universes report already constructs nonconjugate real forms using old surreal fields in different set-theoretic universes and studies their omitted-cut spectra \cite{CollUniv}. The contribution proposed here is an \emph{intrinsic, single-universe support hierarchy}: all set-presented gaps are classified by first-$\kappa$ normal-form prefixes, and uncountable regular support bounds yield a proper-class family without forcing or large-cardinal hypotheses. This is not a claim to have first discovered nonconjugate surreal real forms.

\section{Examples and a synthesis of the invariants}\label{fkc:sb:sec:examples}

\begin{example}[Countable supports]\label{fkc:sb:ex:countable}
$\KN{\aleph_1}\R$ contains exactly the surreal normal forms with countable or finite support. It is real closed, $\aleph_1$-saturated, and closed under countable ambient strong sums. Every countable nest of closed valuation balls intersects. A nest of size $\aleph_1$ can fail, and every set-presented gap has character $(\aleph_1,\aleph_1)$. It is closed under the surreal exponential but has no surjective ordered exponential. Every set-sized ring representation of $\Okl{\aleph_1}$ factors through $\Z$.
\end{example}

\begin{example}[A larger regular bound]\label{fkc:sb:ex:regular}
In $\KN{\aleph_2}\R$ the missing series $h_{\aleph_1}$ from the previous example is present. No gap of character $(\aleph_1,\aleph_1)$ remains, but the series $h_{\aleph_2}$ creates one of character $(\aleph_2,\aleph_2)$. The two fields and their omnific integer parts are nonisomorphic, even as abstract fields and rings respectively. Their complexifications are abstractly isomorphic and their set-sized omnific quotients are identical.
\end{example}

\begin{example}[A singular bound]\label{fkc:sb:ex:singular}
Let $\kappa=\aleph_\omega$. This field contains each $\KN{\aleph_n}\R$ for finite $n\ge1$, yet its least omitted gap character is $(\aleph_0,\aleph_0)$. Take an increasing cofinal sequence of cardinals $\beta_n<\aleph_\omega$ and the corresponding centers $h_{\beta_n}$. Their balls form an empty countable nest. Each center is allowed because its support is below $\aleph_\omega$; a common point would need their full union. Thus increasing the support allowance can \emph{decrease} the saturation and spherical-intersection thresholds. The field remains $\aleph_0$-saturated over finite parameter sets and remains closed under the inherited exponential.
\end{example}
\mergetag{} \Cref{fkc:sb:ex:singular} is the class form of Part~I's \cref{fkc:ex:singular,fkc:ex:complete}, whose surreal realization uses $X_\kappa=\sum_{\alpha<\kappa}\omega^{-\omega^\alpha}$ \eqref{fkc:eq:actualx} inside a set workspace; here the missing element is $h_\kappa=\sum_{\alpha<\kappa}\omega^{-\alpha}$ and the ambient is all of $\No$. Both have exactly $\kappa$ terms.

\begin{center}
\small
\begin{tabular}{@{}p{0.49\textwidth}p{0.43\textwidth}@{}}
\toprule
Invariant or property of $\KN\kappa\R$ & Value or conclusion \\
\midrule
Normal-form support allowance & strictly less than $\kappa$ \\
Natural value group; residue field & $\No$; $\R$ \\
Set-presented omitted gap spectrum & $\{(\cf\kappa,\cf\kappa)\}$ \\
Largest saturation cardinal & $\cf\kappa$ \\
Least failing strong-sum family size & $\cf\kappa$ \\
Least empty closed-ball nest size & $\cf\kappa$ \\
Set-indexed Cauchy nets & all eventually constant \\
Inherited surreal exponential & defined internally, not onto \\
Any surjective ordered exponential & impossible \\
Set-sized quotients of $\Ok$ & exactly quotients of $\Z$ \\
Pure complexification & isomorphic to $\SC$ under Global Choice \\
Complexification with conjugation & remembers the fixed-field gap spectrum \\
\bottomrule
\end{tabular}
\end{center}

\paragraph{A single mechanism behind several thresholds.}
A missing point requires $\kappa$ coefficients. A family of fewer than $\cf\kappa$ individually sub-$\kappa$ supports cannot force that many coefficients. At size $\cf\kappa$, successive prefixes can exhaust $\kappa$ terms. This common counting mechanism governs the gap, strong-summation, and spherical results. Saturation additionally needs the small exponent field and the first-outside-exponent lemma; the exponential obstruction additionally compares convex quotients. The small-target arithmetic theorem goes in the opposite direction: it uses arbitrarily many \emph{different} countably supported elements, not one element with a very long support.

\section{Further research questions and proposed routes}\label{fkc:sb:sec:questions}

The following questions are not asserted to be previously published open problems. They are explicit next problems arising from the theorems above. Where a question involves a wider literature, its priority and current status require a dedicated review. \mergetag{} All eight remain open; the notes after some of them record related material in Part~I and in the collection.

\begin{question}[Isomorphism at equal cofinality]\label{fkc:sb:q:equalcf}
For uncountable $\kappa,\lambda$, does $\cf\kappa=\cf\lambda$ imply $\KN\kappa\R\cong\KN\lambda\R$ as ordered fields? If not, what intrinsic field invariant separates them?
\end{question}
The gap spectrum settles the different-cofinality case only. A monomial-preserving map that also preserves all appropriate strong sums has much more information than an arbitrary field isomorphism. A promising route is to study families of simultaneous cuts and independent nests: a single cut records cofinality, whereas the least size of a jointly realizable system may retain the original support cardinal. A proof would have to define the additional invariant without naming the chosen monomials or normal forms. \mergetag{} Part~I's \cref{fkc:cor:tower} is the set-sized shadow of this question: over a fixed set $\Gamma$, the omission numbers of $K_{\kappa,k}\preccurlyeq K_{\lambda,k}$ also see only $\cf\kappa$ and $\cf\lambda$.

\begin{question}[Theories with an omnific predicate]\label{fkc:sb:q:omnific}
What are the complete first-order theories of $(\KN\kappa\R,\Ok)$? For which pairs $\kappa<\lambda$, if any, is this inclusion elementary in the ordered-field language expanded by an integer-part predicate? Are the pure rings $\Ok$ elementarily equivalent for all uncountable $\kappa$?
\end{question}
Pure fields are elementarily equivalent, and all set module actions factor through the same residue map. Neither fact proves elementarity of the expanded structures or rings. One must examine formulas capable of probing divisibility, support-dependent principal ideals, or relations between the integer part and its fraction field. The absence of a set-sized quotient distinction is a restriction on possible arguments, not a positive answer. \emph{Related (batch 34).} The automorphism report asks the analogous question for the fixed rings of its valued Gaussian-omnific involutions (\lbl{saut:gr:q:theories} \cite{CollSAUT}): for Archimedean real closed fields $F,F'$ of transcendence degree $2^{\aleph_0}$, when are its descended integer parts $R_F$ and $R_{F'}$ elementarily equivalent as rings? Those rings have the same finite quotients and the same solvability of ordinary polynomial equations (\lbl{saut:gr:prop:finite}, \lbl{saut:gr:prop:dio}) but need not be isomorphic, a restriction of the same kind as the one just described. Both questions are open.

\begin{question}[A general support-ideal classification]\label{fkc:sb:q:supportideal}
Replace the cardinal condition $\card S<\kappa$ by a class of allowed set supports closed under subsets and finite algebraic support operations. Which additional hypotheses characterize real closedness, strong-sum closure, and the least empty-nest cardinal?
\end{question}
The proofs suggest an additivity invariant of the allowed-support system. The main technical tasks are to separate abstract support unions from reverse well-ordered unions and to show when algebraic closure remains controlled by a permitted divisible exponent group. Such a theory could produce intrinsic gap spectra more complicated than one symmetric pair. \mergetag{} Part~I's completion $L_{\kappa,k}$ of \eqref{fkc:eq:localfield} is one support system of this kind over a set $\Gamma$ (supports with fewer than $\kappa$ terms below every exponent); the unlabelled question of the hidden-negative-directions report on support-family subfields is related.

\begin{question}[Real-form invariants beyond one spectrum]\label{fkc:sb:q:realforms}
Can one construct, without forcing, families of real forms of $\SC$ with prescribed multiple set-cut characters, or with equal gap spectra but provably nonisomorphic fixed fields?
\end{question}
\Cref{fkc:sb:thm:involutions} supplies a proper-class family indexed by regular cardinals, but not a classification of all real forms. \Cref{fkc:sb:q:equalcf} is already a concrete test case. The support-ideal program is another possible route. Any resulting involution classification must distinguish pure-field conjugacy from conjugacy constrained to preserve a valuation or a coefficient field. \mergetag{} This question is the part of \lbl{univ:q:families} that \cref{fkc:sb:rem:univ} leaves open, in a single universe; the automorphism report's real form with a countable cofinal subset (\lbl{saut:thm:inequivalentrealform}) has a different obstruction, and its batch-32 addition treats the valued case.

\begin{question}[Derivations and integration]\label{fkc:sb:q:derivations}
For which $\kappa$ does a specified surreal derivation, in particular the Berarducci--Mantova derivation, preserve $\KN\kappa\R$? What is the smallest subfield containing $\KN\kappa\R$ and closed under that derivation or under antiderivatives?
\end{question}
No derivative-support estimate is used or proved here. In particular, Taylor closure for an external analytic variable does not imply closure under a derivation acting on the surreal coefficients. The natural first step is a precise support bound for the derivative of a single monomial, followed by the summability and support union analysis for a general normal form. The logarithmic hull collapse warns that an operation on single monomials may already expose unrestricted supports. \mergetag{} Part~I's non-claim N7 disclaims compatibility with this derivation for the set fields; the question is open for both.

\begin{question}[One-sided exponentials and their images]\label{fkc:sb:q:onesided}
How much of $\KN\kappa\R$ is determined by the image of its inherited exponential? Can one classify order-preserving exponential embeddings $\KN\kappa\R\to\KN\kappa\R^{>0}$ up to ordered-field automorphism, or determine which ordered subgroups can occur as their leading-exponent images?
\end{question}
\Cref{fkc:sb:thm:noexp} forbids surjectivity for every such choice, but does not classify non-surjective choices. The inherited map has the particularly strong property of agreeing with real exponentiation and with Gonshor's positive-infinitesimal Taylor behavior. A useful first comparison is the explicit image in \cref{fkc:sb:prop:expimage}. The remaining problem is to decide which other images can occur and whether they can be distinguished up to ordered-field automorphism, including their interaction with positive finite units.

\begin{question}[Detecting the rings beyond set targets]\label{fkc:sb:q:beyondset}
Which naturally defined proper-class targets or class-sized module structures distinguish $\Ok$ from $\Okl\lambda$ when their cofinalities are equal? Is there a useful replacement for a universal small quotient that retains controlled support information?
\end{question}
The small-target theorem eliminates every set-sized target, not just finite rings or countable representations. Any successful invariant of this form must change the size condition or use a richer kind of structure, for example a system of partial maps whose compatibility records the omitted nests. It cannot merely increase the set cardinality of the target.

\begin{question}[Formal proof architecture]\label{fkc:sb:q:formal}
Can the first-outside-exponent lemma, the prefix-gap classification, and the support-controlled ball intersection be formalized over an abstract Hahn structure with explicit universe bounds, then transferred to the surreal construction in the repository?
\end{question}
A practical order is: formalize comparison by the largest coefficient difference; establish support cardinal estimates; prove the separator lemma; prove compatible-prefix gluing; and finally add the cardinal and model-theoretic corollaries. The set-sized reservoir proof can be formalized independently of saturation. Class-wide claims should be expressed through their universe-indexed or set-parameter formulations, not by treating a proper class as an ordinary small type. \mergetag{} Part~I's proof-assistant decomposition (\cref{fkc:sec:audit}) proposes the same layering for the set fields. No \lbl{fkc:} statement is formalized in Lean.

\section{Provenance, novelty boundary, and verification of source 02}\label{fkc:sb:sec:provenance}

\subsection{The supplied repository}
The working baseline is the main-branch snapshot
\begin{center}\small\texttt{343dc2c471212bb9b53ff4623bace2e1943f255b}\end{center}
retrieved on 23 September 2026 \cite{RepoTwo}. The root README and the report catalogue were inspected, together with the guides to set-sized omnific quotients, surreal fields across universes, and the order-three threshold source audit in the holonomic report. These were read through the connected GitHub interface. The catalogue describes a large collection of AI-assisted research drafts and carefully distinguishes their written proofs from mapped Lean coverage.

The review was \emph{targeted}, not a complete audit of every repository source, branch, or prior manuscript. In particular, one GitHub code search returned an incomplete-results flag; it is not evidence that an unmatched term or equivalent theorem is absent. No assertion in this part depends on a claimed successful search of the entire repository.

The across-universes report already contains real-form distinctions, set-cut spectra, and the stabilization of set-indexed Cauchy nets. The omnific quotient report already contains the full-$\Oz$ constant-term universality theorem and scaled-field collision arguments. These are explicitly credited. The present candidate contribution is the support-bounded hierarchy, its singular-cardinal and prefix classification results, and the consequences obtained by combining them with those existing methods. \mergetag{} The comparison is incomplete; see \cref{fkc:sb:sub:corrections}.

\subsection{Public antecedents}
Conway's normal forms and Gonshor's exponential theory are foundational imports; the precise forms needed here are accessible in the primary survey \cite{MM}. Cardinally bounded Hahn fields are not new objects. Kuhlmann and Shelah \cite{KS} construct bounded-support exponential fields by arranging the value group to support a logarithm. Their main construction uses a regular uncountable cardinal; it is not the fixed proper-class value group $\No$ used here.

The impossibility of exponentiation on full nontrivial set-group Hahn fields over $\R$ is classical \cite{KKS}. The general mechanism of comparing ordered additive and positive multiplicative structures is therefore not new. \Cref{fkc:sb:thm:noexp} instead uses a computed set-cut mismatch in an exp-closed surreal subfield. The omega-field results in \cite{BKMM} provide a related positive comparison when the value group has the required additive self-similarity. Bournez and Guilmant \cite{BG} study exp--log stable surreal subfields under different restrictions, including birthday-based and other Hahn constructions.

The user-supplied Wikipedia article \cite{Wikipedia} was consulted as orientation, not used as the authority for a technical proof. No claim of a literature breakthrough rests on search-engine non-results. The primary-source review did not include an exhaustive MathSciNet or zbMATH search or a full subscription-literature audit. Conway's and Gonshor's books \cite{Conway,Gonshor} are credited for foundational results accessed through the survey \cite{MM}, not represented as newly read in full.

\subsection{Claim-by-claim contribution boundary}\label{fkc:sb:sub:claims}
The table is source~02's, as delivered; the merge's corrections to it are in \cref{fkc:sb:sub:corrections}.
\begin{center}
\small
\begin{longtable}{@{}p{0.45\textwidth}p{0.47\textwidth}@{}}
\toprule
Statement or method & Status in source 02 \\
\midrule\endhead
Conway normal forms; real closed Hahn fields; surreal exponential on purely infinite numbers & Classical imports, cited explicitly. \\
Cardinally bounded Hahn construction & Classical framework; fixed full surreal exponent class and all uncountable bounds treated here. \\
First-$\kappa$ prefix classification and exact gap character & Proposed contribution, with a direct proof including singular $\kappa$. \\
Support-sensitive realization of types & Proposed contribution in this formulation; uses classical real closed field quantifier elimination. \\
Exact strong-sum and empty-nest thresholds & Proposed hierarchy results; support gluing and cardinal union estimates are elementary Hahn mechanisms. \\
Set Cauchy nets eventually constant & Existing proper-class-rank observation; not claimed as new. \\
Inherited exponential closure but no exponential isomorphism & Proposed specific obstruction; classical quotient philosophy credited. \\
All-monomials logarithmic hull equals $\No$ & Direct corollary of Gonshor's classical formula and a monomial shift; included as a useful structural result, without a strong priority claim. \\
Universal set quotient of full $\Oz$ & Prior repository result, not reproved as a new discovery. \\
Uniform extension to every $\Ok$ & Proved here using a countably supported large reservoir; candidate extension. \\
Class back-and-forth and existence of nonconjugate real forms & Prior mechanisms; not new in themselves. \\
Intrinsic proper-class family indexed by regular support bounds & Proposed consequence of the exact gap spectrum; no forcing needed. \\
\bottomrule
\end{longtable}
\end{center}

\subsection{Finite checks and their limits}
The accompanying program \texttt{code/02-support-bounded-fields-verify.py} passed 31,672 exact rational checks of finite analogues of the comparison, truncation, bracket, valuation-ball, floor, and collision identities. It uses no floating point arithmetic and records its actual run in \texttt{data/02-support-bounded-fields-verification.json}. These checks are intended to catch sign, strictness, and boundary errors in the displayed formulas.

No finite enumeration establishes the existence of a missing $\kappa$-term point, the cofinality of a singular cardinal, class back-and-forth, or the impossibility of an exponential isomorphism. Those claims depend on the written proofs. No new theorem here is claimed to have been checked by Lean, and no existing repository build is represented as verifying this part.

\subsection{Corrections to source 02's repository statements \mergetag}
\label{fkc:sb:sub:corrections}
Source~02's statements above describe the collection at its pin \texttt{343dc2c}, which is kept as provenance. Checked against the collection at the pin and as it now stands, they need the following corrections.
\begin{enumerate}[label=\textup{(\arabic*)}]
\item \emph{The set version was present.} Part~I of this report was in the collection at the pin. Its \cref{fkc:thm:main} is the set-sized version, over an arbitrary set $\Gamma$, of source~02's field, type, cofinality and nest results, with the same small-workspace and gluing lemmas. Source~02 does not cite it. The merge prints the common arguments once and credits both sources (\cref{fkc:sb:sub:provenance,fkc:sb:sub:relation}).
\item \emph{The large-cardinal report.} At the pin, the large-cardinal report already defined the same class, as its field of short normal forms $\mathscr S_{<\kappa}$, at the critical point $\kappa$ of an elementary embedding, with real closedness, the integer part $\mathscr I_{<\kappa}$ and its fraction field, and closure under the exponential but not the logarithm with the witness $z_\kappa$ (then with bare labels, now \lbl{lce:cor:expclosure}, \lbl{lce:eq:Ozfrac}). After the pin it gained the short-field theorem for every uncountable $\kappa$ (\lbl{lce:sg:thm:shortfield}), closedness and nowhere density in $\No$ (\lbl{lce:sg:prop:closed-support}), the integer part for every uncountable $\kappa$ (\lbl{lce:sg:prop:integerpart}) and the logarithmic-closure collapse (\lbl{lce:thm:twoclosures}). Source~02 does not cite it; the corresponding results are credited where they are printed. What remains source~02's is the treatment at every uncountable $\kappa$ without large cardinals of the exponential closure and the log-hull collapse, and everything not listed here.
\item \emph{The small-target theorem is stale.} Source~02's row ``Uniform extension to every $\Ok$ \ldots\ candidate extension'' and its sentence after \cref{fkc:sb:thm:smallquotient} imply that the quotient report treats only the full $\Oz$. At the pin that report already contained \lbl{osq:thm:support}~(i), which with $D=\Z$, $k=\R$ is \cref{fkc:sb:thm:smallquotient}, and with $D=\Z[i]$, $k=\C$ the Gaussian factorization of \cref{fkc:sb:sub:gaussian}. The theorem is therefore not new; source~02's reservoir proof is a new route to it (\cref{fkc:sb:rem:osq}), and \cref{fkc:sb:cor:rings} is new.
\item \emph{Two automorphism-report statements.} The class back-and-forth is \lbl{saut:thm:acfhomogeneity}, and non-conjugacy to standard conjugation also follows from \lbl{saut:prop:conjugacycriterion}; both were present at the pin (\cref{fkc:sb:rem:univ}).
\item \emph{Statements found accurate.} The across-universes report does contain forcing-based real forms, set-cut spectra and set-Cauchy stabilization (\lbl{univ:thm:transport}, \lbl{univ:thm:manygrounds}, \lbl{univ:cor:spectrum}, \lbl{univ:prop:topology}); the quotient report does contain the full-$\Oz$ theorem and the scaled-field collision (\lbl{osq:thm:universal}, \lbl{osq:lem:collision}); and the foundations report contains the set-Cauchy stabilization (\lbl{found:thm:discrete}).
\end{enumerate}

\subsection{What Part II does not claim}
\label{fkc:sb:sub:nonclaims}
Source~02 states its limitations where they arise. They are collected here, continuing the numbering of Part~I's list (\cref{fkc:subsec:nonclaims}); the last two were added in the merge. None is withdrawn.
\begin{enumerate}[label=\textup{N\arabic*.}, leftmargin=2.8em, start=13]
\item Source~02 is an AI-assisted, unrefereed research manuscript. Its proofs are written mathematical arguments, not Lean proofs, and have had no independent peer review; no existing repository build verifies it.
\item Bibliographic priority is not certified. The repository review was targeted and guide-level, an incomplete code search was not used as evidence, and no exhaustive MathSciNet, zbMATH or subscription-literature audit was made. The Wikipedia article was orientation only; Conway and Gonshor were accessed through the Mantova--Matusinski survey.
\item No named published conjecture is claimed solved, and the eight questions of \cref{fkc:sb:sec:questions} are not claimed to be published open problems.
\item Only the set-presented gaps of the fields $\KN\kappa\R$ are classified. Class-cofinal gaps, all surreal gaps, all real forms of $\SC$ and all bounded Hahn fields are not.
\item Whether $\cf\kappa=\cf\lambda$ implies $\KN\kappa\R\cong\KN\lambda\R$ is open (\cref{fkc:sb:q:equalcf}).
\item The class isomorphisms $\psi_\kappa$ and the involutions $\tau_\kappa$ are not claimed to preserve monomials, the valuation, the exponential, Hahn sums or the omnific ring.
\item $\kappa=\aleph_0$ is excluded: finite-support series do not form a field.
\item Strong sums are support operations, not topological limits.
\item The exact exponential image (\cref{fkc:sb:prop:expimage}) and the logarithmic hull collapse (\cref{fkc:sb:thm:loghull}) are direct consequences of Gonshor's classical formulas, with no strong priority claim.
\item Credited, not claimed: cardinally bounded Hahn fields (Kuhlmann--Shelah; F.-V.~Kuhlmann, S.~Kuhlmann and Shelah), the general additive-versus-multiplicative exponential obstruction, the class back-and-forth, set-Cauchy stabilization, forcing-based nonconjugate real forms (the across-universes report), and the full-$\Oz$ quotient theorem with the scaled-field collision (the quotient report).
\item \Cref{fkc:sb:thm:noexp} does not contradict the $\kappa$-bounded exponential constructions of Kuhlmann--Shelah or the omega-field results of Berarducci--Kuhlmann--Mantova--Matusinski.
\item \Cref{fkc:sb:thm:loghull} needs every monomial; nothing is claimed about exp--log closed subfields in general.
\item The Gaussian rings $\Ok[i]$ are not claimed to be pairwise nonisomorphic.
\item No derivative-support estimate is proved, and no closure under a surreal derivation is claimed.
\item The 31,672 finite checks cover finite identities and boundary conventions only: not cardinal or cofinality arguments, infinite Hahn summability, class constructions or back-and-forth, saturation, the nonexistence of a surjective exponential, novelty, or any Lean formalization.
\item Global Choice is used only in \cref{fkc:sb:sec:realforms}; no class form of Zorn's lemma and no truth predicate for a class structure is used, and the ``proper-class family'' is a single class relation.
\item Added in the merge: source~02's novelty claim for \cref{fkc:sb:thm:smallquotient} is stale, and several of its results were in the collection in set-sized or large-cardinal form (\cref{fkc:sb:sub:corrections}); only the route and the uses of the result in \cref{fkc:sb:cor:rings} are new there.
\item Added in the merge: the comparisons with Part~I and with other reports, and the status notes they motivate, prove no new theorem; \cref{fkc:sb:rem:univ}(2) reads \cref{fkc:sb:thm:involutions} against another report's question and answers one clause of it in a single universe only.
\end{enumerate}

\section{Conclusion of Part II}\label{fkc:sb:sec:conclusion}

Imposing a support cardinal bound while retaining every surreal monomial produces real closed fields with sharply different global structure. Their set-presented gaps can be read exactly from missing normal-form prefixes. The cofinality of the bound controls saturation, strong summation, and nested-ball intersection, while ordinary set-indexed Cauchy completeness is insensitive to it.

Exponentiation makes the contrast more pronounced: the inherited surreal exponential remains defined internally, but the ordered field cannot admit any onto exponential. On the arithmetic side, all set-sized targets still collapse the associated omnific integer part to its ordinary integer constant term. On the complex side, forgetting conjugation erases the distinctions, whereas remembering the real form recovers them.

These are precise results about a specified hierarchy, not a complete solution to surreal model theory or real-form classification. They provide explicit test objects and proof mechanisms for the further questions above, with especially concrete next steps at singular support bounds and in isomorphism questions with equal cofinality.

\appendix
\crefalias{section}{appendix}
\section{Proof-risk checklist and finite verification}
\label{fkc:app:checks}

\paragraph{Cardinal bounds.}
The entire discussion assumes $\kappa>\aleph_0$. At $\kappa=\aleph_0$, finite-support series are generally not a field: $(1-t^g)^{-1}=\sum_{n\ge0}t^{ng}$ for $g>0$ has infinite support. Regularity is not assumed for $\kappa$; it is used only for $\mu=\cf(\kappa)$, which is regular. The small-workspace and small-union arguments are separate and use different bounds.

\paragraph{Prefix versus full support.}
A missing series can have support order type larger than $\kappa$. Its type still remembers only the first $\kappa$ positions. By contrast, a newly added completion element must have support order type \emph{exactly} $\kappa$, not merely cardinality $\kappa$. The example $x_\kappa+t^{e_\kappa}$ shows why these statements differ.

\paragraph{Valuation inequalities.}
A closed ball of radius value $\gamma$ fixes coefficients \emph{below} $\gamma$. A strict valuation inequality fixes coefficients \emph{at and below} $\gamma$. An ordered inequality $|z|<t^\gamma$ forces $v(z)\ge\gamma$, not necessarily $v(z)>\gamma$. The real-type proofs use later cofinal thresholds to recover the coefficient at a given position; they do not make the stronger, false implication.

\paragraph{Cauchy versus pseudo-Cauchy.}
The identity $v(a_j-a_i)=s_{\alpha_i}$ proves a pseudo-Cauchy condition. It gives a Cauchy net only when those values are cofinal in the whole group. This is why the countable family in the complete-field example is not a missing Cauchy limit. The topology induced from full $\No$ is a third, separate matter.

\paragraph{Scope of quantifier elimination.}
The type theorems use exactly the stated ordered or valued language. In the pure complex ring language the answer changes, as proved. A monomial appearing as a parameter in a formula is not a named global cross-section symbol. No quantifier-elimination result for a language with a prefix operation is assumed.

\paragraph{Computational checks.}
The companion \texttt{code/verify.py} uses exact rational coefficients in finite Laurent-polynomial supports over $\mathbb Z^2$ with lexicographic order. It checks first-disagreement identities, invariance of polynomial leading data under sufficiently high tails, nested-ball coefficient agreement, and the local trichotomy underlying the retraction-loss formula. A fixed seed makes its checks reproducible, and it writes their assertion counts to a file \texttt{verification.json} in its own directory; the recorded run is \texttt{data/verification.json} (9,622 assertions, all passed). These checks test finite algebraic mechanisms only. They do not prove any infinite-support, cardinality, completeness, first-order, or novelty statement. In particular, a finite support cutoff is \emph{not} treated as a field model of an uncountable cardinal cutoff.

The manuscript was delivered with its TeX source, compiled PDF, build script, build record and a separate research-audit file. This report's directory contains the source and PDF of this report, the delivered script \texttt{code/build.sh}, the delivered build record \texttt{data/build\_audit.json} (which describes the delivered 21-page PDF) and \texttt{research\_audit.md}. The script was written for the delivery's flat layout; the README of this directory explains how to rerun it on a copy. The written proofs, finite checks, and document-layout checks are distinct forms of evidence. \mergetag{} The directory now also contains the delivered files of source~02, with the prefix \texttt{02-support-bounded-fields-}; they are described in \cref{fkc:sb:sec:provenance} and in the README.

\section{Proof dependencies and sensitive points of Part II}\label{fkc:sb:app:audit}

This appendix and the next are source~02's appendices.

The arguments have a short dependency structure. The field theorem uses only normal forms, small rational spans, and real closedness of set Hahn fields. The gap theorem uses lexicographic comparison and the surreal set-cut property. Saturation adds set compactness and quantifier elimination. The spherical theorem is a direct gluing argument. The nonexistence of an exponential depends on the gap theorem and the normalization lemma, not on an unproved uniqueness of Gonshor's exponential. The omnific quotient theorem uses an independent scaled reservoir. The real-form theorem adds class back-and-forth.

Several details warrant particular attention in independent review. First, the support cardinality condition is not an ordinal length bound on arbitrary initial segments: the first excluded segment has order type the \emph{initial ordinal} $\kappa$, and earlier segments have cardinality strictly below it. Second, both a missing coefficient and an inserted exponent are handled by first-difference comparison; it is not enough merely to compare the $n$-th nonzero terms of two series. Third, a ball of radius $\gamma$ fixes coefficients at valuation \emph{strictly below} $\gamma$, not at the boundary. This is why $B(h_\beta,\beta)$ is the correct closed ball.

Fourth, $\cf\kappa$ rather than $\kappa$ controls unions of varying supports, but a Taylor series in one fixed support is controlled by its finite-sum monoid. Fifth, a hypothetical exponential is normalized by multiplying its input by $\mathcal E^{-1}(2)$; the proof does not silently assume that it agrees with the real exponential. Sixth, the reservoir field is a rational-expression field with countably supported individual elements, not an arbitrarily large full Hahn field. Finally, abstract field isomorphisms of real closed fields preserve order, whereas abstract complex-field isomorphisms need not preserve any chosen real form.

\mergetag{} The third point is Part~I's ``Valuation inequalities'' paragraph of \cref{fkc:app:checks}, in the decreasing convention: a closed ball of radius $\gamma$ fixes the coefficients at the exponents $e>-\gamma$, that is, at valuation below $\gamma$.

\section{A general fixed-support realization principle}\label{fkc:sb:app:realization}

For clarity, the main separator argument can be isolated independently of cardinal arithmetic. Let $\mathcal S$ be any collection of set supports for which every relevant divisible set group $H$ supports a real closed subfield $\R((\omega^H))$ inside a class field $F$, and suppose $F$ contains every truncation supported in $H$ plus one extra exponent. Then every one-variable type over parameters supported in such an $H$ is realized in $F$. The proof is exactly \cref{fkc:sb:thm:supporttypes}: extend to the set real closed Hahn field, fill the resulting cut in $\No$, and truncate at the first exponent outside $H$.

This formulation identifies the local hypothesis behind saturation. It does not require that the whole field be closed under arbitrary unions of supports, or that the collection of allowed supports be regular-cardinal bounded. Establishing useful support systems satisfying these local conditions is one route to the research program of \cref{fkc:sb:sec:questions}.

\begin{thebibliography}{99}
\addcontentsline{toc}{section}{References}

\bibitem{KS}
S.~Kuhlmann and S.~Shelah,
\emph{$\kappa$-bounded Exponential-Logarithmic Power Series Fields},
Annals of Pure and Applied Logic \textbf{136} (2005), 284--296
(journal reference as given by source~02; DOI 10.1016/j.apal.2005.04.001).
arXiv:math/0512220v1 (2005).
\url{https://arxiv.org/abs/math/0512220}.

\bibitem{BKMM}
A.~Berarducci, S.~Kuhlmann, V.~Mantova, and M.~Matusinski,
\emph{Exponential fields and Conway's omega-map},
Proceedings of the American Mathematical Society \textbf{151} (2023), 2655--2669.
DOI: \href{https://doi.org/10.1090/proc/14577}{10.1090/proc/14577}.
Preprint: \url{https://arxiv.org/abs/1810.03029}.

\bibitem{FK}
F.-V.~Kuhlmann,
\emph{Approximation types describing extensions of valuations to rational function fields},
arXiv:2111.10529v1 (2021), especially Sections~2.4--2.6 and~4.2.
\url{https://arxiv.org/abs/2111.10529}.

\bibitem{Poonen}
B.~Poonen,
\emph{Maximally complete fields},
L'Enseignement Math\'ematique \textbf{39} (1993), 87--106,
especially Lemma~4 and Corollary~4.
Author's paper hosted at Stanford:
\url{https://math.stanford.edu/~conrad/Perfseminar/refs/poonencomplete.pdf}.

\bibitem{Tarski}
A.~Tarski,
\emph{A Decision Method for Elementary Algebra and Geometry},
prepared for publication with the assistance of J.~C.~C.~McKinsey,
RAND Corporation, R-109 (1951).
\url{https://www.rand.org/pubs/reports/R109.html}.

\bibitem{Robinson}
A.~Robinson,
\emph{Complete Theories},
North-Holland, Amsterdam (1956).
The one-sorted valued-field quantifier-elimination attribution is also recorded in the introduction of \cite{Yin}.

\bibitem{Yin}
Y.~Yin,
\emph{Quantifier elimination and minimality conditions in algebraically closed valued fields},
arXiv:1006.1393v1 (2010).
\url{https://arxiv.org/abs/1006.1393}.
The introduction distinguishes the classical one-sorted divisibility language from the richer languages studied in that paper.

\bibitem{Gonshor}
H.~Gonshor,
\emph{An Introduction to the Theory of Surreal Numbers},
London Mathematical Society Lecture Note Series \textbf{110},
Cambridge University Press (1986).
Classical reference for Conway normal forms and their arithmetic; the complete book was not freshly audited for this manuscript.

\bibitem{Repo}
\begingroup\raggedright
V.~Reshetnikov and contributors,
\emph{Surreal}, GitHub repository, revision
\texttt{465a54b479a1ee842cbf7db1689a7d2f6bfe25e1}.
Documentation catalogue:
\url{https://github.com/VladimirReshetnikov/Surreal/blob/465a54b479a1ee842cbf7db1689a7d2f6bfe25e1/docs/README.md}.
Closest report's audit:
\url{https://github.com/VladimirReshetnikov/Surreal/blob/465a54b479a1ee842cbf7db1689a7d2f6bfe25e1/docs/surreal/transcendence-over-bounded-support/research_audit.md}.
See \cref{fkc:sec:audit} for the actual inspection boundary.
\par\endgroup

\begingroup\raggedright
\bibitem{CollUniv}
Surreal collection, \emph{Surreal Fields Across Set-Theoretic Universes},
report \nolinkurl{docs/foundations-and-computation/surreal-fields-across-universes}
(label prefix \texttt{univ:}). Added after the manuscript's pin.

\bibitem{CollBST}
Surreal collection, \emph{Bounded-Support Hahn Arithmetic},
report \nolinkurl{docs/surreal/transcendence-over-bounded-support}
(label prefix \texttt{bst:}).

\bibitem{CollDuals}
Surreal collection, \emph{Three Duals at Surreal Scales},
report \nolinkurl{docs/surcomplex/three-duals-of-hahn-vector-spaces}
(label prefix \texttt{duals:}).

\bibitem{CollBerk}
Surreal collection, \emph{Rank-One Non-Archimedean Analytic Geometry inside the Surcomplex Numbers},
report \nolinkurl{docs/surcomplex/rank-one-berkovich} (bare labels).

\bibitem{CollEnt}
Surreal collection, \emph{Cofinality, Factorization and Scalar Extension for Entire Hahn Functions at Arbitrary Rank},
report \nolinkurl{docs/surcomplex/entire-functions-at-arbitrary-rank}
(label prefix \texttt{ent:}).

\bibitem{CollFound}
Surreal collection, \emph{Foundations for Surreal and Surcomplex Mathematics},
report \nolinkurl{docs/foundations-and-computation/foundations}
(label prefix \texttt{found:}).

\bibitem{CollHND}
Surreal collection, \emph{Hidden Negative Directions in Surcomplex Linear Algebra},
report \nolinkurl{docs/surcomplex/hidden-negative-hermitian-directions}
(label prefix \texttt{hnd:}).
\par\endgroup

\bibitem{Conway}
J.~H.~Conway,
\emph{On Numbers and Games},
Academic Press (1976); second edition, A K Peters (2001).
Foundational normal-form results used in Part~II through the exposition in \cite{MM} (source~02).

\bibitem{MM}
V.~Mantova and M.~Matusinski,
\emph{Surreal numbers with derivation, Hardy fields and transseries: a survey},
arXiv:1608.03413v2 (2016).
\url{https://arxiv.org/abs/1608.03413}.

\bibitem{KKS}
F.-V.~Kuhlmann, S.~Kuhlmann, and S.~Shelah,
\emph{Exponentiation in power series fields},
Proceedings of the American Mathematical Society \textbf{125} (1997), 3177--3183.
arXiv:math/9608214.
\url{https://arxiv.org/abs/math/9608214}.

\bibitem{BG}
O.~Bournez and Q.~Guilmant,
\emph{Surreal fields stable under exponential and logarithmic functions},
arXiv:2201.08199 (2022).
\url{https://arxiv.org/abs/2201.08199}.

\bibitem{Wikipedia}
\emph{Surreal number}, Wikipedia, user-supplied orientation page, consulted 23 September 2026 by source~02; not used as a technical proof source.
\url{https://en.wikipedia.org/wiki/Surreal_number}.

\bibitem{RepoTwo}
\begingroup\raggedright
V.~Reshetnikov,
\emph{Surreal}, GitHub repository, root README and report catalogue, revision
\texttt{343dc2c471212bb9b53ff4623bace2e1943f255b}, inspected by source~02 on 23 September 2026.
\url{https://github.com/VladimirReshetnikov/Surreal/tree/343dc2c471212bb9b53ff4623bace2e1943f255b}.
See \cref{fkc:sb:sec:provenance} for the inspection boundary.
\par\endgroup

\begingroup\raggedright
\bibitem{CollLCE}
Surreal collection, \emph{Critical-Point Defects in Surreal Arithmetic},
report \nolinkurl{docs/foundations-and-computation/large-cardinal-embeddings-and-normal-forms}
(label prefix \texttt{lce:}). Cited in Part~II (merge).

\bibitem{CollOSQ}
Surreal collection, \emph{Set-Sized Quotients of the Omnific Integers},
report \nolinkurl{docs/surreal/set-sized-quotients-of-omnific-integers}
(label prefix \texttt{osq:}). Source~02 inspected its guide at the pin.

\bibitem{CollOPA}
Surreal collection, \emph{Omnific-Preserving Automorphisms},
report \nolinkurl{docs/surreal/omnific-preserving-automorphisms}
(label prefix \texttt{opa:}). Cited in Part~II (merge).

\bibitem{CollSAUT}
Surreal collection, \emph{Automorphisms of the Surcomplex Numbers},
report \nolinkurl{docs/surcomplex/surcomplex-field-automorphisms}
(label prefix \texttt{saut:}). Cited in Part~II (merge).
\par\endgroup

\end{thebibliography}
\end{document}
